\documentclass[11pt,oneside]{book}

\usepackage[T1]{fontenc}
\usepackage[utf8]{inputenc}
\usepackage[english]{babel}
\usepackage{lmodern}
\usepackage{microtype}
\usepackage[a4paper,margin=2.8cm]{geometry}
\usepackage{xcolor}
\usepackage{xurl}
\usepackage[hidelinks]{hyperref}
\usepackage{enumitem}
\usepackage{amsmath,amssymb}
\usepackage{braket}
\usepackage{pitia}

\setlist{itemsep=2pt}

\title{Pítia\\[0.5em]\large How far are we from a useful quantum computer?}
\author{}
\date{}

\begin{document}

\frontmatter
\maketitle
\tableofcontents
\chapter*{How to Read This Book}
\addcontentsline{toc}{chapter}{How to Read This Book}
\markboth{HOW TO READ THIS BOOK}{HOW TO READ THIS BOOK}

\begin{flushright}
\itshape Lasciate ogne fretta, voi ch'intrate.\\[2pt]
\upshape\small not quite Dante, \emph{Inferno}, canto III
\end{flushright}

Over the gate of Hell, Dante wrote: abandon all hope, ye who enter here. This book asks
for less: abandon all haste.

This is a dense book. It is packed with names, dates, formulas and numbers, and almost
every page leans on something said before it. It assumes you will stop, go back a few
pages, check a number, and come back. It was written for curious people: the ones who
want to know why a thing is the way it is, who would rather spend an evening on one page
than skim twenty, and who are not satisfied until the idea and the history behind it
both make sense.

One question runs through the whole book: how far are we from a quantum computer
that solves a useful problem, and which problem would that be? Answering it means
building the language first and then going down a stack, one layer at a time. The
book has one part for the language and one for each layer, and each part answers one
question the main one forces on you.

The first part is the foundations, and the question is what a quantum computer is.
It starts where the field started, with Richard Feynman asking in 1981 why we
simulate nature with machines that do not work the way nature does. Then it builds
the rest: what a computer can do at all, what a qubit is, where quantum mechanics
came from, and the linear algebra underneath it.

The second part is about the problems: which real problems are waiting for such a
machine, and how each one turns into something it can act on. There are several, and
they are not equally promising.

The third part is about the algorithm: what the circuit extracts from the problem,
and at what cost.

The fourth part is about hardware: which technology delivers physical qubits, and
with how much error.

The fifth part is about error correction: how many of those physical qubits it takes
to make one qubit you can trust.

The sixth part is about the market: who sells each layer, and what you can actually
use today.

The parts are in the order in which the questions depend on each other, but the
chapters are written in the order the study reaches each subject. A part can grow
while the one before it is still thin, and when a chapter needs something from
another part, it says where to find it.

\section*{Every idea has a parent}

No idea in this book showed up because it was elegant. Max Planck put energy in
packets in 1900 because nothing else fitted the light coming out of a hot body, and
he called it an act of desperation. Kenneth Wilson put quarks on a lattice in 1974
because the strong force was too strong for the approximations that had worked for
electrodynamics. Peter Shor found his factoring algorithm in 1994 by thinking about a
paper by Daniel Simon that he had helped reject. Shor's nine-qubit code exists because
Rolf Landauer said quantum computers could never correct their errors. The transmon
was designed in 2007 because the superconducting qubits before it were ruined by stray
electric charges. Each idea fixed something its predecessor got wrong and brought
costs of its own. If you know the lineage, you know when the idea applies. If you only
know the rule, you will apply it where it does not belong.

\section*{Five kinds of boxes}

Stories, concepts, questions, reflections and words are kept apart, so you
always know which one you are reading.

\begin{history}{year \textperiodcentered{} place}
An amber box tells a story: who did it, when, and what forced the idea into
existence. The people, the dates and the problems in it are real. No invented
incidents.
\end{history}

\begin{concept}{name}
A blue box defines an idea you will need again. It always ends the same way:

\textbf{What it fixes:} the problem that made the idea worth having.

\textbf{What it costs:} what you give up by using it.
\end{concept}

\begin{question}{the question you were embarrassed to ask}
A teal box asks the question that occurs to anyone reading the paragraph above
it and that nobody asks out loud, and then answers it.
\end{question}

\begin{reflection}{the thing that should bother you}
A purple box stops on a statement that sounds absurd, or that the text just said
in passing as if it were obvious, and works out what it means and why it is
true.
\end{reflection}

\begin{etymology}
A gray box says where a word comes from. The origin of a word is not its
definition, but it often shows what the people who coined it were looking at.
\end{etymology}

The text between the boxes is the argument that connects them.

\section*{Notation}

Quantum states are written in the notation Paul Dirac introduced in 1939: $\ket{\psi}$
for a state, $\bra{\phi}$ for its dual, $\braket{\phi|\psi}$ for the inner product of
the two. You do not need to know it before starting. Chapter~\ref{ch:quantum} says
where it came from, and Chapters~\ref{ch:hilbert} and~\ref{ch:state} say what it
means.

\section*{Numbers and figures}

There is almost no code in this book, but every number in it was computed before it
went in. When a figure shows data, such as random points or a curve, the data come
from a real run with a fixed seed, and the script that produced them sits next to the
figure in the book's sources.

\section*{Exact facts, firm opinions}

This is not a scientific book, and it does not pretend to be one.

The history is checked. When the book says who did something and when, that is
what happened. Hardware numbers, error rates and resource estimates change every
few months, so every one of them comes with the date it was true.

The opinions are another matter. When a press release announces a qubit nobody has
shown to exist, the book calls it marketing. When a quantum machine learning paper does
not say how the data get into the qubits, the book says the speedup is not there. When
an analogy misleads, like the coin that is heads and tails at once, the book says so,
even when an earlier chapter used it. Sometimes that will be unfair or imprecise, and
we accept that.

A reader who leaves with a position can argue with it, refine it, or throw it
away. A reader who leaves with ``it depends'' has nothing to hold. Between being
precise and leaving you with a position, this book prefers to give you a
position.


\mainmatter

% The foundations first, starting with Feynman.
\part{Foundations}\label{part:foundations}
\chapter{A Short History of an Idea}\label{ch:history}

\begin{etymology}
\textbf{Algorithm} comes from the name of the Persian mathematician
Muhammad ibn Musa al-Khwarizmi, who worked in Baghdad in the ninth century.
Latin translations of his book on Hindu--Arabic numerals began with
``\emph{Dixit Algorizmi}'', ``thus spoke al-Khwarizmi'', and his name became the
word for a step-by-step method.

\textbf{Oracle} comes from the Latin \emph{oraculum}, from \emph{orare}, ``to
speak'' or ``to pray'': the place where you asked a god a hard question. In
quantum computing an oracle is a black box that answers questions. You know what
it answers, not how.
\end{etymology}

In 1981 Richard Feynman asked why we try to simulate nature with machines that do
not work the way nature does. Quantum computing starts with that question. It was
built in layers over the next fifteen years, and each layer answered a question
the previous one had left open. After Feynman, someone asked whether there was a
universal quantum computer. Then came the first problems where such a computer
would win, the first problem anyone outside physics cared about, the objection
that it could never work, and the answer to the objection.

\section{Feynman's question (1981)}

\begin{history}{May 1981 \textperiodcentered{} Endicott House, MIT}
The conference on the Physics of Computation, organized by MIT and IBM, brought
together physicists and computer scientists at Endicott House, near Boston.
Richard Feynman gave a talk called ``Simulating Physics with Computers'',
published in 1982. Paul Benioff had already shown, in 1980, that a Turing
machine could be described with quantum mechanics, and Yuri Manin had raised a
similar point in Russian the same year. Feynman asked the question the other
way around.
\end{history}

\begin{history}{1981 \textperiodcentered{} Caltech}
Feynman's great work on quantum electrodynamics was thirty years behind him, and
had earned him the Nobel Prize in 1965. In 1981 the hard problem in particle
physics was quantum chromodynamics, the theory of quarks and gluons, and Feynman
published a paper on it that year, on Yang--Mills theory in two space dimensions.
The theory was known, but almost nothing could be computed from it. Since Kenneth
Wilson's lattice formulation of 1974, physicists had tried to extract numbers such
as the mass of the proton by simulating the theory on a grid of points, and by
1980 Michael Creutz was doing it with Monte Carlo methods on the largest computers
available. At the same time, Feynman had been talking computation for years with
his friend Edward Fredkin, who taught him computing in exchange for lessons in
quantum mechanics. When Fredkin invited him to give the opening talk at Endicott
House, Feynman at first refused.
\end{history}

The talk was about simulating physics in general, not quarks. Simulating a
molecule means describing its quantum state, and the number of numbers needed to
describe that state grows exponentially with the number of particles. Each new
electron multiplies the memory by a constant factor. A few dozen particles are
enough to exceed any classical memory.

Feynman went further than memory. A classical computer does not have to store the
whole state if it can sample it instead, and that is what Creutz's Monte Carlo
calculations did.

\begin{history}{1946--1949 \textperiodcentered{} Los Alamos}
Stanisław Ulam, recovering from an illness in 1946, wondered about the chances of
winning a game of solitaire. Counting every possible game was hopeless, so he
thought of playing a few hundred games and counting the wins. He saw at once that
the same trick could follow neutrons through the material of a bomb, a calculation
nobody could do exactly. John von Neumann turned the idea into programs for ENIAC.
Nicholas Metropolis named it \emph{Monte Carlo}, after the casino in Monaco where
Ulam's uncle used to gamble, and Metropolis and Ulam published it in 1949.
\end{history}

The idea fits in one picture. To find the area of a quarter circle, you don't need
geometry: throw random points into the square around it and count how many land
inside. The fraction that lands inside approaches the fraction of the area, which
is $\pi/4$.

\begin{figure}[htbp]
\centering
% Generated by a seeded run (numpy default_rng(1946)); do not edit by hand.
\def\mcInside{0.895/0.375, 0.671/0.491, 0.666/0.562, 0.438/0.218, 0.950/0.075, 0.536/0.721, 0.787/0.025, 0.282/0.476, 0.296/0.056, 0.695/0.051, 0.309/0.883, 0.373/0.506, 0.451/0.702, 0.380/0.325, 0.932/0.196, 0.454/0.544, 0.631/0.179, 0.255/0.108, 0.703/0.515, 0.509/0.585, 0.136/0.394, 0.794/0.409, 0.363/0.316, 0.335/0.552, 0.795/0.238, 0.322/0.721, 0.178/0.902, 0.541/0.175, 0.670/0.454, 0.105/0.555, 0.259/0.864, 0.239/0.480, 0.884/0.145, 0.761/0.547, 0.201/0.626, 0.307/0.208, 0.392/0.896, 0.562/0.532, 0.505/0.320, 0.503/0.365, 0.640/0.729, 0.132/0.086, 0.160/0.160, 0.586/0.532, 0.093/0.297, 0.262/0.265, 0.278/0.845, 0.161/0.372, 0.745/0.349, 0.386/0.661, 0.480/0.793, 0.665/0.248, 0.532/0.250, 0.452/0.248, 0.838/0.405, 0.609/0.401, 0.096/0.914, 0.452/0.312, 0.029/0.117, 0.261/0.019, 0.456/0.130, 0.093/0.935, 0.868/0.154, 0.222/0.458, 0.022/0.149, 0.443/0.588, 0.501/0.407, 0.542/0.293, 0.093/0.903, 0.399/0.058, 0.143/0.219, 0.560/0.179, 0.101/0.672, 0.910/0.083, 0.659/0.464, 0.718/0.014, 0.074/0.064, 0.872/0.323, 0.469/0.841, 0.747/0.186, 0.402/0.062, 0.731/0.423, 0.770/0.347, 0.141/0.583, 0.547/0.824, 0.093/0.673, 0.205/0.955, 0.724/0.564, 0.311/0.083, 0.293/0.661, 0.583/0.120, 0.216/0.215, 0.752/0.072, 0.588/0.228, 0.772/0.261, 0.895/0.372, 0.040/0.915, 0.354/0.489, 0.029/0.127, 0.696/0.378, 0.370/0.040, 0.897/0.025, 0.311/0.890, 0.382/0.624, 0.776/0.315, 0.718/0.313, 0.240/0.033, 0.543/0.502, 0.073/0.409, 0.343/0.032, 0.957/0.044, 0.287/0.300, 0.373/0.370, 0.164/0.809, 0.502/0.308, 0.851/0.460, 0.100/0.363, 0.501/0.161, 0.984/0.142, 0.680/0.711, 0.060/0.746, 0.327/0.598, 0.391/0.801, 0.470/0.021, 0.618/0.259, 0.543/0.754, 0.820/0.312, 0.749/0.587, 0.802/0.539, 0.517/0.479, 0.256/0.568, 0.038/0.927, 0.645/0.505, 0.566/0.302, 0.487/0.497, 0.881/0.390, 0.965/0.150, 0.853/0.184, 0.494/0.058, 0.542/0.185, 0.758/0.040, 0.229/0.312, 0.004/0.705, 0.238/0.926, 0.896/0.394, 0.488/0.483, 0.080/0.104, 0.297/0.592, 0.343/0.028, 0.795/0.375, 0.264/0.087, 0.154/0.905, 0.034/0.338, 0.069/0.571, 0.621/0.162, 0.344/0.900, 0.196/0.323, 0.623/0.339, 0.880/0.061, 0.326/0.404, 0.163/0.413, 0.574/0.565, 0.315/0.152, 0.561/0.635, 0.417/0.361, 0.489/0.468, 0.185/0.910, 0.440/0.091, 0.021/0.186, 0.377/0.251, 0.585/0.639, 0.464/0.069, 0.358/0.551, 0.255/0.272, 0.592/0.634, 0.554/0.500, 0.775/0.584, 0.189/0.518, 0.032/0.841, 0.254/0.089, 0.520/0.540, 0.022/0.209, 0.109/0.882, 0.278/0.297, 0.407/0.453, 0.624/0.051, 0.860/0.509, 0.514/0.822, 0.026/0.915, 0.419/0.822, 0.063/0.750, 0.033/0.837, 0.139/0.942, 0.578/0.794, 0.058/0.454, 0.155/0.587, 0.917/0.065, 0.455/0.153, 0.447/0.620, 0.107/0.313, 0.381/0.856, 0.355/0.461, 0.502/0.575, 0.412/0.401, 0.065/0.229, 0.246/0.625, 0.684/0.201, 0.906/0.131, 0.409/0.196, 0.486/0.124, 0.168/0.254, 0.586/0.182, 0.444/0.719, 0.627/0.220, 0.086/0.605, 0.044/0.203, 0.169/0.138, 0.057/0.173, 0.884/0.273, 0.893/0.436, 0.382/0.029, 0.318/0.382, 0.014/0.592, 0.459/0.691, 0.559/0.590, 0.388/0.334, 0.431/0.790, 0.165/0.577, 0.352/0.751, 0.257/0.180, 0.587/0.299, 0.732/0.242, 0.190/0.374, 0.385/0.849, 0.860/0.130, 0.225/0.856, 0.061/0.778, 0.142/0.110, 0.041/0.973, 0.556/0.487, 0.450/0.856, 0.011/0.202, 0.435/0.347, 0.256/0.075, 0.660/0.323, 0.028/0.867, 0.642/0.139, 0.259/0.584, 0.122/0.238, 0.216/0.693, 0.170/0.522, 0.226/0.908, 0.511/0.488, 0.193/0.205, 0.545/0.065, 0.170/0.119, 0.034/0.778, 0.506/0.836, 0.015/0.768, 0.121/0.823, 0.409/0.475, 0.446/0.657, 0.576/0.750, 0.821/0.133, 0.161/0.134, 0.677/0.001, 0.240/0.642, 0.393/0.850, 0.757/0.182, 0.463/0.304, 0.182/0.171, 0.763/0.428, 0.238/0.314, 0.650/0.470, 0.233/0.143, 0.228/0.313, 0.873/0.076, 0.444/0.176, 0.062/0.561, 0.510/0.578, 0.415/0.702, 0.069/0.399, 0.712/0.336, 0.515/0.578, 0.419/0.242, 0.650/0.085, 0.793/0.513, 0.602/0.676, 0.170/0.884, 0.576/0.096, 0.014/0.632, 0.680/0.699, 0.276/0.329, 0.719/0.617, 0.325/0.418, 0.472/0.100, 0.673/0.190, 0.851/0.320, 0.146/0.896, 0.284/0.921, 0.640/0.231, 0.490/0.171, 0.356/0.425, 0.160/0.094, 0.496/0.650, 0.484/0.348, 0.482/0.579, 0.013/0.249, 0.128/0.826, 0.506/0.654, 0.288/0.565, 0.091/0.299, 0.699/0.149, 0.784/0.284, 0.751/0.184, 0.864/0.154, 0.396/0.727, 0.875/0.334, 0.360/0.334, 0.519/0.307, 0.581/0.751, 0.098/0.634, 0.722/0.099}
\def\mcOutside{0.807/0.764, 0.907/0.525, 0.844/0.830, 0.950/0.956, 0.999/0.274, 0.811/0.932, 0.910/0.966, 0.609/0.809, 0.596/0.911, 0.864/0.546, 0.957/0.596, 0.917/0.462, 0.703/0.768, 0.996/0.968, 0.546/0.844, 0.756/0.967, 0.733/0.938, 0.810/0.742, 0.627/0.905, 0.810/0.934, 0.864/0.553, 0.829/0.582, 0.697/0.969, 0.713/0.912, 0.570/0.962, 0.735/0.949, 0.835/0.946, 0.528/0.879, 0.791/0.973, 0.931/0.938, 0.987/0.502, 0.894/0.810, 0.793/0.918, 0.624/0.844, 0.816/0.647, 0.983/0.665, 0.975/0.378, 0.668/0.847, 0.894/0.958, 0.604/0.846, 0.338/0.997, 0.760/0.982, 0.784/0.941, 0.969/0.725, 0.545/0.978, 0.872/0.694, 0.815/0.965, 0.603/0.833, 0.497/0.871, 0.729/0.764, 0.878/0.549, 0.848/0.943, 0.786/0.817, 0.637/0.869, 0.798/0.763, 0.874/0.956, 0.705/0.932, 0.170/0.991, 0.951/0.908, 0.424/0.942, 0.690/0.863, 0.612/0.846, 0.919/0.869, 0.788/0.731, 0.638/0.982, 0.944/0.804, 0.925/0.813, 0.488/0.878, 0.830/0.791, 0.391/0.959, 0.693/0.946, 0.931/0.743, 0.688/0.788, 0.773/0.943, 0.622/0.887, 0.914/0.908, 0.700/0.789}
\def\mcCurve{(10,3.2000) (12,3.0000) (14,3.1429) (16,3.2500) (19,3.3684) (22,3.4545) (26,3.3846) (30,3.3333) (35,3.4286) (41,3.3171) (48,3.2500) (56,3.2857) (65,3.2000) (76,3.2632) (89,3.3708) (104,3.3462) (122,3.4098) (142,3.3239) (166,3.2771) (194,3.2990) (227,3.2423) (265,3.2906) (310,3.2645) (363,3.2397) (424,3.2264) (495,3.1919) (579,3.1986) (677,3.1965) (791,3.2263) (925,3.2043) (1081,3.2155) (1264,3.1930) (1477,3.1984) (1727,3.1940) (2019,3.1778) (2360,3.1780) (2759,3.1751) (3225,3.1702) (3769,3.1722) (4406,3.1720) (5151,3.1722) (6021,3.1729) (7038,3.1583) (8227,3.1530) (9617,3.1552) (11242,3.1560) (13141,3.1568) (15362,3.1519) (17957,3.1546) (20991,3.1539) (24538,3.1532) (28683,3.1569) (33529,3.1610) (39194,3.1639) (45816,3.1564) (53557,3.1522) (62605,3.1527) (73182,3.1479) (85547,3.1489) (100000,3.1476)}
\def\mcShowN{400}
\def\mcShowIn{323}
\def\mcShowEst{3.23}
\def\mcFinalEst{3.1476}

\begin{tikzpicture}
  \begin{scope}[scale=4.6]
    \fill[hfConcept!6] (0,0) -- (1,0) arc (0:90:1) -- cycle;
    \draw[black!40] (0,0) rectangle (1,1);
    \draw[hfConcept, line width=0.8pt] (1,0) arc (0:90:1);
    \foreach \x/\y in \mcInside { \fill[hfConcept] (\x,\y) circle (0.008); }
    \foreach \x/\y in \mcOutside { \fill[black!45] (\x,\y) circle (0.008); }
    \node[hf/note, below=1mm, text width=44mm] at (0.5,0) {\mcShowIn{} of \mcShowN{}
      points inside\\$\pi \approx 4 \times \mcShowIn/\mcShowN = \mcShowEst$};
  \end{scope}
  \begin{semilogxaxis}[at={(66mm,0)}, width=66mm, height=52mm,
      xmin=10, xmax=100000, ymin=2.9, ymax=3.5,
      xlabel={number of points}, ylabel={estimate of $\pi$},
      ylabel style={yshift=2mm},
      tick label style={font=\sffamily\scriptsize},
      label style={font=\sffamily\small},
      axis lines=left, axis line style={black!50}]
    \addplot[hfHistory, dashed, domain=10:100000, samples=2] {pi};
    \addplot[hfConcept, line width=0.8pt, mark=none] coordinates {\mcCurve};
  \end{semilogxaxis}
\end{tikzpicture}
\caption{Monte Carlo. Left: random points in a square, counted inside and outside a
quarter circle. Right: the estimate as the number of points grows, against the true
value of $\pi$ (dashed). After 100\,000 points it reads \mcFinalEst. The error
shrinks like $1/\sqrt{N}$: each extra correct digit costs a hundred times more
points.}
\end{figure}

The same trick works on Wilson's lattice. In 1974 Kenneth Wilson replaced
continuous space and time with a grid of points a small distance $a$ apart. Quarks
live on the points, the \emph{sites}. The gluon field lives on the \emph{links}
between neighboring sites: each link carries a $3\times3$ unitary matrix $U$ that
says how a quark's color is rotated when it hops from one site to the next. Wilson
put the gluons on the links, and not on the sites, because that keeps the symmetry
of the theory exact even on a coarse grid. The smallest closed loop, a square of
four links called a \emph{plaquette}, multiplies four of those matrices. How far
the product is from the identity measures the strength of the field inside the
square, and adding it up over all plaquettes gives the \emph{action} $S$ of the
configuration.

\begin{figure}[htbp]
\centering
\begin{tikzpicture}[x=13mm, y=13mm]
  % plaquette
  \fill[hfHistory!12] (2,1) rectangle (3,2);
  \draw[hfHistory, line width=0.7pt, -{Stealth[length=1.8mm]}]
    (2.5,1.5) ++(-40:0.27) arc (-40:250:0.27);
  % links
  \foreach \x in {0,...,5} \foreach \y in {0,...,3} {
    \ifnum\x<5 \draw[hfConcept!70, line width=0.8pt] (\x,\y) -- (\x+1,\y); \fi
    \ifnum\y<3 \draw[hfConcept!70, line width=0.8pt] (\x,\y) -- (\x,\y+1); \fi
  }
  % one highlighted link
  \draw[hfConcept, line width=2pt, -{Stealth[length=2.6mm]}] (4,2) -- (4.92,2);
  \node[font=\sffamily\small, text=hfConcept, above] at (4.5,2) {$U$};
  % sites
  \foreach \x in {0,...,5} \foreach \y in {0,...,3}
    \fill[hfSource] (\x,\y) circle (2.4pt);
  % spacing
  \draw[black!60, decorate, decoration={brace, amplitude=4pt, mirror}]
    (0,-0.12) -- (1,-0.12) node[midway, below=4pt, font=\sffamily\small, text=black] {$a$};
  % axes
  \draw[black!50, -{Stealth[length=2mm]}] (-0.45,-0.7) -- (1.0,-0.7)
    node[right, hf/note] {space};
  \draw[black!50, -{Stealth[length=2mm]}] (-0.45,-0.7) -- (-0.45,0.7)
    node[above, hf/note] {time};
  % labels
  \node[hf/note, anchor=west, text width=40mm, align=left] (ls) at (5.45,3) {\textbf{site}: a point of spacetime, where the quarks live};
  \draw[black!40] (ls.west) -- (5.08,3);
  \node[hf/note, anchor=west, text width=40mm, align=left, text=hfConcept] (ll) at (5.45,2) {\textbf{link}: the gluon field, a $3\times3$ matrix $U$};
  \draw[hfConcept!50] (ll.west) -- (5.0,2);
  \node[hf/note, anchor=west, text width=40mm, align=left, text=hfHistory] (lp) at (5.45,1) {\textbf{plaquette}: four links around a square; their product measures the field};
  \draw[hfHistory!60] (lp.west) -- (3.05,1.5);
\end{tikzpicture}
\caption{A slice of Wilson's lattice, two of its four dimensions. One
\emph{configuration} is a choice of matrix for every link. Monte Carlo samples
configurations instead of summing over all of them.}
\end{figure}

A configuration is one choice of matrix for every link, and the theory says to sum
over all of them. A modest lattice of $32^4$ sites has 4\,194\,304 links, and each
$3\times3$ matrix of this kind has 8 free real parameters: the sum is an integral
over more than 33 million variables. No computer can do that integral directly.
Instead, you sample configurations at random, each with probability proportional to
$e^{-S}$, and average over the samples, exactly like the points in the square.

It works because $e^{-S}$ is a positive number, so it behaves like a probability.
And it is positive only thanks to a trick: Wilson's lattice runs in \emph{imaginary}
time. In real time, quantum mechanics weighs each configuration with $e^{iS}$, a
complex number that turns around the circle as $S$ changes. Replacing $t$ with $-i\tau$
turns that rotating weight into a positive one. For a proton at rest, which does not
change in time, nothing is lost. For anything that moves, like a collision, the
trick is not available.

Feynman showed that this cannot work for quantum mechanics in general. Quantum
amplitudes are not positive numbers; they carry signs and phases, and they cancel
(Chapter~\ref{ch:qubit}). When the weights can be negative, random samples mostly
cancel each other and the noise swamps the answer. Physicists later named this the
\emph{sign problem}. The correlations that quantum mechanics predicts, the ones John
Bell had turned into an inequality in 1964, cannot be reproduced by any local
classical process running on ordinary probabilities. ``Nature isn't classical, dammit,'' he
said, ``and if you want to make a simulation of nature, you'd better make it quantum
mechanical.''

Feynman's proposal was simple: use a quantum system you can control to imitate
the quantum system you want to study. Prepare it, let it evolve according to the
dynamics of interest, measure at the end. The physics does the work instead of
the processor. At that point the idea was simulation, not general-purpose
computing.

\section{A universal machine (1985)}

In 1985, David Deutsch, at Oxford, asked the general question: is there a model
of quantum computation that is \emph{universal}, able to simulate any physical
system efficiently? He defined the \emph{quantum Turing machine} and argued that
it was. In the same paper he gave the first problem where a quantum machine does
better than a classical one: deciding whether a function of one bit is constant
or balanced. A classical machine has to ask the function twice. Deutsch's
version asked once and got the answer half of the time; the version that always
works came in 1998.

The problem is artificial and the gain is tiny, but the separation is clean. This
is where quantum computing becomes a model of computation, not just a way of
simulating physics.

\section{The algorithms (1992--1996)}

\begin{history}{1992--1993 \textperiodcentered{} Oxford and Berkeley}
In 1992, Deutsch and Richard Jozsa extended the 1985 problem to functions of $n$
bits: the quantum algorithm needs two queries, while a deterministic classical
algorithm needs, in the worst case, $2^{n-1}+1$. Richard Cleve, Artur Ekert,
Chiara Macchiavello and Michele Mosca brought it down to one query in 1998,
which is the form taught today. In 1993,
Ethan Bernstein and Umesh Vazirani put the quantum Turing machine on rigorous
ground and defined BQP, the class of problems a quantum computer solves
efficiently (Chapter~\ref{ch:turing}).
\end{history}

Deutsch--Jozsa is a toy. A classical algorithm that is allowed to guess at random
solves it with high probability in a handful of queries, and nobody needs to know
whether a function is constant. Its value was to show that the separation could
be large, not to solve anything.

\begin{history}{1994 \textperiodcentered{} STOC, FOCS and Bell Labs}
Daniel Simon found a problem with an exponential separation, even against
randomized classical algorithms: finding a hidden period in a function defined
on the corners of an $n$-dimensional cube. His method was, in essence, a Fourier
transform over those corners. He sent the paper to STOC, one of the two main
conferences in theoretical computer science, and it was rejected. Peter Shor was
on the program committee. He later said he argued for accepting it, though not
hard enough.

Shor had seen what Simon's paper meant. Fourier transforms find periods; the
discrete logarithm problem behind Diffie--Hellman is a period problem on a torus
instead of a cube. He solved the discrete logarithm first and presented it in an
internal seminar at Bell Labs in the spring of 1994. A few days later he had
factoring.
By the time Umesh Vazirani called him, the rumor had already turned ``discrete
log'' into ``factoring''. Simon's paper and Shor's were both accepted at FOCS, in
November 1994.
\end{history}

Shor's insight was that factoring and the discrete logarithm are period-finding
problems, the same kind as Simon's, only on a cyclic group instead of a cube.
Chapter~\ref{ch:factoring} shows the reduction. The quantum Fourier transform that finds the
period is the subject of Chapter~\ref{ch:fourier}.

\begin{history}{1994--1996 \textperiodcentered{} Santa Fe and Bell Labs}
In May 1994, at a workshop at the Santa Fe Institute, Rolf Landauer of IBM
raised the objection that seemed to end the story: a quantum computer cannot
correct its errors. Every classical technique for fixing errors involves copying
something, and quantum states cannot be copied. Without correction, a
computation of a billion steps needs each step to be accurate to one part in a
billion. In 1995 Shor answered with a code that protects one qubit using nine.
In 1996, Robert Calderbank and Shor, and independently Andrew Steane, generalized
it into a whole family, the CSS codes. Part~\ref{part:qec} tells this story in detail.
\end{history}

In 1996, Lov Grover, also at Bell Labs, published a quantum algorithm for
unstructured search: finding the one marked item among $N$ takes about
$\sqrt{N}$ queries instead of $N$. The gain is only quadratic, much less dramatic
than Shor's. But it applies to almost anything that can be phrased as a search,
which made it the second pillar of the field.

\section{What the sequence shows}

By 1996 every piece of the argument was on the table. Feynman showed that
classical machines struggle with quantum physics. Deutsch showed there is a
universal model that does not. Simon showed that the Fourier transform is the
mechanism behind exponential gains. Shor showed that one of those gains breaks
the cryptography of the internet, which turned an academic curiosity into
something governments fund. And error correction showed that the model can, in
principle, survive noise.

What did not exist in 1996 was a machine. The thirty years since then are the
story of Parts~\ref{part:hardware} and~\ref{part:qec}.

\chapter{What a Computer Can Do}\label{ch:turing}

\begin{etymology}
\textbf{Compute} comes from the Latin \emph{computare}: \emph{com-},
``together'', and \emph{putare}, ``to reckon'' (and, originally, ``to prune'',
to clean something up in order to assess it). To compute was to count by
putting things together; the word only later came to mean processing
information by rules.

\textbf{Complexity} comes from \emph{complexus}: \emph{com-} plus
\emph{plectere}, ``to braid''. Something complex is braided together and hard to
pull apart. In computing it means the cost of solving a problem and how that
cost grows with the size of the input.
\end{etymology}

Saying a quantum computer is ``faster'' means nothing until you say faster than
what, at what, and measured how. That needs a precise model of what a classical
computer is. The model was built in 1936, ten years before the first electronic
computer, to answer a question about the foundations of mathematics.

\section{The Turing machine}

\begin{history}{1928--1936 \textperiodcentered{} Göttingen, Cambridge and Princeton}
In 1928, David Hilbert and Wilhelm Ackermann posed the
\emph{Entscheidungsproblem}, the decision problem: is there a mechanical
procedure that decides, for any statement of first-order logic, whether it is
provable? To answer ``no'', someone had to define what a mechanical procedure
is. Alonzo Church, at Princeton, did it with the lambda calculus in 1936. Alan
Turing, a 23-year-old fellow at King's College, Cambridge, did it the same year
with an imaginary machine, in ``On Computable Numbers''. Both answers were no.
\end{history}

Turing's machine has four parts: an infinite tape divided into cells, a head that
reads and writes one cell at a time, a finite set of internal states, and a
transition table. At each step the table looks at the current state and the
symbol under the head and says three things: which symbol to write, which way to
move, and which state to go to next.

A machine that inverts every bit of a binary string needs one working state and
three rules: on a 0, write 1 and move right; on a 1, write 0 and move right; on a
blank, stop. Fed \texttt{110100}, it walks across the tape and leaves
\texttt{001011}. Nothing about this looks like a computer, and that is the point:
anything a laptop does can be broken down into steps this small.

\begin{concept}{Church--Turing thesis}
Anything that can be computed by a mechanical procedure can be computed by a
Turing machine.

\textbf{What it fixes:} it gives ``computable'' a precise meaning, so you can
prove that some things are not.

\textbf{What it costs:} it is a thesis, not a theorem. It cannot be proved,
because ``mechanical procedure'' is an informal idea. It has held up for ninety
years.
\end{concept}

Turing also showed that one machine can imitate any other if you write the other
machine's table on its tape. That \emph{universal} machine is the idea behind
software: a program is the description of a machine, and the computer is the
universal machine that runs it.

\section{What no machine can do}

Turing's real target was a limit. He proved that no machine can decide, for every
machine $M$ and input $x$, whether $M$ eventually stops on $x$. This is the
\emph{halting problem}. The proof is a diagonal argument, like Cantor's: suppose
such a decider exists, build a machine that asks the decider about itself and
then does the opposite, and you get a contradiction.

So there are well-posed questions that no algorithm can answer, ever. A quantum
computer does not change this. It computes exactly the same functions as a Turing
machine; what it might change is how fast.

\section{Cost: P, NP and their cousins}

Computable is not the same as practical. A problem can have an algorithm that
would take longer than the age of the universe. Complexity theory sorts problems
by how the cost grows with the size $n$ of the input.

\begin{history}{1965--1972 \textperiodcentered{} the birth of P and NP}
In 1965, Jack Edmonds and Alan Cobham, independently, proposed that ``efficient''
should mean ``polynomial time'': the number of steps grows like $n^k$ for some
fixed $k$. In 1971, Stephen Cook showed that a large family of problems that are
easy to check are all equally hard to solve, if any of them is. Richard Karp
listed 21 of them the next year. P versus NP was born, and it is still the most
famous open problem in computer science.
\end{history}

The classes that matter here:

\begin{itemize}
\item \textbf{P}: problems solvable in polynomial time by a deterministic machine.
\item \textbf{NP}: problems where a proposed solution can be checked in polynomial
time. Factoring is in NP: given $p$ and $q$, multiply and compare.
\item \textbf{BPP}: problems solvable in polynomial time by a machine that may flip
coins, with error probability below $1/3$ (repeating the run makes it as small as
you like). Defined by John Gill in 1977.
\item \textbf{BQP}: the same thing for a quantum computer. Defined by Bernstein and
Vazirani in 1993.
\end{itemize}

It is proved that $\mathrm{P} \subseteq \mathrm{BPP} \subseteq \mathrm{BQP}$:
anything a classical machine does efficiently, a quantum machine does too. What
nobody has proved is that any of these inclusions is strict. Factoring is in BQP
thanks to Shor, and nobody knows whether it is in P. It is also not believed to
be NP-complete, which matters: Shor's algorithm says nothing about the hard
problems in Karp's list.

The gap between polynomial and exponential is the whole story. At $n = 50$, $n^2$
is 2500 steps and $2^n$ is about $1.1 \times 10^{15}$.

\section{The quantum Turing machine}

Deutsch's 1985 machine keeps Turing's parts but lets the tape, the head and the
state be in superposition, and replaces the transition table with a unitary
transformation (Chapter~\ref{ch:state}). Nobody programs with it. In 1989
Deutsch also introduced quantum circuits, sequences of gates acting on qubits,
and in 1993 Andrew Yao proved that the two models are equivalent. Circuits are
what everybody uses.

The machine can be in a superposition of many configurations, but a measurement
returns only one of them. The advantage is never ``try everything and read
everything''. It is to arrange the computation so that wrong answers cancel each
other before the measurement and the right one remains. Chapter~\ref{ch:qubit}
shows what that means with one qubit.

\section{What the classes say about quantum speedups}

Shor put factoring in BQP. Grover gave a quadratic speedup for search, and in
1997 Bennett, Bernstein, Brassard and Vazirani proved that quadratic is the best
possible for unstructured search. A quantum computer cannot solve NP-complete
problems by brute force in polynomial time.

So the exponential speedups come from structure, like the periodicity Shor
exploits. Where there is no structure, you get Grover's square root at best.
Whether BQP is really bigger than BPP, outside of oracle models, is still open.
Shor's algorithm is the strongest evidence that it is, but it is evidence, not
proof: nobody has shown that factoring is classically hard.

\chapter{Why You Can Compute with a Qubit}\label{ch:qubit}

\begin{etymology}
\textbf{Abstraction} comes from the Latin \emph{abstrahere}: \emph{ab-},
``away'', and \emph{trahere}, ``to pull''. To abstract is to pull something out
of its context and keep only what matters from one point of view.

\textbf{Superposition} comes from \emph{superponere}: \emph{super-}, ``on top
of'', and \emph{ponere}, ``to put''. Things put on top of each other, occupying
the same place and adding up into one whole.
\end{etymology}

The whole chapter unpacks one sentence:

\begin{quote}
You can compute with qubits because you can encode the candidate answers of a
problem in a superposition, and then apply unitary operations that make the
amplitudes interfere, so that the right answer is favored when you measure.
\end{quote}

Five ideas are in that sentence: the qubit, superposition, interference, unitary
operations and encoding. This chapter takes them one at a time, with as little
mathematics as possible. The next chapters bring the mathematics back.

\section{A qubit is an abstraction}

A qubit is a device with two distinguishable states, which we call $\ket{0}$ and
$\ket{1}$. That is the only thing it shares with a bit. A bit is always one of the
two. A qubit can be in a combination of both, written
\[
  \alpha\ket{0} + \beta\ket{1}.
\]
The numbers $\alpha$ and $\beta$ are \emph{amplitudes}. They are not
probabilities, but they produce them: when you measure, you get $0$ with
probability $|\alpha|^2$ and $1$ with probability $|\beta|^2$, and the two have to
add up to 1.

The usual picture is a coin. On the table, it shows heads or tails. In the air,
you could say it is ``both''. Suppose the coin is weighted so that its state in the
air is $0.866\ket{\text{heads}} + 0.5\ket{\text{tails}}$. Then
$0.866^2 = 0.75$ and $0.5^2 = 0.25$: three times out of four it lands heads. Throw
it a thousand times and the fraction of heads settles near 75\%.

\section{Superposition is not speed}

The coin is a bad picture, and it is worth saying why. A spinning coin is always
either heads or tails at each instant; we just cannot follow it. That is
ignorance, and ignorance is described by ordinary probabilities. A qubit in
superposition is not secretly one of the two. Both components are present at the
same time, and, unlike probabilities, they can cancel each other.

A better image is two waves added together. Take one oscillation and another with
a different frequency, and add them: you get a single wave. Neither of the
originals can be seen in it, but both are in there. This is still an analogy (a
qubit is not a wave in space), so use it only to hold the idea of coexistence.

\section{Interference}

Here is the mechanism, in three steps, with the operation every textbook uses
first. The \emph{Hadamard} operation $H$ acts like this:
\[
  H\ket{0} = \frac{\ket{0} + \ket{1}}{\sqrt{2}}, \qquad
  H\ket{1} = \frac{\ket{0} - \ket{1}}{\sqrt{2}}.
\]
Start from $\ket{0}$ and apply $H$: you get an even superposition, 50/50. Now
flip the sign of the $\ket{1}$ component (a \emph{phase flip}, called $Z$). The
state becomes $(\ket{0} - \ket{1})/\sqrt{2}$. Measure now and you still get 50/50,
because a sign does not change the size of an amplitude. Then apply $H$ again. The
result is $\ket{1}$, with certainty.

Without the phase flip, the same recombination gives $\ket{0}$ with certainty.

\begin{figure}[htbp]
\centering
\begin{tikzpicture}[x=5mm, y=11mm, font=\sffamily\scriptsize]
  \foreach \px/\title/\a/\b in {0/{after $H$}/0.707/0.707, 9/{after the flip}/0.707/-0.707, 18/{after $H$ again}/0/1} {
    \draw[black!40] (\px-0.5,0) -- (\px+5.5,0);
    \fill[hfConcept!70] (\px+0.5,0) rectangle (\px+2,\a);
    \fill[hfHistory!80] (\px+3,0) rectangle (\px+4.5,\b);
    \node[below] at (\px+1.25,{min(0,\a)}) {$\ket{0}$};
    \node[below] at (\px+3.75,{min(0,\b)}) {$\ket{1}$};
    \node[above, font=\sffamily\small] at (\px+2.5,1.25) {\title};
  }
  \node[hf/note] at (2.5,-1.25) {50\% / 50\%};
  \node[hf/note] at (11.5,-1.25) {50\% / 50\%};
  \node[hf/note] at (20.5,-1.25) {0\% / 100\%};
\end{tikzpicture}
\caption{Amplitudes, not probabilities. Flipping the sign of $\ket{1}$ leaves the
probabilities at 50/50, but the second $H$ turns that sign into a certain answer.}
\end{figure}

Why does the sign decide the outcome? Because $H$ adds and subtracts. Applied to
$a\ket{0} + b\ket{1}$, it gives
\[
  \frac{a + b}{\sqrt{2}}\ket{0} + \frac{a - b}{\sqrt{2}}\ket{1}.
\]
When $a$ and $b$ have the same sign, they pile up on $\ket{0}$ and cancel on
$\ket{1}$. When the signs are opposite, the reverse happens. The total
probability is conserved; it is only moved around. This is \emph{interference},
and it is the engine of every quantum algorithm: mark the right answer with a sign
the measurement cannot see, then recombine so the sign turns into probability.

\begin{concept}{Interference}
Amplitudes add like signed (in general, complex) numbers before they become
probabilities. Equal signs reinforce; opposite signs cancel.

\textbf{What it fixes:} an algorithm can steer probability toward the right
answer without ever looking at the intermediate state.

\textbf{What it costs:} the phases must survive the whole computation. Any
contact with the environment scrambles them, and the interference is gone. This
is the reason qubits are so hard to build (Part~\ref{part:hardware}).
\end{concept}

The hard question is how to flip the sign of the right answer without knowing it
in advance. That is what an oracle does, and Chapter~\ref{ch:encoding} returns
to it.

\section{Unitary operations}

$H$ and $Z$ are \emph{unitary} operations. For now, two of their properties
matter. They act on all the components of the superposition at once. And they
are reversible: nothing is thrown away, so you can always undo them. They have to
be. An operation that destroyed information along the way would destroy
amplitudes, and there would be nothing left to interfere. The conservation you
saw above, where 50/50 became 0/100 and the total stayed at 1, is a consequence.
Chapter~\ref{ch:state} says exactly what unitary means.

\section{Encoding}

In the coin, the candidate answers were heads and tails, the two states of the
qubit. A real problem needs more room. One qubit holds 2 candidates, two qubits
hold 4, three hold 8, and $n$ qubits hold $2^n$. Encoding means mapping the
candidates of your problem onto those states, so that all of them can be put into
superposition and take part in the interference. Chapter~\ref{ch:encoding} shows
how it is done for a real problem.

\section{How much does that buy?}

Grover's search finds one marked item among $N = 2^n$ in about $\sqrt{N}$ steps.
A classical machine has to check items one by one; give it 24 cores and it checks
24 at a time.

\begin{center}
\begin{tabular}{rrrrr}
\hline
qubits $n$ & candidates $2^n$ & 24 cores, $N/24$ & Grover, $\sqrt{N}$ & ratio \\
\hline
10 & $1.0 \times 10^{3}$ & 43 & 32 & 1.3 \\
20 & $1.0 \times 10^{6}$ & $4.4 \times 10^{4}$ & $1.0 \times 10^{3}$ & 43 \\
30 & $1.1 \times 10^{9}$ & $4.5 \times 10^{7}$ & $3.3 \times 10^{4}$ & $1.4 \times 10^{3}$ \\
40 & $1.1 \times 10^{12}$ & $4.6 \times 10^{10}$ & $1.0 \times 10^{6}$ & $4.4 \times 10^{4}$ \\
50 & $1.1 \times 10^{15}$ & $4.7 \times 10^{13}$ & $3.4 \times 10^{7}$ & $1.4 \times 10^{6}$ \\
\hline
\end{tabular}
\end{center}

Each extra qubit doubles the classical work and multiplies the quantum work by
only $\sqrt{2}$, so the ratio keeps growing.

The table is also misleading, and you should know how. It compares
\emph{steps}, and a step is not the same thing on both sides. A classical core
checks a candidate in nanoseconds. A Grover step on an error-corrected machine is
a long sequence of logical gates, each of which costs many physical operations
and microseconds of time. In 2021, Ryan Babbush and colleagues at Google worked out
that, with realistic error-correction overheads, a quadratic speedup will not beat
a classical machine on the first generations of fault-tolerant hardware. Square-root
speedups are real, but they are the weakest argument for building quantum
computers.

\chapter{From Newton to Hamilton}\label{ch:classical}

\begin{etymology}
\textbf{Mechanics} comes from the Greek \emph{mēkhanikē}, the art of machines,
from \emph{mēkhanē}, ``machine'' or ``contrivance''. It came to mean the study of
how things move under fixed rules.

\textbf{Energy} comes from the Greek \emph{enérgeia}: \emph{en-}, ``in'', and
\emph{érgon}, ``work''. Something at work, in action.

\textbf{Hamiltonian} is named after William Rowan Hamilton (1805--1865), the Irish
mathematician who rewrote Newton's mechanics in terms of position and momentum.
\end{etymology}

Quantum mechanics was not invented from scratch. It grew out of classical
mechanics at the moment classical mechanics stopped matching experiments, and it
kept a surprising amount of the old structure. To see what changed, you need to
know what was there before.

Classical mechanics was written three times, by Newton, Lagrange and Hamilton.
The three versions say the same thing about the physical world. What changes is
the language, and each new language is more abstract and more powerful than the
last. Hamilton's is the one quantum mechanics kept.

\section{Newton: forces and trajectories}

\begin{history}{1687 \textperiodcentered{} Cambridge}
Isaac Newton published the \emph{Philosophiæ Naturalis Principia Mathematica} in
1687, after Edmond Halley pushed him to write up his work on orbits and paid for
the printing. The three laws of motion and the law of gravitation explained, with
one set of equations, why apples fall and why planets follow ellipses.
\end{history}

The second law, $F = ma$, says that force changes motion. Written out, the
acceleration is the rate of change of the rate of change of the position,
\[
  m\,\frac{d^2x}{dt^2} = F .
\]
Notice what this equation does not say. It does not give the position. It gives a
relation between the position and how fast it changes, and the position itself has to
be found from that relation. An equation of this kind is a \emph{differential
equation}, and Newton had to invent the mathematics to write it.

\begin{history}{1665--1736 \textperiodcentered{} Woolsthorpe, Hanover and London}
In 1665 and 1666, with Cambridge closed by the plague, the young Newton worked out at
his family's farm in Woolsthorpe what he called the \emph{method of fluxions}: a way
of computing rates of change (fluxions) of quantities that flow in time. In a treatise
written around 1671 he sorted the equations that relate a quantity to its fluxions into
types and showed how to solve some of them, but he kept it unpublished; it appeared
only in 1736, after his death. Gottfried Wilhelm Leibniz, in Hanover, found the same
calculus independently, published first in 1684, and gave it the notation $dx/dt$ that
is still used. He also coined the name, \emph{aequatio differentialis}. The quarrel over
who came first poisoned English and Continental mathematics for decades.
\end{history}

\begin{concept}{Differential equation}
An equation that relates an unknown function to its rates of change. The equation
only says how the system changes from one instant to the next. Solving it means
finding the function itself: the one that says where the system will be at every later
moment, how it evolves. That function is the prize, and finding it is the whole point of
writing the equation. When the function depends on a single variable, usually time, it
is an \emph{ordinary} differential equation (ODE): the spring below, and every equation
of motion in this chapter. When it depends on several variables at once, such as
position and time, the rates of change are \emph{partial} derivatives and it is a
\emph{partial} differential equation (PDE): the vibrating string (d'Alembert, 1747), the
heat equation (Fourier, 1807, Chapter~\ref{ch:fourier}) and Schrödinger's equation
(Chapter~\ref{ch:quantum}).

\textbf{What it fixes:} physical laws are local. They say what happens in the next
instant, here, given what is happening now. A differential equation is exactly that
kind of statement, and it predicts the whole future from it.

\textbf{What it costs:} most differential equations have no formula for their
solution. You either find one of the few that do, or you compute the solution step by
step.
\end{concept}

The state of a particle is its position and velocity, $(x, v)$. Give the state at one
instant and the forces, and the differential equation fixes the whole future. This is
classical determinism.

The standard example is a mass on a spring. The spring pulls back with a force
proportional to the displacement, $F = -kx$, so the equation of motion is the ODE
$m\ddot{x} = -kx$ (a dot means a derivative in time, Newton's own notation). Solving it
means finding a function whose second derivative is $-k/m$ times itself, and the
cosine is such a function:
\[
  m\ddot{x} = -kx, \qquad x(t) = A\cos(\omega t + \phi), \qquad
  \omega = \sqrt{k/m}.
\]
The mass oscillates forever with frequency $\omega$. This \emph{harmonic
oscillator} comes back again and again in physics, and it comes back in
quantum computing too: a superconducting qubit is a harmonic oscillator with a
small defect added on purpose (Chapter~\ref{ch:hardware}).

\begin{figure}[htbp]
\centering
\begin{tikzpicture}[x=1mm, y=1mm]
  % equilibrium line
  \draw[black!35, dashed] (41,8) -- (41,-39) node[below, hf/note] {rest position, $x = 0$};
  \foreach \row/\d/\lab in {0/12/a, 1/0/b, 2/-12/c} {
    \pgfmathsetmacro{\y}{-\row*16}
    \pgfmathsetmacro{\bx}{36+\d}
    \pgfmathsetmacro{\seg}{(\bx-6)/9}
    % wall and floor
    \draw[line width=1pt] (0,\y-5) -- (0,\y+5);
    \foreach \h in {-4,-2,...,4} \draw[black!45] (0,\y+\h) -- (-2,\y+\h-2);
    \draw[black!40] (0,\y-5) -- (64,\y-5);
    % spring
    \draw[line width=0.7pt] (0,\y) -- (2,\y);
    \draw[line width=0.7pt, decorate, decoration={coil, aspect=0.5,
      segment length=\seg mm, amplitude=2.2mm}] (2,\y) -- (\bx-3,\y);
    \draw[line width=0.7pt] (\bx-3,\y) -- (\bx,\y);
    % block
    \filldraw[fill=hfConcept!15, draw=hfConcept, line width=0.8pt]
      (\bx,\y-5) rectangle (\bx+10,\y+5);
    \node[font=\small] at (\bx+5,\y) {$m$};
    \node[font=\sffamily\small\bfseries, text=black!60] at (-6,\y) {(\lab)};
  }
  % forces
  \draw[hfAlert, line width=1pt, -{Stealth[length=2.2mm]}] (58,7) -- (47,7)
    node[midway, above, hf/note, text=hfAlert] {$F = -kA$};
  \node[hf/note, anchor=west] at (62,0) {$x = +A$};
  \node[hf/note, anchor=west] at (62,-16) {$x = 0$};
  \draw[hfAlert, line width=1pt, -{Stealth[length=2.2mm]}] (24,-25) -- (35,-25)
    node[midway, above, hf/note, text=hfAlert] {$F = +kA$};
  \node[hf/note, anchor=west] at (62,-32) {$x = -A$};
  % x(t)
  \begin{axis}[at={(82mm,-37mm)}, anchor=south west, width=52mm, height=50mm,
      xmin=0, xmax=1.6, ymin=-1.3, ymax=1.3,
      xtick={0,0.25,0.5,1}, xticklabels={$0$,$\frac{T}{4}$,$\frac{T}{2}$,$T$},
      ytick={-1,0,1}, yticklabels={$-A$,$0$,$+A$},
      xlabel={$t$}, ylabel={$x(t)$},
      tick label style={font=\scriptsize}, label style={font=\small},
      axis lines=middle, axis line style={black!50},
      xlabel style={at={(ticklabel* cs:1.02)}, anchor=west},
      ylabel style={at={(ticklabel* cs:1.02)}, anchor=south}]
    \addplot[hfConcept, line width=0.9pt, domain=0:1.55, samples=120] {cos(360*x)};
    \addplot[only marks, mark=*, mark size=1.6pt, hfAlert] coordinates {(0,1) (0.25,0) (0.5,-1)};
    \node[font=\sffamily\scriptsize\bfseries, text=black!60, anchor=south west] at (axis cs:0.02,1) {(a)};
    \node[font=\sffamily\scriptsize\bfseries, text=black!60, anchor=south east] at (axis cs:0.23,0.02) {(b)};
    \node[font=\sffamily\scriptsize\bfseries, text=black!60, anchor=north west] at (axis cs:0.52,-1) {(c)};
  \end{axis}
\end{tikzpicture}
\caption{A mass on a spring. (a) Stretched by $A$, the mass stops for an instant and
is pulled back with force $-kA$. (b) At the rest position the force is zero, but the
mass is moving at its fastest and overshoots. (c) Compressed by $A$, it stops again and
is pushed out. The position
traces $x(t) = A\cos\omega t$, and one full cycle takes $T = 2\pi/\omega$.}\label{fig:spring}
\end{figure}

Newton's method works well when you can write down every force. It becomes painful
when the system has constraints: a pendulum on a rod, a bead on a wire, a ball rolling
on a surface. What hurts is writing the equations of motion, and the pendulum shows
why.

A mass $m$ hangs from a rigid rod of length $\ell$ and swings. Newton's law works with
forces and with the coordinates $x$ and $y$ of the mass, so you write $F = ma$ once for
each direction. Two forces act: gravity, $mg$, which you know, and the tension $T$ of
the rod, which you do not. The tension is whatever it takes to keep the mass on its
circle, and it changes all the time as the mass swings. With the pivot at the origin,
\[
  m\ddot x = -T\,\frac{x}{\ell}, \qquad
  m\ddot y = -T\,\frac{y}{\ell} - mg, \qquad
  x^2 + y^2 = \ell^2 .
\]
Three unknown functions, $x(t)$, $y(t)$ and $T(t)$, two differential equations and a
constraint. You have to use the constraint to eliminate the tension by hand before you
can solve anything, and the tension is a quantity you never wanted to know. For a
single pendulum this is a page of algebra. For a double pendulum, or a machine with a
dozen joints, it becomes unmanageable.

\begin{figure}[htbp]
\centering
\begin{tikzpicture}[x=1mm, y=1mm]
  \fill[black!15] (-12,0) rectangle (12,2.5);
  \draw[black!60, line width=0.8pt] (-12,0) -- (12,0);
  \draw[black!30, dashed] (0,0) -- (0,-36);
  \draw[black!40] (0,-36) arc (-90:-55:36);
  \draw[black!40] (0,-36) arc (-90:-125:36);
  % rod at theta = 28 degrees
  \coordinate (m) at ({36*sin(28)},{-36*cos(28)});
  \draw[black!70, line width=1.2pt] (0,0) -- (m);
  \fill[black!70] (0,0) circle (0.9);
  \draw[black!60] (0,-10) arc (-90:-62:10);
  \node[font=\small] at (2.6,-12.5) {$\theta$};
  \node[font=\small, left] at ($(0,0)!0.55!(m)$) {$\ell$};
  % tension and weight
  \draw[hfAlert, line width=1pt, -{Stealth[length=2.2mm]}] (m) -- ($(m)!11mm!(0,0)$)
    node[pos=0.8, left=2pt, hf/note, text=hfAlert] {$T$ (unknown)};
  \draw[hfConcept, line width=1pt, -{Stealth[length=2.2mm]}] (m) -- ++(0,-12)
    node[below, hf/note, text=hfConcept] {$mg$};
  \filldraw[fill=hfConcept!15, draw=hfConcept, line width=0.8pt] (m) circle (2.6);
  \node[font=\small] at (m) {$m$};
  % axes
  \draw[black!45, -{Stealth[length=1.8mm]}] (22,-30) -- (30,-30) node[right, hf/note] {$x$};
  \draw[black!45, -{Stealth[length=1.8mm]}] (22,-30) -- (22,-22) node[above, hf/note] {$y$};
  % comparison
  \node[hf/note, anchor=west, align=left, text width=58mm] at (38,-8)
    {\textbf{Newton:} unknowns $x$, $y$ and $T$; two equations plus the constraint
     $x^2 + y^2 = \ell^2$; eliminate $T$ by hand.};
  \node[hf/note, anchor=west, align=left, text width=58mm] at (38,-26)
    {\textbf{Lagrange:} one coordinate, $\theta$; one equation,
     $\ddot\theta = -(g/\ell)\sin\theta$; $T$ never appears.};
\end{tikzpicture}
\caption{A pendulum. Newton's law needs every force, including the tension of the
rod, which is unknown. Lagrange's method uses the angle, which the rod allows to change
freely, and the tension drops out.}\label{fig:pendulum}
\end{figure}

\section{Lagrange: energy instead of forces}

\begin{history}{1788 \textperiodcentered{} Paris}
Joseph-Louis Lagrange published the \emph{Mécanique analytique} in 1788. Its
preface boasts that the book contains no figures: mechanics had become pure
analysis, operations on equations. The book starts from a principle, not from
forces: the principle of virtual work, which Jean d'Alembert had extended to bodies in
motion in 1743. The principle of least action, which Pierre-Louis Maupertuis had
stated vaguely in 1744 and Leonhard Euler and Lagrange had made precise, comes out of
it in the book as a consequence.
\end{history}

Newton and Lagrange looked at the same pendulum and saw different things. Newton saw a
body driven by forces: list every force, add them up, and the sum gives the
acceleration. Lagrange saw that, under those same forces, the motion obeys a
principle: of all the ways the system could go from here to there, it takes the one
that satisfies a single condition. Newton described the system with the forces, and
Lagrange described it with the principle. The forces do not vanish in Lagrange's
version; they are folded into the energies. Define the \emph{Lagrangian} as kinetic
energy minus potential energy, $L = T - V$. The path the system actually takes is
the one that makes the \emph{action} $S = \int L\,dt$ stationary.

\begin{concept}{Principle of least action}
Fix where the system starts and where it ends, and when. Between those two moments it
could, in imagination, take any path at all. For each imaginable path, add up $T - V$
along the way; that sum is the action,
\[
  S = \int_{t_1}^{t_2} (T - V)\,dt .
\]
The path nature actually takes is the one where $S$ is \emph{stationary}: nudge the
path slightly in any direction and, to first order, $S$ does not change. Usually that
is a minimum, which is where the word \emph{least} comes from, but strictly it is only
stationary.

\begin{center}
\begin{tikzpicture}[x=1mm, y=1mm]
  \draw[black!45, -{Stealth[length=1.8mm]}] (0,0) -- (62,0) node[right, hf/note] {$t$};
  \draw[black!45, -{Stealth[length=1.8mm]}] (0,0) -- (0,30) node[above, hf/note] {$q$};
  \coordinate (A) at (6,6);
  \coordinate (B) at (54,20);
  \draw[black!35, dashed, line width=0.6pt] (A) .. controls (20,30) and (35,-4) .. (B);
  \draw[black!35, dashed, line width=0.6pt] (A) .. controls (18,-6) and (40,34) .. (B);
  \draw[black!35, dashed, line width=0.6pt] (A) .. controls (30,8) and (34,6) .. (B);
  \draw[hfConcept, line width=1.2pt] (A) .. controls (22,22) and (38,26) .. (B);
  \fill[black!70] (A) circle (1) node[above left, hf/note] {start, $t_1$};
  \fill[black!70] (B) circle (1) node[below right, hf/note] {end, $t_2$};
  \node[hf/note, text=hfConcept, anchor=south] at (30,25) {the path taken: $S$ stationary};
  \node[hf/note, anchor=west] at (64,9) {other imaginable paths};
  \draw[black!35] (63.5,9) -- (45,9.5);
\end{tikzpicture}
\end{center}

Pierre-Louis Maupertuis proposed in 1744 that nature always spends the least
``action'', a vague quantity he partly justified by theology. Euler made it precise the
same year for a single particle, and Lagrange, in 1760, extended it to any system.

\textbf{What it fixes:} instead of following forces instant by instant, you compare whole
histories using one number. That number is the same in any coordinates, so you can
choose the convenient ones, and forces that do no work never appear.

\textbf{What it costs:} it is stated with both ends fixed, as if the system knew where it
was going. It does not; the Euler--Lagrange equations turn it back into a rule for the
next instant. And it only applies, in this simple form, to systems that do not lose
energy to friction.
\end{concept}

From that one principle come the Euler--Lagrange equations:
\[
  \frac{d}{dt}\frac{\partial L}{\partial \dot{q}} - \frac{\partial L}{\partial q} = 0 .
\]
Read it as a balance. $\partial L/\partial\dot q$ is how much $L$ depends on the
velocity; for the usual kinetic energy it is the momentum, and the first term is how
fast that momentum changes. $\partial L/\partial q$ is how much $L$ depends on the
position; for the usual potential energy it is the force. The equation says that at
every instant the two are equal: the momentum changes exactly as fast as the force
pushes. That is Newton's law again, but $q$ can now be any quantity you like, a
distance or an angle, and choosing it well is the whole trick. It is an ordinary
differential equation, one for each coordinate, because the only variable is time.

Nothing in it says the system has to be mechanical. All it needs is a Lagrangian: a
function of some quantity and of its rate of change, whose integral the system makes
stationary. In an electric circuit the quantity is the charge on the capacitor. The
energy in the coil plays the part of kinetic energy, the energy in the capacitor
plays the part of potential energy, and the same equation gives the circuit's
equation; Chapter~\ref{ch:hardware} starts from the energy of such a circuit to
describe a superconducting qubit. For a field, such as the electromagnetic field,
there is one coordinate at every point of space, the equation becomes a partial
differential equation, and out come Maxwell's equations. David Hilbert obtained the
equations of general relativity this way in 1915, and the Standard Model of particle
physics is written today as a single Lagrangian. What does not fit, in this simple
form, is friction: a system that loses energy has no action to make stationary.

For the pendulum of Figure~\ref{fig:pendulum}, Lagrange chooses the angle $\theta$
instead of $x$ and $y$, and two things disappear. The rod's constraint disappears,
because any angle keeps the bob at distance $\ell$. The tension disappears, because
it pulls along the rod while the bob moves across it, so it does no work and never
enters the energies. What is left is a single equation:
\[
  \ddot\theta = -\frac{g}{\ell}\sin\theta .
\]
For small swings it is the spring's equation, and its solution is again a cosine,
$\theta(t) = \theta_0\cos\bigl(\sqrt{g/\ell}\,t\bigr)$: the function that says where
the pendulum is at every moment.

\section{Hamilton: position and momentum}

\begin{history}{1827--1835 \textperiodcentered{} Dublin}
William Rowan Hamilton was made Astronomer Royal of Ireland in 1827, at 21, while
still an undergraduate. His first big work was not about mechanics at all but about
light. In 1827 he read his \emph{Theory of Systems of Rays} to the Royal Irish
Academy: a whole family of light rays, passing through lenses and mirrors, could be
described by a single function. Optics already had its own principle of this kind,
Pierre de Fermat's from 1662: light goes from one point to another along the path
that takes the least time. Hamilton noticed that mechanics, with Maupertuis's least
action, had the same shape, and in two papers of 1834 and 1835 he carried his optical
method over to moving bodies.
\end{history}

So Hamilton was not fixing a failure of Lagrange. Lagrange's equations solved every
problem of mechanics that Hamilton's would solve, and in the nineteenth century
nobody computed an orbit more easily because of him. What Hamilton was after was the
analogy: one theory for rays of light and for the paths of bodies. Following it, he
ended up with a different way of writing the state of a system, and that turned out
to matter more than the analogy itself.

In Lagrange, the unknown is the position $q(t)$, and the velocity is only its
derivative. To know where the system goes next you need both, so you solve one
second-order equation. Hamilton promoted the momentum to an unknown of its own,
with the same standing as the position. The price is twice as many equations; the
gain is that each one is first-order, and that the pair $(q, p)$ at one instant is
the whole state of the system, nothing missing.

\begin{concept}{Momentum}
Momentum is mass times velocity, $p = mv$, and it measures how much motion a body
carries: a truck and a bicycle at the same speed do not stop the same way. René
Descartes called it the quantity of motion in 1644, but he multiplied size by speed
and ignored direction, and his rules for collisions came out wrong. In 1668 and 1669,
answering a call from the Royal Society, John Wallis, Christopher Wren and Christiaan
Huygens showed that direction counts: momenta add as arrows, and in a collision their
total does not change. Newton put it at the center of his second law in 1687: force is
the rate at which momentum changes.

Since then the idea has stretched far beyond mass times velocity. In Lagrange and
Hamilton the momentum of a coordinate is $p = \partial L/\partial\dot q$, whatever the
coordinate is; for an angle it is the angular momentum. Light has no mass and still
carries momentum, $p = E/c$: James Clerk Maxwell predicted in 1873 that light pushes on
what it hits, and Pyotr Lebedev measured the push in 1900. In 1924 Louis de Broglie
went the other way and gave every particle a wave, with a wavelength set by its
momentum, $\lambda = h/p$. In 1918 Emmy Noether showed what all of these have in
common: momentum is the quantity that is conserved because the laws of physics are
the same here as one meter to the left. In quantum mechanics it becomes the operator
that moves a state through space.

\textbf{What it fixes:} it gives one conserved quantity to follow through any
collision, and in Hamilton's mechanics it completes the state: position and momentum
together fix the whole future.

\textbf{What it costs:} the picture of mass times velocity stops working quickly. The
momentum of an angle has different units, the momentum of light has no mass in it,
and in quantum mechanics position and momentum can no longer both be sharp.
\end{concept}

Write $q$ for the coordinate, whatever it is, and $p$ for its momentum,
$p = \partial L/\partial\dot q$. For a mass on a spring $p = m\dot x$; for a pendulum
it is the angular momentum $m\ell^2\dot\theta$. Hamilton then built a new function,
the \emph{Hamiltonian}:
\[
  H(q, p) = p\dot{q} - L .
\]
Read it this way: the Hamiltonian of $q$ and $p$ is the total energy of the system,
written as a function of where the system is and how much momentum it carries. The
Lagrangian is $T - V$, written with position and velocity; the Hamiltonian is
$T + V$, written with position and momentum. The term $p\dot q$ is only the
bookkeeping that makes the swap. For the spring, $p\dot x = m\dot x^2 = 2T$, so
$H = 2T - (T - V) = T + V$. Swapping a variable for a slope in this way is older than
Hamilton: Adrien-Marie Legendre used it in 1787, and it is called the Legendre
transform.

The motion becomes two first-order equations, symmetric in $q$ and $p$:
\[
  \dot q = \frac{\partial H}{\partial p}, \qquad
  \dot p = -\frac{\partial H}{\partial q}.
\]
The first says the position changes at the rate set by how the energy depends on the
momentum; for the spring, $\partial H/\partial p = p/m$, which is the velocity. The
second says the momentum changes at the rate set by how the energy depends on the
position, with a minus sign; for the spring, $-\partial H/\partial x = -kx$, which is
the force. The minus sign is what pushes the system downhill in potential energy.

The state of the system is a point $(q, p)$ in \emph{phase space}, and time moves
that point along a flow. For the spring, $H = p^2/2m + kx^2/2$. Energy is
conserved, so each trajectory stays on a curve of constant $H$, and for this $H$ the
curve is an ellipse.

\begin{figure}[htbp]
\centering
\begin{tikzpicture}[scale=0.9]
  \draw[black!40, -{Stealth[length=2mm]}] (-3.4,0) -- (3.6,0) node[right, font=\sffamily\small, text=black] {$x$};
  \draw[black!40, -{Stealth[length=2mm]}] (0,-2.6) -- (0,2.8) node[above, font=\sffamily\small, text=black] {$p$};
  \foreach \r/\c in {0.7/hfConcept, 1.2/hfConcept, 1.7/hfBus, 2.2/hfHistory} {
    \draw[\c, line width=0.9pt] (0,0) ellipse ({1.35*\r} and {\r});
  }
  \draw[hfHistory, line width=0.9pt, -{Stealth[length=2.4mm]}] (0.05,2.2) -- (0.3,2.18);
  \draw[black!35] (3.05,1.75) -- (2.3,1.45);
  \node[hf/note, anchor=west] at (3.1,1.75) {higher energy};
  \draw[black!35] (3.05,-1.75) -- (0.72,-0.5);
  \node[hf/note, anchor=west] at (3.1,-1.75) {lower energy};
  \fill[hfHistory] (2.97,0) circle (2.2pt) node[above right, hf/note, text=hfHistory] {(a)};
  \fill[hfHistory] (0,2.2) circle (2.2pt) node[above left, hf/note, text=hfHistory] {(b)};
  \fill[hfHistory] (-2.97,0) circle (2.2pt) node[above left, hf/note, text=hfHistory] {(c)};
\end{tikzpicture}
\caption{Phase space of the harmonic oscillator. Each ellipse is one trajectory,
at one fixed energy, and the system goes around its ellipse forever, clockwise. The
points (a), (b) and (c) are the three moments of Figure~\ref{fig:spring}: stretched
and at rest, passing the rest position at full speed, compressed and at
rest.}\label{fig:phase}
\end{figure}

\begin{reflection}{Is a point of this diagram an energy, or a rate of change?}
Neither. A point of the diagram is a state: one position and one momentum at the
same instant. The point (a) is the spring stretched with the mass at rest; the point
(b) is the mass passing the rest position at full speed.

The energy is not on either axis. It is a number attached to each point, $H(x, p)$,
the way altitude is attached to each point of a map, and the ellipses are the
contour lines of that map. Every point on one ellipse has the same energy, and points
on bigger ellipses have more. At (a) all of it is potential energy, because $p = 0$;
at (b) all of it is kinetic, because $x = 0$; in between it is shared, and the sum
stays the same.

The rate of change is a third thing: how fast, and in which direction, the point
moves. Hamilton's equations give it, $\dot x = \partial H/\partial p$ and
$\dot p = -\partial H/\partial x$, and those are the slopes of the energy map turned
by a right angle. Someone who always walks at right angles to the slope of a hill
follows a contour line and never climbs or descends; that is why the system stays on
its ellipse. So the energy is not the rate of change. The energy is the height of the
map, and the rate of change comes from how steep the map is.
\end{reflection}

Hamilton's form has one more piece that matters later. For any two quantities
$f$ and $g$ of the system, define the \emph{Poisson bracket}, introduced by
Siméon Denis Poisson in 1809:
\[
  \{f, g\} = \frac{\partial f}{\partial q}\frac{\partial g}{\partial p}
           - \frac{\partial f}{\partial p}\frac{\partial g}{\partial q}.
\]
Then the time evolution of any quantity is $\dot f = \{f, H\}$, and position and
momentum satisfy $\{q, p\} = 1$.

What the new form bought was structure, and it paid off over the next ninety years.
Conserved quantities show up on sight: if $H$ does not depend on some $q$, then
$\dot p = -\partial H/\partial q = 0$ and that momentum never changes. Carl Jacobi took
the method up in 1837 and made it the standard tool of celestial mechanics. The flow
in phase space keeps areas fixed, a fact that goes back to Joseph Liouville in 1838,
and Josiah Willard Gibbs built statistical mechanics on it in 1902. In 1926 Erwin
Schrödinger followed Hamilton's own analogy one step further: if the paths of bodies
are to mechanics what rays are to optics, and rays are only an approximation to light
waves, then there should be waves behind the paths of bodies too. And the Poisson
bracket is what Dirac would turn into the quantum commutator.

\begin{question}{What are quaternions, and why do they come up with Hamilton?}
Hamilton is as famous for quaternions as for mechanics. Complex numbers, $a + bi$,
describe rotations in the plane: multiplying by $i$ turns a point by $90^\circ$.
Hamilton spent years looking for numbers that would do the same for rotations in
space. On 16 October 1843, walking with his wife along the Royal Canal in Dublin, he
saw that it needed three imaginary units, not two, and carved the rule into the stone
of Broom Bridge:
\[
  i^2 = j^2 = k^2 = ijk = -1 .
\]
The carving is gone; a plaque marks the place. A quaternion is
$a + bi + cj + dk$, and it does not commute: $ij = k$ but $ji = -k$. It was the first
system of numbers in which the order of a product matters, more than eighty years
before Heisenberg's $xp \neq px$.

They come up in this book for a precise reason. Quaternions of length one describe
every rotation in space, and Chapter~\ref{ch:state} shows that every gate on one
qubit is a rotation of the Bloch sphere. The Pauli matrices, multiplied by $-i$,
obey exactly Hamilton's rule. A single-qubit gate is, up to a phase, a unit
quaternion. That is also why video games and spacecraft use quaternions to store
orientations.
\end{question}

\section{The bridge}

\begin{history}{1925 \textperiodcentered{} Cambridge}
In the autumn of 1925, Paul Dirac, a 23-year-old research student, was reading
Heisenberg's new paper, in which physical quantities no longer commuted: $xp$ was
not equal to $px$. On a Sunday walk he remembered something about Poisson
brackets, but not the details. The libraries were closed, so he waited
impatiently until Monday morning to check. The bracket had the same algebraic
structure as Heisenberg's commutator.
\end{history}

Dirac's rule, now called \emph{canonical quantization}, is a translation table:

\begin{center}
\begin{tabular}{ll}
\hline
classical & quantum \\
\hline
quantities $q$, $p$ & operators $\hat q$, $\hat p$ \\
$\{q, p\} = 1$ & $[\hat q, \hat p] = \hat q\hat p - \hat p\hat q = i\hbar$ \\
$\dot f = \{f, H\}$ & $i\hbar\,\dot{\hat A} = [\hat A, \hat H]$ \\
Hamiltonian $H(q, p)$ & Hamiltonian operator $\hat H(\hat q, \hat p)$ \\
\hline
\end{tabular}
\end{center}

The geometry of phase space does not survive quantization; there are no more
trajectories. The algebra does. The rule is not rigorous for every system, but it
works for every system in this book.

\begin{concept}{The Hamiltonian}
The function, and later the operator, that gives the total energy of a system and
generates its evolution in time.

\textbf{What it fixes:} it puts all the dynamics of a system into a single object,
in a form that survives the move from classical to quantum.

\textbf{What it costs:} you have to know the energy of every interaction in the
system, and for a molecule that is already most of the problem.
\end{concept}

This is why the Hamiltonian is everywhere in quantum algorithms. Hamiltonian
simulation implements $e^{-i\hat H t}$ as a circuit, which is Feynman's original
proposal. Phase estimation extracts the eigenvalues of $\hat H$, which are energies.
Variational algorithms look for the state that minimizes the energy
$\bra{\psi}\hat H\ket{\psi}$. Without the road from Newton to Hamilton, none of
these makes sense.

\chapter{From Planck to Dirac}\label{ch:quantum}

\begin{etymology}
\textbf{Quantum} is Latin for ``how much''. Planck used it for the smallest
amount of energy an oscillator can emit or absorb: a packet that cannot be
divided.

\textbf{Spectrum} is Latin for ``appearance'' or ``image'', from \emph{specere},
``to look''. Newton used it for the band of colors a prism makes from white
light. In quantum mechanics, the spectrum of an operator is the set of values a
measurement can return: what the system can show when you look.

\textbf{Uncertainty} comes from \emph{incertus}: \emph{in-}, ``not'', and
\emph{certus}, ``settled'', from \emph{cernere}, ``to separate''. Something that
cannot be sharply separated from what surrounds it.
\end{etymology}

Nobody set out to invent quantum mechanics. Classical physics was forced to give
way by experiments it could not explain, and the fixes proposed between 1900 and
1927 converged on a formalism that Dirac and von Neumann put in its final form.
That formalism is the native language of the qubit.

\section{Three experiments classical physics could not explain}

\begin{history}{1900 \textperiodcentered{} Berlin}
A hot body glows, and the color of the glow depends only on its temperature.
Classical physics predicted that a perfect absorber, a \emph{black body}, would
radiate more and more energy at higher frequencies, without limit; Paul
Ehrenfest later named this the \emph{ultraviolet catastrophe}. Experiments
showed the opposite: the intensity peaks and falls off. On 14 December 1900, Max
Planck presented a formula that fitted the data perfectly. To derive it, he had to
assume that energy is emitted in packets of size $E = h\nu$, proportional to the
frequency. In a letter of 1931 he called it ``an act of desperation''. He did not
believe the packets were real.
\end{history}

\begin{history}{1905 \textperiodcentered{} Bern}
Light shining on a metal knocks electrons out of it, but only if its frequency is
above a threshold; below it, no intensity is enough. Philipp Lenard had measured
this in 1902, and classical waves could not explain it. Albert Einstein, then a
patent clerk, proposed that light itself comes in packets of energy $h\nu$, later
called photons. One photon either has enough energy to free an electron or it
does not. Robert Millikan spent ten years trying to prove Einstein wrong and ended
up confirming him in 1916. Einstein's Nobel Prize of 1921 was for this, not for
relativity.
\end{history}

But Thomas Young had shown in 1801 that light makes interference patterns, like a
wave. Light was behaving as a wave in some experiments and as a particle in others.

\begin{history}{1913 \textperiodcentered{} Manchester and Copenhagen}
Atoms heated in a flame emit light only at specific frequencies, a barcode of
sharp lines unique to each element. Niels Bohr proposed that the electron in a
hydrogen atom can only occupy certain orbits, and emits a photon when it jumps
from one to another, with energy equal to the difference. The model reproduced
the hydrogen lines exactly, and nobody, Bohr included, could say why the orbits
were allowed.
\end{history}

The three cases point the same way: energy in the microscopic world comes in
discrete steps. A theory was needed in which discreteness is built in.

\section{Waves of matter}

In 1924 Louis de Broglie argued in his doctoral thesis that if light, a wave,
behaves like particles, then particles might behave like waves, with wavelength
$\lambda = h/p$. Clinton Davisson and Lester Germer saw electrons diffract off a
nickel crystal in 1927.

Two versions of the new mechanics appeared within months of each other. In 1925
Werner Heisenberg, recovering from hay fever on the island of Helgoland, built a
mechanics out of tables of numbers that Max Born recognized as matrices. In 1926
Erwin Schrödinger wrote a wave equation for a function $\psi(x, t)$:
\[
  i\hbar\,\frac{\partial \psi}{\partial t} = \hat H\psi
  = \Big(-\frac{\hbar^2}{2m}\nabla^2 + V(x)\Big)\psi .
\]
The two looked nothing alike. Schrödinger himself showed in 1926 that they give
the same predictions.

Put a particle in a box of width $L$ with impenetrable walls. The solutions of
Schrödinger's equation are sine waves that fit exactly in the box, and only some
energies are allowed:
\[
  \psi_n(x) = \sqrt{2/L}\,\sin(n\pi x / L), \qquad
  E_n = \frac{n^2 \pi^2 \hbar^2}{2 m L^2}, \qquad n = 1, 2, 3, \dots
\]
Discreteness is no longer a postulate. It comes out of the equation, from the
requirement that the wave fit.

\begin{history}{1926 \textperiodcentered{} Göttingen}
What is $\psi$? Schrödinger thought it was a real wave of matter. Max Born, in a
paper on collisions, proposed instead that it gives probabilities. The famous
rule (the probability is the square of the amplitude) appears in a footnote
added while the paper was in proof. The main text had said the probability was
proportional to $\psi$ itself. Born got the Nobel Prize for it in 1954.
\end{history}

\section{Dirac and von Neumann}

Paul Dirac's \emph{Principles of Quantum Mechanics} (1930) unified matrices and
waves as two descriptions of the same abstract objects: states and operators.
John von Neumann's \emph{Mathematical Foundations of Quantum Mechanics} (1932)
gave it the rigorous form still used today, built on Hilbert spaces
(Chapter~\ref{ch:hilbert}). The bra--ket notation, $\ket{\psi}$ for a state and
$\bra{\phi}$ for its dual, came later, in a short paper Dirac published in 1939.

In that language, quantum mechanics fits in five postulates:

\begin{enumerate}
\item The state of a system is a unit vector $\ket{\psi}$ in a Hilbert space.
\item Every measurable quantity is a Hermitian operator.
\item The possible results of a measurement are the eigenvalues of that operator.
\item After a measurement that returns the eigenvalue $\lambda_k$, the state
becomes the corresponding eigenvector $\ket{\lambda_k}$.
\item Between measurements, the state evolves as
$\ket{\psi(t)} = e^{-i\hat H t/\hbar}\ket{\psi(0)}$.
\end{enumerate}

Every quantum algorithm is these postulates at work. Preparing $\ket{0}$ on every
qubit is preparing a known state. Applying gates is applying
$e^{-i\hat Ht}$ for carefully chosen Hamiltonians. Reading the result is a
measurement. Chapter~\ref{ch:state} goes through them again with one qubit.

\section{The uncertainty principle}

\begin{history}{1927 \textperiodcentered{} Copenhagen}
Heisenberg's first explanation of uncertainty was a thought experiment: to see
where an electron is, you hit it with light, and the light kicks it, disturbing
its momentum. The same year, Earle Kennard proved the precise inequality
$\Delta x\,\Delta p \geq \hbar/2$, where $\Delta$ is the standard deviation over
many identically prepared systems. In 1929 Howard Percy Robertson generalized it to
any pair of quantities.
\end{history}

Heisenberg's microscope is the wrong explanation, and it is still the one most
people learn. The disturbance story suggests that the particle has a definite
position and momentum and that measuring spoils one of them. It does not have
them. The real reason is that position and momentum are operators that do not
commute:
\[
  [\hat x, \hat p] = i\hbar .
\]
Two operators that do not commute have no common set of eigenvectors, so there is
no state in which both have a definite value. For any two quantities $A$ and $B$,
Robertson's inequality says
\[
  \Delta A\,\Delta B \geq \tfrac12 \big|\langle [\hat A, \hat B] \rangle\big| .
\]
A Gaussian wave packet reaches the minimum exactly: squeeze it in position and it
spreads in momentum, and the product stays at $\hbar/2$. This is a statement
about the geometry of the state space, not about instruments.

\section{What this means for a qubit}

\textbf{Measurement destroys superposition.} Measuring forces the system into an
eigenstate of what you measured. You gain a definite value and lose the
superposition you had.

\textbf{Superposition is not ignorance.} Before a measurement, the qubit is not in
a definite state that we happen not to know. In 1964 John Bell turned this into a
testable inequality; Alain Aspect's experiments of 1982 and many since have
confirmed that no theory of hidden definite values reproduces quantum
mechanics. Aspect, John Clauser and Anton Zeilinger shared the 2022 Nobel Prize
for these tests. This is what separates a quantum computer from a classical
machine flipping coins.

\textbf{Unknown states cannot be copied.}

\begin{history}{1982 \textperiodcentered{} FLASH}
Nick Herbert proposed a device he called FLASH that would use quantum
correlations to send signals faster than light. The scheme needed to make many
copies of an unknown quantum state. Refuting it, William Wootters and Wojciech
Zurek, and independently Dennis Dieks, proved in 1982 that no physical process
can copy an arbitrary unknown state.
\end{history}

\begin{concept}{No-cloning theorem}
There is no operation that takes an unknown state $\ket{\psi}$ and a blank qubit
and produces $\ket{\psi}\ket{\psi}$. The proof uses only linearity: a machine that
copied $\ket{0}$ and $\ket{1}$ correctly would turn their superposition into an
entangled state, not into two copies.

\textbf{What it fixes:} it is why quantum key distribution is secure, and why
Herbert's machine cannot work.

\textbf{What it costs:} every classical way of protecting data against errors
involves copies. Quantum error correction had to be invented from scratch
(Part~\ref{part:qec}).
\end{concept}

No-cloning is often presented as a consequence of the uncertainty principle. It
is better to see it as a consequence of linearity. The two are consistent (if you
could clone, you could measure position on one copy and momentum on another), but
the proof needs nothing more than the fact that quantum evolution is linear.

\chapter{Vectors, Spaces and Hilbert Space}\label{ch:hilbert}

\begin{etymology}
\textbf{Algebra} comes from the Arabic \emph{al-jabr}, ``restoration''. It is a
word from the title of al-Khwarizmi's treatise on equations, written in Baghdad
around 820: \emph{al-jabr} was the operation of moving a subtracted term to the
other side of an equation, restoring the balance.

\textbf{Vector} is Latin for ``carrier'', from \emph{vehere}, ``to carry''.
Something that carries you from one point to another.

\textbf{Space} comes from \emph{spatium}, ``extent'', ``interval''. The room in
which things can exist and move.

\textbf{Phase} comes from the Greek \emph{phásis}, ``appearance''. How something
shows itself at a given moment; the moon has phases. For anything that repeats,
it means where in the cycle it is.
\end{etymology}

The state of a qubit is a vector. Before that sentence can mean anything, you need
to know what a vector is, which turns out to be much more general than an arrow,
and what extra structure a space of vectors needs to talk about probabilities.
The answer is a Hilbert space, and getting there took about a thousand years of
algebra.

\section{From equations to structures}

\begin{history}{around 820 \textperiodcentered{} Baghdad}
Muhammad ibn Musa al-Khwarizmi wrote a treatise on solving equations of the
second degree, without symbols, in words: ``a square and ten roots are equal to
thirty-nine''. Today that is $x^2 + 10x = 39$. The answer was less important than
the method. Al-Khwarizmi noticed that the same two operations, moving terms across
(\emph{al-jabr}) and combining like terms (\emph{al-muqabala}), solve every
equation of a given type, whatever the numbers are.
\end{history}

That is the seed of algebraic thinking: the structure of the manipulation
matters more than the particular numbers. Symbols came slowly. François Viète,
around 1591, was the first to use letters systematically for unknown quantities.
René Descartes, in 1637, fixed the convention still used: letters from the end of
the alphabet ($x, y, z$) for unknowns, from the start ($a, b, c$) for known
quantities.

\begin{history}{30 May 1832 \textperiodcentered{} Paris}
Évariste Galois, 20 years old, spent the night before a duel writing a letter that
summarized his work on why equations of degree five have no general solution by
radicals. The answer, he had found, depends on the symmetries of the roots, not
on the roots themselves. To describe those symmetries he used what we now call
groups. He died of his wound the next day. His papers were published in 1846.
\end{history}

After Galois, algebra stopped being about finding numbers and became the study of
structures: a set, some operations, and the rules they obey. Groups have one
operation. Fields, like the rationals, the reals and the complex numbers, have
addition and multiplication with all inverses; Richard Dedekind introduced them
in 1871 and Heinrich Weber gave the abstract definition in 1893. And
\emph{vector spaces} have addition plus multiplication by numbers taken from a
field.

\section{What a vector is}

Vectors first appeared in physics, for quantities with size and direction: force,
velocity. Hermann Grassmann, a schoolteacher in Stettin, built a general calculus
of directed quantities in 1844; almost nobody read it. Josiah Willard Gibbs and
Oliver Heaviside turned vectors into the everyday tool of physics in the 1880s.
Giuseppe Peano, in 1888, gave the axiomatic definition, and with it the key
step: what makes something a vector is not that it is an arrow, but the rules it
obeys.

\begin{concept}{Vector space}
A set of objects (vectors) that can be added to each other and multiplied by
numbers (scalars), where the result is always in the set and the usual rules
hold: addition is commutative and associative, there is a zero vector, every
vector has an opposite, and scaling distributes over addition.

\textbf{What it fixes:} any theorem proved for vector spaces applies at once to
arrows, lists of numbers, polynomials, functions and quantum states.

\textbf{What it costs:} the intuition of arrows in the plane stops being reliable.
You have to check the rules, not the picture.
\end{concept}

Arrows in the plane are vectors. So are lists of $n$ real numbers, polynomials of
degree at most $n$, continuous functions, and lists of $n$ complex numbers, which
form the space $\mathbb{C}^n$.

The rule that is easiest to forget is \emph{closure}: adding two vectors, or
scaling one, must give something inside the set. Points of the unit disk are not
a vector space: $(0.8, 0) + (0.8, 0) = (1.6, 0)$ falls outside. The positive
quadrant, where both coordinates are positive, is closed under addition (two
positive numbers add to a positive number) but not under scaling: multiply
$(1, 2)$ by $-1$ and you leave it. Without closure, none of the theorems of linear
algebra apply.

\section{Complex numbers}

\begin{history}{1545--1572 \textperiodcentered{} Milan and Bologna}
Girolamo Cardano published the solution of the cubic equation in the \emph{Ars
Magna} of 1545. For some equations with three perfectly real solutions, the
formula passes through square roots of negative numbers along the way. Cardano
called these quantities ``as subtle as they are useless''. Rafael Bombelli, in his
\emph{Algebra} of 1572, worked out rules to compute with them, and showed that
they cancel out and leave the real answer. Complex numbers were not invented to
solve $x^2 = -1$. They were forced on mathematicians by equations that had real
solutions.
\end{history}

For two centuries, these ``imaginary'' numbers were used without being
understood. Caspar Wessel in 1799, Jean-Robert Argand in 1806 and Carl Friedrich
Gauss in 1831 gave them a geometric meaning: $z = a + ib$ is a point in the plane,
with $a$ on the horizontal axis and $b$ on the vertical one. Leonhard Euler had
already found, in 1748, the identity that makes the picture work:
\[
  e^{i\theta} = \cos\theta + i\sin\theta .
\]
So every complex number can be written as $z = re^{i\theta}$, with a size $r$ and
an angle $\theta$, the \emph{phase}. Multiplying by $e^{i\theta}$ rotates by
$\theta$; multiplying by $i$ rotates by 90 degrees. The algebra of complex numbers
was built so that multiplying means rotating.

In $\mathbb{R}^n$, scaling a vector can only stretch it or flip its sign. In
$\mathbb{C}^n$, scaling can also rotate it. A real number has two directions,
positive and negative; a complex number has a whole circle of them.

Angles exist in $\mathbb{R}^2$ too, of course, but there they are something you
compute between two vectors. In $\mathbb{C}$ the angle is part of the number
itself, and it travels with every scalar you use.

You can describe oscillations with real numbers only: $A\cos(\omega t + \theta)$
has an amplitude, a frequency and a phase. But shifting a phase or adding two
oscillations then needs trigonometric identities, and the bookkeeping grows fast.
Written as $Ae^{i(\omega t + \theta)}$, an oscillation is a rotating point, and
shifting its phase is one multiplication. Real numbers can do the job; complex
numbers do it with the structure built in. Quantum amplitudes oscillate, and
interference depends on their phases, which is why they are complex.

\section{Hilbert space}

\begin{history}{1904--1929 \textperiodcentered{} Göttingen and Berlin}
Between 1904 and 1910, David Hilbert worked on integral equations, which appear
all over physics. They could not be handled with finite lists of numbers: the
unknowns were functions, which live in spaces of infinite dimension. Hilbert
showed that length, angle and perpendicularity still make sense there, through
infinite sequences whose squares add up to a finite number. In 1929 John von
Neumann, who had worked with Hilbert in Göttingen, gave the abstract definition
and the name: Hilbert space. He built quantum mechanics on it three years later.
\end{history}

A Hilbert space is a vector space over $\mathbb{C}$ with an \emph{inner product}
$\braket{u|v}$, a way of combining two vectors into one number, that satisfies
four conditions:

\begin{enumerate}
\item it is linear in the second argument;
\item swapping the arguments gives the complex conjugate,
$\braket{u|v} = \overline{\braket{v|u}}$;
\item $\braket{v|v} \geq 0$, and it is zero only for the zero vector;
\item the space is \emph{complete}: every Cauchy sequence converges to a point
inside the space.
\end{enumerate}

A Cauchy sequence, named after Augustin-Louis Cauchy, is one whose terms get
closer and closer to each other, like $1 + 1/n$. Completeness guarantees that
the point they are approaching exists in the space. The rationals are not
complete: a sequence of fractions can approach $\sqrt{2}$, which is not a fraction.
In the finite spaces of quantum computing completeness is automatic, but it is
what makes the infinite-dimensional spaces of physics usable.

In $\mathbb{C}^n$ the inner product is
$\braket{u|v} = \sum_i \overline{u_i}\,v_i$. The conjugation on the first vector is
what makes $\braket{v|v} = \sum_i |v_i|^2$ real and non-negative. Without it, the
``length'' of a complex vector could come out negative or complex.

Everything geometric comes out of this one operation:

\begin{center}
\begin{tabular}{lll}
\hline
concept & from the inner product & meaning \\
\hline
norm & $\|v\| = \sqrt{\braket{v|v}}$ & the length of the vector \\
orthogonality & $\braket{u|v} = 0$ & the vectors do not overlap \\
projection & $\dfrac{\braket{u|v}}{\braket{u|u}}\,u$ & how much of $v$ lies along $u$ \\
\hline
\end{tabular}
\end{center}

\section{The state vector}

A quantum state is a vector $\ket{\psi}$ in a Hilbert space with norm 1:
$\braket{\psi|\psi} = 1$. The system can be a particle's position, a photon's
polarization, or a qubit. For one qubit,
\[
  \ket{\psi} = \alpha\ket{0} + \beta\ket{1}, \qquad |\alpha|^2 + |\beta|^2 = 1 .
\]
The state does not store physical values directly. It stores how the system will
answer when you ask it something: each basis vector is a possible answer, and each
coefficient says how much the state is aligned with it.

Why a vector and not something else? Superposition is linear: if two states are
possible, any normalized combination of them is too, and that is exactly what a
vector space guarantees. Probabilities come from inner products,
$P = |\braket{\phi|\psi}|^2$, and the norm makes them add to 1. And operations on
states (gates, time evolution, measurements) are linear maps on the space.

The state vector is not a probability distribution: the amplitudes are complex,
and their phases, which do not show up in $|\alpha|^2$, decide interference. It is
not the physical system either, only its description. And not all of its phase is
physical.

\begin{concept}{Global and relative phase}
Multiplying the whole state by $e^{i\gamma}$ changes no probability, because
$|e^{i\gamma}\braket{\phi|\psi}|^2 = |\braket{\phi|\psi}|^2$. The states
$\ket{\psi}$ and $e^{i\gamma}\ket{\psi}$ are physically the same. But a phase
between components, as in $\ket{0} + e^{i\varphi}\ket{1}$, changes how they
interfere.

\textbf{What it fixes:} it says which part of the description is physics and which
is bookkeeping.

\textbf{What it costs:} the relative phase is invisible to a direct measurement
and decisive for the next operation. Quantum algorithms work by manipulating
exactly that.
\end{concept}

\chapter{State, Unitaries and Measurement}\label{ch:state}

\begin{etymology}
\textbf{State} comes from the Latin \emph{status}, from \emph{stare}, ``to
stand''. The condition something stands in at a given moment, whatever path
brought it there.

\textbf{Unitary} comes from \emph{unitas}, ``oneness'', from \emph{unus},
``one''. Something that stays whole. A unitary transformation can change the form
of a state but keeps it whole.

\textbf{Measurement} comes from \emph{mensura}, from \emph{metiri}, ``to measure'',
to judge by comparison with a standard.
\end{etymology}

Chapter~\ref{ch:qubit} used the words state, unitary and measurement loosely.
This chapter pins them down for a single qubit:

\begin{quote}
A unitary operation changes the state without destroying quantum information: it
moves amplitudes and phases around but keeps the length of the state vector. A
measurement does the opposite: it returns one result and replaces the
superposition with it.
\end{quote}

\section{The state as a vector}

Draw a vector of length 1 in the plane, at 45 degrees. Its shadows on the two
axes have length $\cos 45^\circ = \sin 45^\circ = 1/\sqrt{2} \approx 0.707$. Tilt it
to 30 degrees and the shadows become $0.866$ and $0.5$. Square them and you get
$0.75$ and $0.25$, which add to 1, because the vector still has length 1.

That is already the picture of a qubit. The vector is a combination of the two
axes, and tilting it moves weight from one to the other while keeping the total
fixed. The squared shadows are the probabilities of the two outcomes.

What the plane cannot show is phase. The state of a qubit is a vector in
$\mathbb{C}^2$: two complex numbers $(\alpha, \beta)$, each with a size and an
angle. Sizes become probabilities; the angle between them, the relative phase,
controls how the components combine. A real number describes a line. A pair of
reals describes a plane. A single complex number already describes a plane, and
two of them describe a qubit.

\section{The Bloch sphere}

Two complex numbers are four real numbers. Normalization,
$|\alpha|^2 + |\beta|^2 = 1$, removes one. The global phase, which no measurement
can see (Chapter~\ref{ch:hilbert}), removes another. Two remain, and two
numbers describe a point on a sphere. Every state of one qubit can be written as
\[
  \ket{\psi} = \cos\tfrac{\theta}{2}\,\ket{0} + e^{i\varphi}\sin\tfrac{\theta}{2}\,\ket{1},
\]
where $\theta$ is the latitude, measured from the north pole, and $\varphi$ the
longitude.

\begin{figure}[htbp]
\centering
\begin{tikzpicture}[scale=2.1]
  \draw[black!50, line width=0.7pt] (0,0) circle (1);
  \draw[black!35, dashed] (1,0) arc (0:180:1 and 0.3);
  \draw[black!50] (-1,0) arc (180:360:1 and 0.3);
  \draw[black!35, -{Stealth[length=1.6mm]}] (0,0) -- (0,1.25) node[above, font=\sffamily\small, text=black] {$\ket{0}$};
  \draw[black!35] (0,0) -- (0,-1) node[below, font=\sffamily\small, text=black] {$\ket{1}$};
  \draw[black!35, -{Stealth[length=1.6mm]}] (0,0) -- (-0.55,-0.42) node[below left, font=\sffamily\scriptsize, text=black!60] {$x$};
  \draw[black!35, -{Stealth[length=1.6mm]}] (0,0) -- (1.25,0) node[right, font=\sffamily\scriptsize, text=black!60] {$y$};
  % latitude of the state theta = 60 deg: z = cos 60 = 0.5, radius sin 60 = 0.866
  \draw[hfHistory, line width=0.7pt] (0.866,0.5) arc (0:-180:0.866 and 0.26);
  \draw[hfHistory, dashed, line width=0.5pt] (0.866,0.5) arc (0:180:0.866 and 0.26);
  \draw[hfConcept, line width=1.2pt, -{Stealth[length=2.2mm]}] (0,0) -- (-0.70,0.35);
  \fill[hfConcept] (-0.70,0.35) circle (0.9pt);
  \node[font=\sffamily\small, text=hfConcept] at (-0.88,0.22) {$\ket{\psi}$};
  \draw[black!60] (0,0.32) arc (90:150:0.32 and 0.3);
  \node[font=\sffamily\scriptsize] at (-0.1,0.42) {$\theta$};
  \node[hf/note, text=hfHistory, anchor=west] at (1.05,0.55) {same $|\alpha|^2$, any $\varphi$};
\end{tikzpicture}
\caption{The Bloch sphere. $\ket{0}$ is at the north pole and $\ket{1}$ at the
south pole. The latitude fixes the probabilities; moving around the circle of
latitude changes only the relative phase.}
\end{figure}

\begin{history}{1946 \textperiodcentered{} Stanford}
Felix Bloch introduced the sphere to describe nuclear spins in his work on
nuclear magnetic induction, the basis of NMR and, much later, of MRI. Henri
Poincaré had drawn the same sphere in 1892 for the polarization of light. The
mathematics is identical, because both are two-level systems.
\end{history}

Fix $|\alpha|^2 = 0.75$ and $|\beta|^2 = 0.25$ and you fix the latitude: the
state is at $\theta = 60^\circ$. But that does not fix the state. The relative
phase $\varphi$ moves it around that circle of latitude without touching the
probabilities. The state
$\sqrt{0.75}\,\ket{0} + 0.5\,\ket{1}$, with $\varphi = 0$, sits at the point
$(0.866, 0, 0.5)$. The sphere invents nothing; it only draws the two things a
qubit state carries, how much and which phase, in a form the eye can follow.

\section{Matrices}

If a state is a vector, an operation on a state is something that takes a vector
and returns another. In linear algebra that object is a \emph{matrix}. It can
rotate, stretch, flip or mix the components of a vector, depending on its
entries. A matrix is not just a table of numbers; it is the algebraic form of a
transformation (Chapter~\ref{ch:linear}).

\section{Unitary transformations}

The state has length 1, and that is not a technicality: it is what makes the
probabilities add to 100\%. So a physical operation cannot change a state any way
it likes. If it stretched or shrank the vector, the probabilities would stop
adding up. The operations allowed are the ones that preserve length.

In the real plane, the transformations that change direction but not length are
rotations. In complex space, they are called \emph{unitary}, and they satisfy
\[
  U^\dagger U = I ,
\]
where $U^\dagger$ is the conjugate transpose of $U$ and $I$ is the identity. A
state evolves as $\ket{\psi'} = U\ket{\psi}$, and the condition guarantees that
$\braket{\psi'|\psi'} = \braket{\psi|U^\dagger U|\psi} = \braket{\psi|\psi} = 1$.
It also guarantees that $U$ has an inverse, $U^\dagger$, so every unitary step can
be undone.

\begin{concept}{Quantum gate}
A unitary transformation used as an elementary step of a circuit, usually acting
on one or two qubits. Deutsch introduced quantum circuits made of gates in 1989.

\textbf{What it fixes:} complicated evolutions can be built from a small menu of
simple, calibrated operations. A few gates (for example $H$, $T$ and CNOT) are
enough to approximate any unitary.

\textbf{What it costs:} every gate on real hardware has an error, and a useful
algorithm needs millions or billions of them.
\end{concept}

The gates of Chapter~\ref{ch:qubit}, in matrix form:
\[
  H = \frac{1}{\sqrt 2}\begin{pmatrix} 1 & 1 \\ 1 & -1 \end{pmatrix},
  \qquad
  Z = \begin{pmatrix} 1 & 0 \\ 0 & -1 \end{pmatrix},
  \qquad
  X = \begin{pmatrix} 0 & 1 \\ 1 & 0 \end{pmatrix}.
\]
$X$ swaps $\ket{0}$ and $\ket{1}$, the quantum NOT. $Z$ flips the sign of
$\ket{1}$. On the Bloch sphere, each one is a rotation. And $HZH = X$: the
interference trick of Chapter~\ref{ch:qubit}, sign flip between two
Hadamards, is a NOT gate in disguise.

\section{Measurement}

Unitary evolution is smooth, reversible and deterministic. Measurement is none of
those things. Measuring a qubit in the basis $\{\ket{0}, \ket{1}\}$ returns 0
with probability $|\alpha|^2$ and 1 with probability $|\beta|^2$ (Born's rule,
Chapter~\ref{ch:quantum}). Afterwards, the state is no longer
$\alpha\ket{0} + \beta\ket{1}$: it is $\ket{0}$ or $\ket{1}$, whichever came out.
Von Neumann wrote this \emph{projection postulate} into his 1932 book. Measuring
again gives the same answer; the amplitudes are gone.

So you never see a superposition. You see one outcome per run, and to estimate the
amplitudes you have to prepare the same state and measure it many times. The goal
of an algorithm is not to look at the superposition but to arrange the unitaries
so that, at the moment of measurement, the useful answer has a high probability.

Measuring in a different basis is the same as rotating first and then measuring in
the standard one. Apply $H$ and measure, and you are asking whether the qubit was
$\ket{+} = (\ket{0} + \ket{1})/\sqrt2$ or $\ket{-} = (\ket{0} - \ket{1})/\sqrt2$,
a question about the relative phase that the standard measurement cannot answer.
Every measurement answers some questions and destroys the answers to others.

\section{Why complex amplitudes?}

The interference of Chapter~\ref{ch:qubit} only used signs, $+$ and $-$. Real
amplitudes with signs already reinforce and cancel. So why complex?

Signs give two phases, $0$ and $180$ degrees. Complex numbers give a continuous
circle of them, and quantum mechanics uses all of it: the time evolution
$e^{-i\hat Ht}$ turns phases continuously, and the longitude $\varphi$ of the Bloch
sphere would collapse to a single great circle without it.

For a long time physicists suspected that the complex numbers might be a
convenience, and that an equivalent theory with real numbers only could exist.
For single systems it can. In 2021, Marc-Olivier Renou and colleagues showed that
in networks with independent sources, real-number quantum theory makes different
predictions from the standard one, and experiments in 2022 sided with complex
numbers. The $i$ in quantum mechanics is not bookkeeping.

\chapter{Matrices, Eigenvalues and the Tensor Product}\label{ch:linear}

\begin{etymology}
\textbf{Transformation} comes from \emph{transformare}: \emph{trans-},
``across'', and \emph{formare}, ``to shape''. To carry something from one form
to another according to rules.

\textbf{Tensor} comes from the Latin \emph{tendere}, ``to stretch''. William
Rowan Hamilton used the word in 1846 for something else (the length of a
quaternion). Woldemar Voigt gave it its modern sense in 1898, for the quantities
that describe how a crystal stretches and deforms differently in different
directions.
\end{etymology}

Chapter~\ref{ch:hilbert} gave the space where states live. This one is about what
acts on them, and about how the spaces of several qubits combine. Three ideas carry
almost all of quantum computing: linear transformations, which are the gates;
eigenvalues, which are what measurements and many algorithms extract; and the
tensor product, which is where both the power of many qubits and entanglement
come from.

\section{Linear transformations and matrices}

A transformation $T$ between vector spaces is \emph{linear} if it respects their
structure:
\[
  T(\alpha u + \beta v) = \alpha\,T(u) + \beta\,T(v).
\]
Transforming a combination gives the same combination of the transformed
vectors. Most functions are not linear, but the ones that are have a remarkable
property: they are fully determined by what they do to the vectors of a basis.
Write $T(e_1), \dots, T(e_n)$ as columns, and you have a matrix.

\begin{history}{1850--1858 \textperiodcentered{} London}
James Joseph Sylvester coined the word \emph{matrix}, Latin for womb, in 1850,
for the rectangular array from which determinants are born. His friend Arthur
Cayley, a lawyer who did mathematics on the side, published ``A Memoir on the
Theory of Matrices'' in 1858. It treated matrices as objects in their own right,
with addition and multiplication, and defined multiplication so that it composes
transformations.
\end{history}

That is why matrix multiplication has its odd row-times-column formula: it was
designed so that the matrix of ``apply $S$, then $T$'' is the product $TS$. A
quantum circuit is a sequence of gates, so the matrix of a whole circuit is the
product of the matrices of its gates.

The determinant of a matrix measures how it scales areas (or volumes), with a
sign for orientation. A determinant of zero means the transformation flattens the
space and loses a dimension, so it cannot be undone. Unitary matrices always have
determinants of absolute value 1.

\section{Eigenvalues and eigenvectors}

Most transformations mix directions: a vector along $e_1$ may be sent to some
combination of $e_1$ and $e_2$. But some directions are special. They are only
stretched or shrunk, never turned. A vector $v \neq 0$ with
\[
  T v = \lambda v
\]
is an \emph{eigenvector} of $T$, and $\lambda$ is its \emph{eigenvalue}: the
direction stays, the scale changes by $\lambda$.

\begin{history}{1829--1904 \textperiodcentered{} Paris and Göttingen}
Cauchy proved in 1829 that symmetric matrices have real eigenvalues, while
studying quadratic forms. The words came later. In 1904 Hilbert, working on
integral equations, used \emph{Eigenwert}, ``own value'', and called the set of
those values the \emph{spectrum} of the operator. He later said he had chosen
the name with no idea that it would turn out to describe the spectra of atoms.
Twenty years later it did: the energy levels of an atom are the eigenvalues of its
Hamiltonian.
\end{history}

Operators for physical quantities are \emph{Hermitian} ($A = A^\dagger$), and
their spectra are always real, as a measurement result must be. The
\emph{spectral decomposition} says such an operator is entirely given by its
eigenvectors and eigenvalues:
\[
  A = \sum_k \lambda_k \ket{v_k}\bra{v_k} .
\]
In the basis of its eigenvectors, $A$ is diagonal, and diagonal is simple.

Take $A = \begin{pmatrix} 3 & 1 \\ 1 & 2 \end{pmatrix}$. Its eigenvalues are
$(5 \mp \sqrt5)/2$, about $1.382$ and $3.618$. Along the two eigenvectors, $A$
only stretches, by those factors. Every other vector gets turned. Rebuilding
$\lambda_1\ket{v_1}\bra{v_1} + \lambda_2\ket{v_2}\bra{v_2}$ gives back $A$ exactly.

\begin{concept}{Spectrum}
The set of eigenvalues of an operator.

\textbf{What it fixes:} it reduces an operator to the numbers that matter. For a
Hamiltonian, they are the energies; for a measurement, the possible results.

\textbf{What it costs:} for a matrix of size $2^n$, computing the spectrum
classically takes time that grows with $2^n$. Extracting it on a quantum computer
is what phase estimation does, and a large part of useful quantum computing is,
underneath, spectral processing.
\end{concept}

\section{Tensors}

A vector has one index, $v_i$; a matrix has two, $M_{ij}$. A \emph{tensor} can
have any number. But a tensor is not just a table with many indices. It is a table
that changes in a specific way when you change the basis, so that equations
written with tensors look the same in every coordinate system.

\begin{history}{1887--1915 \textperiodcentered{} Padua and Zurich}
Gregorio Ricci-Curbastro developed an ``absolute differential calculus'' in the
1880s and 1890s, a way of writing geometry that does not depend on coordinates.
With his student Tullio Levi-Civita he published the full method in 1900. Almost
nobody used it until Marcel Grossmann showed it to his friend Albert Einstein in
1912. Einstein needed exactly that to describe curved spacetime, and general
relativity, in 1915, is written in tensors.
\end{history}

Quantum computing uses less of this machinery than the name suggests. What it
uses all the time is one operation that builds new spaces out of old ones.

\section{The tensor product}

Given two spaces $V$ and $W$, of dimensions $m$ and $n$, the \emph{tensor product}
$V \otimes W$ is a space of dimension $mn$. If $\{e_i\}$ is a basis of $V$ and
$\{f_j\}$ a basis of $W$, the $mn$ products $e_i \otimes f_j$ are a basis of
$V \otimes W$.

One qubit lives in $\mathbb{C}^2$. Two qubits live in
$\mathbb{C}^2 \otimes \mathbb{C}^2 = \mathbb{C}^4$, with basis
$\ket{00}, \ket{01}, \ket{10}, \ket{11}$. And $n$ qubits live in
$\mathbb{C}^{2^n}$. Dimensions multiply: each new qubit doubles the size of the
space.

That doubling is the source of both the promise and the problem. Describing 50
qubits on a classical computer takes $2^{50} \approx 1.1\times10^{15}$ complex
amplitudes. At 16 bytes each, that is about 18 petabytes of memory, for one state.
This is Feynman's argument from Chapter~\ref{ch:history}, in numbers.

\section{Entanglement}

Some states of two qubits can be written as a product of one state for each:
$\ket{0}\otimes\ket{1} = \ket{01}$, or
$\ket{+}\otimes\ket{+} = \tfrac12(\ket{00} + \ket{01} + \ket{10} + \ket{11})$.
These are \emph{separable}. But most vectors of $V \otimes W$ cannot be split this
way. The simplest example is the Bell state
\[
  \ket{\Phi^+} = \frac{\ket{00} + \ket{11}}{\sqrt 2} .
\]
Try to write it as $(a\ket0 + b\ket1)\otimes(c\ket0 + d\ket1)$: you need $ac$ and
$bd$ to be nonzero and $ad$ and $bc$ to be zero, which is impossible. The two qubits
have no state of their own. Only the pair does.

\begin{history}{1935 \textperiodcentered{} Princeton and Oxford}
Einstein, Boris Podolsky and Nathan Rosen published a paper arguing that states
like this show quantum mechanics is incomplete. Schrödinger answered the same year
and gave the phenomenon its name, \emph{Verschränkung}, which he translated as
\emph{entanglement}. He called it ``not one but rather the characteristic trait of
quantum mechanics''.
\end{history}

\begin{concept}{Entanglement}
A state of several systems that cannot be written as a product of states of each
system. Mathematically, non-separability in the tensor product.

\textbf{What it fixes:} correlations between qubits live in the state itself,
with no extra bits to store them. Every algorithm with an exponential speedup
produces large amounts of it.

\textbf{What it costs:} entangled states are fragile. Disturbing one qubit
disturbs the whole state, and the environment is always disturbing something.
\end{concept}

\section{Measuring entanglement: the Schmidt decomposition}

Any state of two systems can be written as
\[
  \ket{\psi} = \sum_{k=1}^{r} \lambda_k\,\ket{a_k}\otimes\ket{b_k},
\]
with orthonormal $\ket{a_k}$, orthonormal $\ket{b_k}$ and positive coefficients
$\lambda_k$. The number of terms $r$ is the \emph{Schmidt rank}. If $r = 1$, the
state is separable. If $r > 1$, it is entangled, and the more evenly spread the
$\lambda_k$, the more entangled it is.

Finding the decomposition is easy: arrange the four amplitudes of a two-qubit
state as a $2\times2$ matrix and compute its singular values. That is the
\emph{singular value decomposition}, found by Eugenio Beltrami in 1873 and Camille
Jordan in 1874. Erhard Schmidt rediscovered it in 1907 for integral equations,
which is why physicists use his name. For $\ket{01}$ and for
$\ket{+}\otimes\ket{+}$, the singular values are $(1, 0)$: rank 1, separable. For
the Bell state they are $(0.707, 0.707)$: rank 2, as entangled as two qubits can be.

Entanglement, the resource with no classical analogue, turns out to be a question
about the rank of a matrix, a tool from nineteenth-century linear algebra.

\section{Identical particles: bosons and fermions}\label{sec:identical}

The tensor product keeps its factors apart: there is a first qubit and a second
qubit, and $\ket{01}$ is not $\ket{10}$. For two qubits on a chip that is right,
because they sit in different places. For two electrons it is wrong. Every
electron in the universe is identical to every other one. Nothing carries a label
saying ``electron 1'', so swapping two electrons cannot change anything that can be
measured.

Write the swap as an operator, $P(\ket a \otimes \ket b) = \ket b \otimes \ket a$.
For identical particles, $P\ket\psi$ must be the same physical state as $\ket\psi$,
and Chapter~\ref{ch:hilbert} showed that two vectors describe the same state only if
they differ by a global phase: $P\ket\psi = e^{i\theta}\ket\psi$. Swapping twice
puts everything back, $P^2 = I$, so $e^{2i\theta} = 1$, and the phase can only be
$+1$ or $-1$. Nature uses both.
\[
  \text{bosons: } \frac{\ket{ab} + \ket{ba}}{\sqrt2},
  \qquad
  \text{fermions: } \frac{\ket{ab} - \ket{ba}}{\sqrt2}.
\]
Now put both particles in the same state, $a = b$. The boson state becomes
$\sqrt2\,\ket{aa}$, a perfectly good state once normalized. The fermion state
becomes $\ket{aa} - \ket{aa} = 0$, which is not a state at all. Two fermions can
never occupy the same state. That is the \emph{exclusion principle}, and here it
fell out of one minus sign.

\begin{history}{1924--1926 \textperiodcentered{} Dacca, Hamburg, Florence, Cambridge}
In June 1924 Satyendra Nath Bose, a lecturer at the University of Dacca, sent
Einstein a short paper that derived Planck's law by counting photons as truly
indistinguishable. Einstein translated it into German, had it published, and
extended it to atoms. In January 1925 Wolfgang Pauli, in Hamburg, explained the
shells of electrons in atoms with a rule: no two electrons in the same state. In
1926 Enrico Fermi, in Florence, and Paul Dirac, in Cambridge, worked out
independently how particles obeying that rule fill up energy levels. Particles of
the first kind became \emph{bosons}, after Bose, and of the second \emph{fermions},
after Fermi. In 1940 Pauli proved which is which: particles with integer spin are
bosons, particles with half-integer spin are fermions.
\end{history}

The sign matters well beyond atoms, and this book runs into it three times.

\textbf{Matter.} Electrons are fermions, so they cannot all drop into the lowest
level of an atom; they stack up in shells, and chemistry is the result. Photons are
bosons, so they can pile into one state by the billion, and a laser is the result.

\textbf{Superconductors.} Two electrons bound together, a \emph{Cooper pair}, behave
as a boson: swapping two pairs swaps two fermions twice, and $(-1)^2 = +1$. Bosons
can share one state, so below a critical temperature all the pairs of a metal fall
into a single collective state. That is superconductivity, and it is what makes the
qubits of Chapter~\ref{ch:hardware} work.

\textbf{Computation.} The minus sign is where the cost of simulating electrons
hides. A classical simulation of many electrons has to keep track of states that
change sign every time two of them are swapped, and in a random sampling the
positive and negative contributions cancel, leaving noise. That is the most famous
form of the sign problem of Chapter~\ref{ch:history}, and the reason simulating
molecules is one of the problems in Part~\ref{part:problems}.

\subsection*{In two dimensions, any phase}

The argument that left only $\pm1$ hid an assumption: that swapping twice is the
same as doing nothing. Swapping twice means carrying one particle all the way around
the other and back to where it started. In three dimensions, that loop can be lifted
over the other particle and shrunk to a point, so it really is the same as staying
put, and $P^2 = I$. In two dimensions, a loop around a particle cannot leave the
plane. To shrink it, it would have to pass through the particle. Nothing then forces
$P^2 = I$, and the phase of an exchange can be anything.

\begin{figure}[htbp]
\centering
\begin{tikzpicture}[x=1mm, y=1mm]
  % ---------- (a) three dimensions ----------
  \begin{scope}
    \node[font=\sffamily\small] at (20,33) {(a) three dimensions};
    \fill[hfSource] (20,8) circle (1.6);
    \node[font=\small, left=1.2mm] at (20,8) {$B$};
    \draw[hfConcept, line width=1pt] (20,8) ellipse (14 and 4.5);
    \fill[hfConcept] (34,8) circle (1.4);
    \node[font=\small, right=1mm] at (34,8) {$A$};
    \draw[hfConcept, -{Stealth[length=2mm]}, line width=0.8pt] (21,12.5) -- (19.5,12.5);
    % lifted loop
    \draw[hfConcept, dashed, line width=0.8pt] (20,20) ellipse (8 and 2.5);
    \draw[hfConcept, dashed, line width=0.8pt] (20,27) ellipse (2.5 and 0.8);
    \draw[black!40, -{Stealth[length=2mm]}, line width=0.8pt] (37,14) to[bend right=20] (29,24);
    \node[hf/note, anchor=west] at (37,22) {lift the loop\\over $B$: it\\shrinks to\\nothing};
    \node[hf/note] at (20,-5) {going around twice is the same as\\staying put, so $P^2 = I$ and the\\phase is $+1$ or $-1$};
  \end{scope}
  % ---------- (b) two dimensions ----------
  \begin{scope}[xshift=78mm]
    \node[font=\sffamily\small] at (20,33) {(b) two dimensions};
    \fill[black!5] (-6,1) -- (40,1) -- (48,17) -- (2,17) -- cycle;
    \draw[black!30, line width=0.6pt] (-6,1) -- (40,1) -- (48,17) -- (2,17) -- cycle;
    \fill[hfSource] (21,9) circle (1.6);
    \node[font=\small, left=1.2mm] at (21,9) {$B$};
    \draw[hfConcept, line width=1pt] (21,9) ellipse (13 and 4.5);
    \fill[hfConcept] (34,9) circle (1.4);
    \node[font=\small, right=1mm] at (34,9) {$A$};
    \draw[hfConcept, -{Stealth[length=2mm]}, line width=0.8pt] (22,13.5) -- (20.5,13.5);
    \node[hf/note, text=hfAlert] at (21,24) {the loop cannot leave the plane:\\to shrink it, $A$ must cross $B$};
    \node[hf/note] at (21,-5) {going around twice is not the same\\as staying put, and the phase of\\an exchange can be anything};
  \end{scope}
\end{tikzpicture}
\caption{Why only two kinds of particle exist in three dimensions and more exist in
two. Exchanging $A$ and $B$ twice is the same as carrying $A$ once around $B$. (a) In
space, that loop can be lifted over $B$ and shrunk to a point, so a double exchange
does nothing. (b) On a plane, the loop is trapped around $B$, and the double exchange
can leave a phase behind. Particles that use that freedom are anyons.}\label{fig:exchange}
\end{figure}

\begin{history}{1977--1982 \textperiodcentered{} Oslo and Santa Barbara}
Jon Magne Leinaas and Jan Myrheim, in Oslo, found the loophole in 1977 by looking at
the space of all possible paths of identical particles: in one and two dimensions, an
exchange can give any phase. In 1982 Frank Wilczek, at Santa Barbara, built a concrete
example, a charged particle bound to a thin tube of magnetic flux, and named such
particles \emph{anyons}, because the phase can be anything.
\end{history}

A two-dimensional world sounds like a mathematical game, but electrons confined in a
thin layer of semiconductor, in a strong magnetic field, behave as one. Their collective
excitations in the fractional quantum Hall effect, discovered in 1982, are anyons, and in
2020 two groups, in Paris and at Purdue, measured their exchange phase directly.

There is one more step, and it is the one topological quantum computing rests on. If
several anyons share a set of states with the same energy, an exchange does not have to
multiply the state by a number. It can apply a matrix, a unitary on that set of states,
and then the order of the exchanges matters: exchanging $A$ with $B$ and then $B$ with
$C$ is not the same as the other way around. Such particles are \emph{non-abelian}
anyons, after the mathematician Niels Henrik Abel, whose name is attached to operations
whose order does not matter. Braiding non-abelian anyons around each other is a quantum
gate, and Section~\ref{sec:topological} is about the hardware that tries to build them.


% Then one part per layer of the stack, from the problems down to the market.
\part{The Problems}\label{part:problems}
\chapter{Undoing a Multiplication}\label{ch:factoring}

\begin{etymology}
\textbf{Factor} comes from the Latin \emph{factor}, ``the one who makes'', from
\emph{facere}, ``to make''. To factor a number is to find the things that make
it: to write it back as the product that produced it.

\textbf{Modulus} is Latin for ``a small measure''. Gauss brought it into
arithmetic in 1801, in the \emph{Disquisitiones Arithmeticae}.
\end{etymology}

There are several problems a quantum computer might solve better than a classical
one, and they are not equally useful. Feynman's original one, simulating molecules
and materials, is the one most likely to pay off. Optimization problems, like
choosing the best investment portfolio, are the ones most often promised and the
least supported by evidence. Factoring is neither. It is the problem that made the
field famous, because it threatens the cryptography the internet runs on, and it is
the one whose cost is best understood. That makes it the best place to start
measuring how far away a useful machine is.

The number 42 can be written as $2 \times 3 \times 7$. Finding those three
numbers is what \emph{factoring} means: given a number, find the numbers that
multiply to give it.

For a computer, the two directions are very different. Multiplying two primes
is trivial: $37 \times 53 = 1961$, and a machine does that with numbers
thousands of digits long in a fraction of a second. Going back, from $1961$ to
$37$ and $53$, is hard. For small numbers you can try divisors one by one. For
numbers with hundreds of digits, nobody knows a fast way.

This chapter is about why that matters, and about the move that made factoring
the most famous problem in quantum computing: Peter Shor showed that factoring
is the same as finding a period.

\section{A lock that anyone can close}

\begin{history}{1976--1977 \textperiodcentered{} Stanford and MIT}
In 1976, Whitfield Diffie and Martin Hellman published ``New Directions in
Cryptography'', proposing that two people could agree on a secret over a public
channel. Their scheme rested on the discrete logarithm: computing $g^a \bmod p$
is easy, recovering $a$ from the result is not. A year later, at MIT, Ron
Rivest, Adi Shamir and Leonard Adleman built a system on a different hard
problem, factoring. RSA was published in 1978. In 1997 the British government
revealed that Clifford Cocks, at GCHQ, had found the same idea in 1973 and kept
it secret.
\end{history}

\begin{concept}{Public-key cryptography}
A pair of keys: a public one that anyone can use to lock a message, and a
private one that unlocks it. In RSA, the public key contains a number
$N = p \times q$, and the private key depends on knowing $p$ and $q$.

\textbf{What it fixes:} two strangers can talk securely without ever meeting to
exchange a secret.

\textbf{What it costs:} all of the security rests on a hard problem that nobody
has proved is hard. If someone finds a fast way to factor, the lock opens.
\end{concept}

\section{How hard is hard?}

The best classical algorithm known for large numbers is the \emph{general number
field sieve}. Its running time grows like
\[
  \exp\!\Big(c\,(\ln N)^{1/3}(\ln\ln N)^{2/3}\Big),
  \qquad c = (64/9)^{1/3} \approx 1.923 .
\]
Read the formula piece by piece. $N$ is the number to factor, and $\ln N$ is
proportional to its number of digits, the real measure of the size of the
input. The main part of the exponent is $(\ln N)^{1/3}$; the
$(\ln\ln N)^{2/3}$ factor grows very slowly. The result is called
\emph{subexponential}: faster than any polynomial in the number of digits,
slower than a pure exponential.

\begin{history}{1988--1993 \textperiodcentered{} Pollard and the Lenstras}
John Pollard proposed the idea of the number field sieve in 1988, in a letter
to Andrew Odlyzko. Arjen Lenstra, Hendrik Lenstra, Mark Manasse and Pollard
turned it into a working algorithm, and in 1990 they used it to factor the ninth
Fermat number, $2^{512}+1$, which had resisted everything else. The general
version, for any number, was published in 1993. Thirty years later it is still
the best classical method for large numbers.
\end{history}

Plugging numbers in shows the shape of the growth. The table drops the
lower-order terms hidden in the exact formula, so only the comparison between
rows means anything.

\begin{center}
\begin{tabular}{lrr}
\hline
number & digits & $\exp(\dots)$ \\
\hline
$42$ & 2 & $\approx e^{3.6}$ \\
$1961$ & 4 & $\approx e^{6.0}$ \\
RSA-250 & 250 & $\approx 10^{24}$ \\
RSA-2048 & 617 & $\approx 10^{35}$ \\
\hline
\end{tabular}
\end{center}

RSA-250 is not a hypothetical row. In February 2020, a team led by Fabrice
Boudot and Paul Zimmermann factored it with the number field sieve, using about
2700 core-years of computation. It is the largest RSA challenge number
factored so far. RSA-2048, the size used in practice, sits roughly eleven orders
of magnitude further. Nobody is going to get there by waiting for faster
classical machines.

\section{Shor's answer}

In 1994 Peter Shor, at Bell Labs, showed that a quantum computer can factor in
polynomial time. His bound for the circuit is
\[
  O\!\big((\ln N)^2 \cdot \ln\ln N \cdot \ln\ln\ln N\big).
\]
The $(\ln N)^2$ term dominates. It comes from modular exponentiation, which
does $O(\ln N)$ multiplications of numbers with $O(\ln N)$ digits. The
$\ln\ln N$ and $\ln\ln\ln N$ terms come from the fast multiplication algorithm
of Schönhage and Strassen used inside each step, and in practice they are small
constants.

There is no exponential anywhere. For RSA-2048, $(\ln N)^2$ is about
$2\times10^{6}$ and the whole expression lands around $10^{7}$. That number counts
elementary operations of an ideal machine. On real hardware every one of those
operations has to be protected against errors, and that is where the real
distance lies (Parts~\ref{part:hardware} and~\ref{part:qec}). In 2019, Craig Gidney and Martin Ekerå
estimated that factoring RSA-2048 would take 20 million noisy physical qubits
running for eight hours. In 2025, Gidney brought that down to under one million
qubits, running for under a week. The largest gate-based machines of 2025 had
on the order of a thousand physical qubits.

So Shor did not make RSA breakable tomorrow. What he showed is that the
difficulty of factoring is a fact about classical computers, not about the
problem. The rest of this chapter is about how he found that out.

\section{Factoring is finding a period}

Shor did not attack factoring head-on. He used a reduction from number theory,
known since Gary Miller's work on primality in 1976: if you can find the
\emph{period} of a certain function, you can factor.

Take the number $N$ to factor and pick some $x$ that shares no divisor with $N$
other than 1 (that is, $\gcd(x, N) = 1$). Now compute the powers of $x$, but
keep only the remainder of the division by $N$ at each step. That is what
``mod $N$'' means: $a \bmod N$ is the remainder of $a$ divided by $N$, and
$a \equiv b \pmod N$ means $a$ and $b$ leave the same remainder. The function is
\[
  f(a) = x^a \bmod N .
\]
Take $N = 15$ and $x = 2$:

\begin{center}
\begin{tabular}{r|rrrrrrrrrrrr}
$a$ & 0 & 1 & 2 & 3 & 4 & 5 & 6 & 7 & 8 & 9 & 10 & 11 \\
\hline
$2^a$ & 1 & 2 & 4 & 8 & 16 & 32 & 64 & 128 & 256 & 512 & 1024 & 2048 \\
$2^a \bmod 15$ & 1 & 2 & 4 & 8 & 1 & 2 & 4 & 8 & 1 & 2 & 4 & 8 \\
\end{tabular}
\end{center}

The powers grow without end, but the remainders repeat: $1, 2, 4, 8$, and again.

\begin{figure}[htbp]
\centering
\begin{tikzpicture}[x=8mm, y=3.2mm]
  \draw[black!40] (-0.3,0) -- (11.6,0);
  \draw[black!40] (-0.3,0) -- (-0.3,9.5);
  \foreach \a/\v in {0/1,1/2,2/4,3/8,4/1,5/2,6/4,7/8,8/1,9/2,10/4,11/8} {
    \draw[hfConcept, line width=1.2pt] (\a,0) -- (\a,\v);
    \fill[hfConcept] (\a,\v) circle (2.2pt);
    \node[font=\sffamily\scriptsize, below] at (\a,0) {\a};
  }
  \foreach \v in {1,2,4,8} \node[font=\sffamily\scriptsize, left] at (-0.3,\v) {\v};
  \foreach \s in {0,4,8} {
    \draw[hfHistory, decorate, decoration={brace, amplitude=4pt}] (\s-0.1,9.8) -- (\s+3.1,9.8)
      node[midway, above=3pt, font=\sffamily\scriptsize, text=hfHistory] {$r = 4$};
  }
  \node[font=\sffamily\small, below=5mm] at (5.5,0) {$a$};
\end{tikzpicture}
\caption{The function $2^a \bmod 15$. The remainders repeat every four steps:
the period is $r = 4$.}
\end{figure}

The table does not show any factor of 15. It shows something else: the
\emph{period} of the function, the smallest $r$ with $x^r \equiv 1 \pmod N$.
Here $r = 4$. Now comes the algebra that turns the period into factors. The
goal is modest: build a product that $N$ divides.

If $x^r \equiv 1 \pmod N$, then $x^r - 1 \equiv 0 \pmod N$, which means $N$
divides $x^r - 1$. So $x^r - 1$ is a bigger number that contains the factors of
$N$. When $r$ is even, it splits as a difference of squares:
\[
  x^r - 1 = \big(x^{r/2} - 1\big)\big(x^{r/2} + 1\big).
\]
$N$ divides the product, so with some luck each factor of $N$ ends up in a
different parenthesis. The greatest common divisor, which Euclid's algorithm
computes fast, pulls them out. With $x = 2$ and $r = 4$, $x^{r/2} = 4$:
\[
  \gcd(4 - 1, 15) = 3, \qquad \gcd(4 + 1, 15) = 5 .
\]
And $15 = 3 \times 5$.

It does not always go this well. If $r$ comes out odd, the difference of squares
is not available. If $x^{r/2} \equiv -1 \pmod N$, all of $N$ falls into one
parenthesis and the gcd returns $N$ itself, which is useless. Take $x = 14$:
$14^2 = 196 \equiv 1 \pmod{15}$, so $r = 2$, but $14 \equiv -1$, and nothing
comes out. In both cases you pick another $x$ and try again. For a number with
at least two distinct odd prime factors, a random $x$ works at least half the
time, so a few attempts are enough. And if the divisor you find is not prime,
you have still split the problem into smaller numbers, and you keep going.

\begin{concept}{Reduction to period finding}
Factoring $N$ becomes: pick $x$, find the period $r$ of $x^a \bmod N$, compute
$\gcd(x^{r/2} \pm 1, N)$.

\textbf{What it fixes:} everything in this recipe is classically easy except one
step, finding $r$. The hard part of factoring is isolated in a single, precise
question.

\textbf{What it costs:} classically, finding $r$ is as hard as factoring, because
the period can be as large as $N$ itself. The reduction only pays off if
something else can find periods fast.
\end{concept}

Periods are what Fourier analysis was made for, and a quantum computer can
perform a Fourier transform on an exponentially large table of values in one go.
That is Shor's algorithm in one sentence. How the table of $x^a \bmod N$ is
written into qubits, and how the transform reads its period, is the subject of
Part~\ref{part:algorithm}.


\part{The Algorithm}\label{part:algorithm}
\chapter{Encoding a Problem in Qubits}\label{ch:encoding}

\begin{etymology}
\textbf{Encode} comes from \emph{code}, from the Latin \emph{codex}, ``tree
trunk'', then the wooden tablets used for writing, then a book, and finally a
book of laws. To encode is to choose a set of symbols and rules to write something
down.

\textbf{Register} comes from \emph{regesta}, ``things written down'', from
\emph{regerere}, ``to record''. The place where information is kept while you
work with it.

\textbf{Basis} is Greek for ``step'' or ``what you stand on''. The set of
references you describe everything else with.
\end{etymology}

\begin{quote}
To encode a problem in qubits is to decide what the basis states of a register
stand for.
\end{quote}

In practice that means choosing a correspondence between configurations of the
problem and strings of bits. A register of $n$ qubits has $2^n$ basis states, so
encoding is deciding how the variables, candidates or answers of the problem are
translated into those states. This chapter shows how that works in general, how
Shor did it for factoring, and what each choice costs.

\section{The computational basis is an alphabet}

The basis states of a register are all the strings of $n$ bits. One qubit has
$\ket0$ and $\ket1$. Two qubits have $\ket{00}, \ket{01}, \ket{10}, \ket{11}$. Three
have eight, from $\ket{000}$ to $\ket{111}$.

Imagine the eight states of three qubits are eight people. The first bit says
whether a person wears a hat, the second whether they wear glasses, the third
whether their hand is raised. Then $\ket{010}$ and $\ket{111}$ are not mysterious;
they are labels for two profiles. Each basis state is one configuration of the
problem.

A classical computer that wants to find the people matching some rule checks them
one by one. A quantum computer can put the register in a superposition of all
eight and apply the rule to all of them in a single operation: ``everyone who
matches, flip your sign''. The operation acts on every component of the
superposition at once.

That does not mean you can read the eight answers. A measurement returns one basis
state. Getting the full record of every person is a question a quantum computer
does not answer any better than a classical one. What it can do is interfere the
signs so that the measurement is likely to return a person who matches. With $n$
features there are $N = 2^n$ profiles. Checking them one by one costs $O(N)$; with
Grover's interference it costs $O(\sqrt N)$.

\section{Registers}

A register is a group of qubits that together represent a number in binary. With
three qubits, $\ket{101}$ means $1\cdot4 + 0\cdot2 + 1\cdot1 = 5$. So far this is
exactly what a classical register does. The difference is that a Hadamard on each
of the three qubits puts the register in an even superposition of all eight
values, $0$ to $7$, at once. Three physical qubits, a space of eight dimensions.
To hold the numbers from $0$ to $N - 1$ you need $\lceil \log_2 N \rceil$ qubits:
four for $N = 15$, since $2^3 = 8 < 15 \leq 16 = 2^4$.

\section{How Shor encodes factoring}

Chapter~\ref{ch:factoring} reduced factoring to finding the period $r$ of $f(a) = x^a \bmod N$. Shor
needed two things in the quantum state at the same time: the inputs $a$ being
tried, and the value of $f$ at each of them. So he used two registers.

\begin{figure}[htbp]
\centering
\begin{tikzpicture}[node distance=7mm]
  \node[hf/reader, minimum width=26mm, minimum height=10mm] (r1) {register 1\\{\footnotesize $t$ qubits, holds $a$}};
  \node[hf/source, minimum width=26mm, minimum height=10mm, below=6mm of r1] (r2) {register 2\\{\footnotesize holds $x^a \bmod N$}};
  \node[hf/bus, minimum width=12mm, minimum height=10mm, right=8mm of r1] (h) {$H^{\otimes t}$};
  \node[hf/bus, minimum width=22mm, minimum height=26mm, right=8mm of h, yshift=-8mm] (u) {modular\\exponentiation};
  \node[hf/bus, minimum width=14mm, minimum height=10mm, right=8mm of u, yshift=8mm] (qft) {QFT};
  \node[hf/reader, minimum width=16mm, minimum height=10mm, right=8mm of qft] (m) {measure};
  \draw[hf/ask] (r1) -- (h);
  \draw[hf/ask] (h) -- (h -| u.west);
  \draw[hf/ask] (h -| u.east) -- (qft);
  \draw[hf/ask] (qft) -- (m);
  \draw[hf/answer] (r2) -- (r2 -| u.west);
  \node[hf/note, below=1mm of m, text width=26mm] {an integer close to a multiple of $q/r$};
\end{tikzpicture}
\caption{The structure of Shor's algorithm. The first register is put in
superposition, the second receives $f(a)$ for every $a$ at once, and the Fourier
transform turns the period hidden in the first register into something a
measurement can see.}
\end{figure}

\textbf{Register 1} holds the exponent $a$. It has $t$ qubits, so it can represent
$q = 2^t$ values, from $0$ to $q - 1$. Shor takes $t \approx 2\log_2 N$, for a
reason that comes at the end. Hadamards on every qubit put it in an even
superposition, while register 2 starts at zero:
\[
  \frac{1}{\sqrt q}\sum_{a=0}^{q-1} \ket{a}\ket{0}.
\]
The factor $1/\sqrt q$ keeps the norm at 1.

\textbf{Register 2} receives $f(a)$. A circuit computes $x^a \bmod N$ for the value
in register 1 and writes it into register 2. Since register 1 holds every value
at once, the result is a table of all the pairs:
\[
  \frac{1}{\sqrt q}\sum_{a=0}^{q-1} \ket{a}\,\ket{x^a \bmod N}.
\]
The period is now encoded in this state. Every $a$ that gives the same value of
$f$ is tied to it, and those $a$ are spaced exactly $r$ apart.

\section{How modular exponentiation is computed}

There is nothing mysterious in the circuit that computes $x^a \bmod N$. It is the
classical method, made reversible. Write $a$ in binary,
$a = a_0 + 2a_1 + 4a_2 + 8a_3 + \cdots$, where each $a_i$ is 0 or 1. Then
\[
  x^a = x^{a_0}\cdot x^{2a_1}\cdot x^{4a_2}\cdot x^{8a_3}\cdots,
\]
and each factor depends on a single bit: it is 1 if the bit is off and
$x^{2^i}$ if it is on. For $a = 13 = 1101_2 = 1 + 4 + 8$, this gives
$x^{13} = x\cdot x^4\cdot x^8$.

The powers $x, x^2, x^4, x^8, \dots \bmod N$ are computed in advance, classically,
by squaring repeatedly. The circuit then reads the bits of register 1 one at a
time, and each bit controls whether register 2 is multiplied by its precomputed
constant. With $t$ bits there are $t$ controlled multiplications, each on numbers
of about $\log_2 N$ bits. That is where the $(\log N)^2$ of Shor's cost comes from
(Chapter~\ref{ch:factoring}).

\begin{history}{1961--1973 \textperiodcentered{} IBM}
In 1961 Rolf Landauer showed that erasing a bit of information has an unavoidable
energy cost. Charles Bennett, also at IBM, showed in 1973 that any computation can
be made reversible, keeping enough information to run it backwards. Both results
were about the thermodynamics of ordinary computers. They became essential to
quantum computing, where every step must be unitary and therefore reversible: the
quantum circuit for $x^a \bmod N$ is Bennett's construction applied to the
schoolbook algorithm.
\end{history}

\section{Reading the period}

The last step is a quantum Fourier transform on register 1 (Chapter~\ref{ch:fourier}).
It turns a periodic pattern into peaks: the probability concentrates on values of
$k$ close to multiples of $q/r$. For $N = 15$, $x = 2$ and a register of four
qubits ($q = 16$), the period is $r = 4$, and the measurement returns $0$, $4$, $8$
or $12$, each with probability exactly $1/4$. Dividing by $q$ gives
$0, 1/4, 1/2, 3/4$: fractions with denominator $r = 4$.

When $r$ does not divide $q$, the peaks are not exact. The measured $k/q$ is then
only close to some $j/r$, and the classical method of continued fractions finds the
fraction with the smallest denominator nearby. Taking $q \geq N^2$, that is,
$t \approx 2\log_2 N$ qubits, makes the approximation good enough for continued
fractions to find $r$ with high probability. That is the reason for the size of
register 1.

\section{What the encoding decided}

Shor's choice of encoding decided everything that follows: which qubits carry which
part of the problem, how the information is processed, and how it is read out. It
also fixed the cost. Register 1 needs $O(\log N)$ qubits, modular exponentiation
needs $O((\log N)^3)$ gates in the schoolbook version, and all of it must run under
error correction, because noise anywhere in the exponentiation or the transform
corrupts the period. That cost belongs to this encoding, not to factoring itself.
Most of the progress in resource estimates since 1994 has come from finding cheaper
ways to do the same arithmetic.

\section{Three ways to encode}

Shor's is one style. There are three broad ones, each with its own costs.

\begin{concept}{Basis encoding}
Each configuration of the problem is a basis state: $n$ binary variables, $n$
qubits, and $\ket{x_1x_2\dots x_n}$ is one assignment of values. Used by Grover's
algorithm, by Shor's, and by QAOA (Farhi, Goldstone and Gutmann, 2014) for
combinatorial optimization.

\textbf{What it fixes:} reading the answer is direct. You measure and get a string
of bits.

\textbf{What it costs:} the number of qubits grows linearly with the number of
variables, and you still need interference to make the right string likely.
\end{concept}

\begin{concept}{Amplitude encoding}
The data go into the values of the amplitudes, not into which basis state is which.
A vector of $2^n$ numbers fits in $n$ qubits. Used by the HHL algorithm for linear
systems (Harrow, Hassidim and Lloyd, 2009) and by many proposals for quantum
machine learning.

\textbf{What it fixes:} an exponential compression: $2^n$ numbers in $n$ qubits.

\textbf{What it costs:} loading $2^n$ numbers into amplitudes generally takes
$2^n$ steps, and reading them back out takes many measurements. Either cost can
erase the advantage.
\end{concept}

Amplitude encoding is where most claims of exponential speedup for machine learning
die. Scott Aaronson listed the fine print of HHL in 2015: the input must load fast,
the matrix must be well behaved, and the output is a quantum state you can only
sample. In 2018 Ewin Tang, an undergraduate at the University of Texas, showed that
the most celebrated quantum recommendation algorithm could be matched by a classical
one under the same assumptions about the input. When a quantum machine learning
paper does not say how the data get into the amplitudes, assume they do not.

\begin{concept}{Hamiltonian encoding}
The problem goes into an operator, not into the initial state: build a Hamiltonian
$\hat H$ whose lowest-energy state is the solution. Used in adiabatic quantum
computing (Farhi and colleagues, 2000) and in the variational quantum eigensolver,
VQE (Peruzzo and colleagues, 2014). Andrew Lucas showed in 2014 how to write many
NP-hard problems this way.

\textbf{What it fixes:} it is natural for optimization, and for chemistry and
physics, where the Hamiltonian is the problem.

\textbf{What it costs:} the Hamiltonian must be one the hardware can implement, and
there is no guarantee that the system reaches its lowest-energy state in reasonable
time.
\end{concept}

\section{The oracle}

In many algorithms the problem shows up as an \emph{oracle}: a unitary $U_f$ that
implements the function $f$ of the problem. Given $\ket{x}$, it marks the state
according to $f(x)$, usually by flipping its sign. This is how Deutsch--Jozsa,
Bernstein--Vazirani and Grover are stated. The answers are not in the initial
state; they are in the effect the oracle has on the superposition. Here the encoding
of the problem is the construction of $U_f$.

That is also the trick hidden in Chapter~\ref{ch:qubit}: you can flip the sign of
the right answer without knowing it, as long as you can \emph{check} an answer.
The oracle is a circuit that checks. But an oracle is an assumption, not a gift.
On real hardware someone has to build $U_f$ out of gates, and its cost has to be
counted.

\section{Three costs}

Every encoding pays three costs, and they pull against each other:

\begin{itemize}
\item \textbf{qubits}: how many the mapping needs;
\item \textbf{depth}: how many sequential operations it takes to prepare the state
or apply the oracle;
\item \textbf{readout}: how much information each measurement gives, and how many
measurements you need.
\end{itemize}

Amplitude encoding uses few qubits but deep preparation and many measurements.
Basis encoding uses more qubits and reads out directly. There is no encoding that
is best in general, only one that suits the problem and the hardware. The thesis
behind this whole book is that this step, turning a real problem into states,
operators and measurements, is the hardest part of applying a quantum computer, and
the one engineers are least prepared for.

\chapter{From Fourier to the Quantum Fourier Transform}\label{ch:fourier}

\begin{etymology}
\textbf{Period} comes from the Greek \emph{periodos}: \emph{peri-}, ``around'',
and \emph{hodos}, ``way''. A way around, a circuit that comes back to where it
started.

\textbf{Frequency} comes from the Latin \emph{frequens}, ``crowded'',
``repeated''. How often something comes back.
\end{etymology}

Shor's algorithm (Chapter~\ref{ch:encoding}) ends with a quantum Fourier transform
that turns a hidden period into a measurable number. To see why that works, you
have to see what the ordinary Fourier transform does. It was invented to study heat,
and it turned out to be the natural tool for anything periodic.

\section{The problem that started it}

\begin{history}{1807--1822 \textperiodcentered{} Grenoble and Paris}
Joseph Fourier, prefect of the Isère department under Napoleon, studied how heat
spreads through a solid. The equation he wrote had simple solutions shaped like
sines and cosines, but real temperature profiles are not sine-shaped. He claimed
that any function can be written as a sum of sines and cosines, and submitted the
memoir to the Institut de France in 1807. Lagrange objected: a sum of smooth
oscillations, he argued, could not produce a function with corners. The claim was
not fully rigorous, and Peter Gustav Lejeune Dirichlet only proved in 1829 under
which conditions it holds. Fourier published his \emph{Théorie analytique de la
chaleur} in 1822. The idea survived because it worked.
\end{history}

Decomposing a signal into frequencies separates its scales. Each frequency can be
handled on its own and the results put back together, which is often much simpler
than handling the whole signal at once. The central question becomes: how much of
each frequency is in this signal?

\section{Fourier series}

A periodic function with period $T$ can be written as
\[
  f(t) = \sum_{k=-\infty}^{\infty} c_k\, e^{i 2\pi k t / T},
  \qquad
  c_k = \frac{1}{T}\int_0^T f(t)\, e^{-i 2\pi k t/T}\, dt .
\]
Each $c_k$ is how much of the frequency $k/T$ is present in $f$. The exponentials
are used instead of separate sines and cosines because each one packs both into a
single rotation (Chapter~\ref{ch:hilbert}).

This is linear algebra in disguise. Functions form a vector space with the inner
product $\braket{f|g} = \int_0^T \overline{f(t)}\, g(t)\, dt$, and in that space the
exponentials $e^{i2\pi kt/T}$ are orthonormal: the inner product of two different
ones is zero. They are a basis, and the $c_k$ are the coordinates of $f$ in it.
Computing $c_k$ is projecting $f$ onto one pure oscillation.

\section{The discrete Fourier transform}

A computer has no continuous functions, only $N$ samples $x_0, \dots, x_{N-1}$. The
\emph{discrete Fourier transform} (DFT) maps them to $N$ complex coefficients:
\[
  X_k = \sum_{j=0}^{N-1} x_j\, \omega_N^{-jk}, \qquad \omega_N = e^{i 2\pi / N},
\]
and the inverse brings them back,
$x_j = \frac1N\sum_k X_k\,\omega_N^{jk}$. $X_k$ measures how much the sequence
correlates with the oscillation of frequency $k/N$.

For $N = 4$ and the sequence $[1, 0, 1, 0]$, the transform is $[2, 0, 2, 0]$. The
sequence repeats every 2 samples, and the spectrum has weight only at $k = 0$ and
$k = 2$.

\section{The DFT is a unitary matrix}

The DFT is linear, so it is a matrix. With the normalization chosen to make it
symmetric,
\[
  (F_N)_{kj} = \frac{1}{\sqrt N}\, \omega_N^{jk},
\]
and $F_N$ is unitary: $F_N^\dagger F_N = I$. It is a change of basis that preserves
length, from the standard basis to the \emph{Fourier basis}, whose vectors are pure
frequencies.

For $N = 2$,
\[
  F_2 = \frac{1}{\sqrt2}\begin{pmatrix} 1 & 1 \\ 1 & -1 \end{pmatrix} = H .
\]
The Hadamard gate is the Fourier transform of one qubit. The first gate in almost
every quantum circuit is already a Fourier transform. That is why $H\ket0$ is an
even superposition: the state $\ket0$ is a single spike, and the transform of a
spike is flat.

\begin{history}{1965 \textperiodcentered{} IBM and Princeton}
Multiplying by $F_N$ directly costs $O(N^2)$ operations. In 1965 James Cooley and
John Tukey published an algorithm that exploits the recursive structure of $F_N$ to
do it in $O(N\log N)$: the fast Fourier transform. Richard Garwin had pushed for it
because detecting Soviet nuclear tests from seismometer data needed fast spectral
analysis. Gauss had found essentially the same method around 1805, in an unpublished
notebook on asteroid orbits.
\end{history}

\section{The DFT finds periods}

This is the property that matters for the rest of the book. If a sequence of length
$N$ has period $r$, and $r$ divides $N$, its transform is zero everywhere except at
multiples of $N/r$. The $N/r$ copies of the period add up in phase at those
frequencies and cancel exactly at all the others. It is interference, the same
mechanism as in Chapter~\ref{ch:qubit}.

Take $N = 8$ and a spike every 4 samples:
\[
  [1, 0, 0, 0, 1, 0, 0, 0] \;\xrightarrow{\ \text{DFT}\ }\; [2, 0, 2, 0, 2, 0, 2, 0].
\]
The peaks are at $0, 2, 4, 6$, the multiples of $8/4 = 2$. The
spectrum of a periodic signal is sparse, and the spacing of its peaks reveals the
period.

Apply this to $f(a) = 2^a \bmod 15$, which has period 4 (Chapter~\ref{ch:factoring}), and the
transform of its first eight values has peaks at $k = 2, 4, 6$ (the constant part
sits at $k = 0$). Spacing 2, so $r = 8/2 = 4$.

\section{The quantum Fourier transform}

The quantum Fourier transform (QFT) is the same matrix $F_N$, with $N = 2^n$, built as
a circuit on $n$ qubits:
\[
  \mathrm{QFT}\,\ket{j} = \frac{1}{\sqrt N}\sum_{k=0}^{N-1} \omega_N^{jk}\,\ket{k}.
\]
Applied to a superposition $\sum_j x_j\ket j$, it produces $\sum_k X_k\ket k$: the DFT
of the amplitudes.

\begin{history}{1994 \textperiodcentered{} IBM}
Shor's first construction needed a Fourier transform over a large range of values,
and Don Coppersmith, at IBM, showed in 1994 how to build it from Hadamards and
small controlled phase rotations. Coppersmith's report was titled ``An approximate
Fourier transform useful in quantum factoring''. He also showed that the tiniest
rotations can be dropped with little loss of accuracy.
\end{history}

The circuit uses $O(n^2)$ gates: one Hadamard per qubit, controlled rotations
$R_k = \mathrm{diag}(1, e^{i2\pi/2^k})$ between pairs, and swaps at the end to
reverse the order of the qubits. For three qubits, that is three Hadamards, three
controlled rotations and one swap, and multiplying their matrices together gives
exactly $F_8$.

Compare the costs. The classical FFT on $N = 2^n$ numbers takes $O(N\log N) = O(2^n n)$
operations. The QFT takes $O(n^2)$ gates. That looks like an exponential speedup, and
in a narrow sense it is.

\begin{concept}{Quantum Fourier transform}
The discrete Fourier transform applied to the $2^n$ amplitudes of an $n$-qubit
state, with a circuit of $O(n^2)$ gates.

\textbf{What it fixes:} it makes periodicity hidden in an exponentially large
superposition visible to a measurement.

\textbf{What it costs:} you cannot read the spectrum. Measuring returns a single
$k$, with probability $|X_k|^2$. The QFT is not a faster FFT.
\end{concept}

\section{What you can and cannot read}

After the QFT the state is $\sum_k X_k\ket k$, and a measurement gives one $k$. If
the spectrum is flat, that $k$ says nothing. If the spectrum has a few sharp peaks,
almost every measurement lands on one of them, and a few runs are enough to infer
the spacing. The QFT is useful exactly when the spectrum is sparse: periodic
functions, or phases of eigenvalues.

This is why every claim of the form ``the QFT gives an exponential speedup for
signal processing'' should be read with suspicion. Loading a classical signal into
the amplitudes costs about as much as the classical FFT, and the output cannot be
read. The QFT pays off only when the periodic state is produced inside the
computer, as in Shor's algorithm, and only a small piece of information (a period,
a phase) needs to come out.

\section{Two uses}

\textbf{Shor.} The state $\sum_a \ket a\ket{x^a \bmod N}$ is periodic in $a$, with
period $r$. A QFT on the first register puts the probability on multiples of $q/r$,
and a measurement plus continued fractions gives $r$ (Chapter~\ref{ch:encoding}).

\textbf{Phase estimation.} A unitary $U$ with eigenvector $\ket{u}$ and eigenvalue
$e^{i2\pi\phi}$ can be used to write the phase $\phi$ into a control register, as
$\frac{1}{\sqrt N}\sum_k e^{i2\pi k\phi}\ket k$: a pure oscillation with frequency
$\phi$. The inverse QFT concentrates it on the integer closest to $N\phi$. Alexei
Kitaev introduced phase estimation in 1995, and in 1998 Cleve, Ekert, Macchiavello
and Mosca showed that Shor's algorithm is a special case of it. Phase estimation is
how a quantum computer reads eigenvalues, including the energies of a Hamiltonian,
which brings the story back to Feynman's original question.


\part{Hardware}\label{part:hardware}
\chapter{What a Qubit Is Made Of}\label{ch:hardware}

\begin{etymology}
\textbf{System} comes from the Greek \emph{sýstēma}: \emph{syn-}, ``together'',
and \emph{histanai}, ``to set up''. Parts set up together as a whole, which we
choose to separate from the rest of the world in order to study it.

\textbf{Coherence} comes from \emph{cohaerere}: \emph{co-}, ``together'', and
\emph{haerere}, ``to stick''. In quantum physics it means that the phase
relations between amplitudes are kept, which is what allows interference.
\emph{Decoherence} is the loss of those relations, when the system gets entangled
with its surroundings and the phases leak out.
\end{etymology}

A qubit is any physical system with two levels that obeys quantum mechanics: the
spin of an electron (up or down), the polarization of a photon (horizontal or
vertical), two energy levels of an atom. Mathematically all of them are the same
vector in $\mathbb{C}^2$. Physically they could hardly be more different, and the
differences decide everything about how far a technology can go.

The hard part is not finding a two-level system. Nature is full of them. The hard
part is finding one that you can prepare, control and read, and that still keeps
its phases long enough to compute. Those demands pull against each other. A qubit
well isolated from the world keeps its coherence, but it is also hard to reach
with the controls. A qubit easy to control is, for the same reason, easy for the
environment to disturb.

This chapter is a tour of the technologies, and a checklist guides it. Every
platform is walked through the same list, item by item, so that by the end each
item means something concrete.

\section{The checklist}

\begin{history}{2000 \textperiodcentered{} IBM}
David DiVincenzo, at IBM, listed what a physical system needs to be a quantum
computer. He had written, with Daniel Loss, one of the first concrete proposals for
qubits in 1998. His list became the checklist every technology is measured against.
\end{history}

\begin{enumerate}
\item \textbf{Characterized and scalable.} You know exactly what your qubit is, and
adding more of them does not make the others worse.
\item \textbf{Initialization.} You can put every qubit in a known starting state,
$\ket{00\dots0}$.
\item \textbf{Coherence.} The qubit keeps its quantum state much longer than a gate
takes.
\item \textbf{Universal gates.} A small menu of operations is enough to build any
computation.
\item \textbf{Measurement.} You can read each qubit, reliably.
\end{enumerate}

He added two more for quantum communication: (6) converting between qubits that
stay put and qubits that travel, which in practice are photons, and (7) sending the
traveling ones faithfully from one place to another.

In 2018, John Preskill gave a name to the machines that exist today: \emph{NISQ},
noisy intermediate-scale quantum. Tens to a few thousand qubits, too noisy for error
correction at scale, too large to simulate exactly. No technology yet meets the five
criteria at the fidelity and scale that Shor's algorithm needs on a cryptographic
number. The tour shows where each one falls short.

\section{Superconducting circuits}

\begin{history}{1962--1985 \textperiodcentered{} Cambridge and Berkeley}
In 1962 Brian Josephson, a 22-year-old PhD student at Cambridge, predicted that a
current can flow without resistance through a thin insulating barrier between two
superconductors. The \emph{Josephson junction} earned him a share of the 1973
Nobel Prize. In 1984 and 1985, at Berkeley, John Clarke, Michel Devoret and John
Martinis showed that an electrical circuit built around such a junction, big enough
to see under a microscope, behaves as a single quantum object: it tunnels through
energy barriers and has discrete energy levels that they could probe with
microwaves. A circuit with billions of electrons had a quantized spectrum. They
received the Nobel Prize in Physics for it in 2025.
\end{history}

\subsection{What the qubit is}

Start with the simplest circuit there is: a capacitor and an inductor in a loop. Charge
sloshes from the capacitor through the inductor and back, and the current oscillates
at a fixed frequency. It is the mass on a spring of Chapter~\ref{ch:classical}, with
charge in place of position and current in place of velocity. Cool it until the metal
superconducts, so nothing is lost to resistance, and it becomes a quantum harmonic
oscillator: its energy can only take values on a ladder of equally spaced rungs.

\begin{reflection}{The energy of a whole circuit comes in steps?}
That sentence should bother you. Quantization was discovered in atoms. A circuit is a
piece of metal you can see under a microscope, made of billions of electrons, and we
are saying its energy cannot change smoothly.

First, which energy. Not the energy of each electron, and not heat. It is the energy
of the oscillation itself, the one that sloshes back and forth. At one instant it sits
in the electric field of the charged capacitor, $Q^2/2C$; a quarter of a period later
it sits in the magnetic field of the current through the inductor, $\Phi^2/2L$, where
$\Phi$ is the magnetic flux. Their sum is the Hamiltonian of the circuit,
\[
  H = \frac{Q^2}{2C} + \frac{\Phi^2}{2L},
\]
the same shape as the spring's $p^2/2m + kx^2/2$ in Chapter~\ref{ch:classical}: charge
plays the role of momentum, flux the role of position, and $1/C$ and $1/L$ the roles of
$1/m$ and $k$. So far this is classical physics. Three steps turn it into quantum
physics, and each one was built earlier in the book.

\textbf{Dirac's rule} (end of Chapter~\ref{ch:classical}). Replace each classical
quantity with an operator, and the Poisson bracket with a commutator. For the circuit,
$\{\Phi, Q\} = 1$ becomes
\[
  [\hat\Phi, \hat Q] = \hat\Phi\hat Q - \hat Q\hat\Phi = i\hbar .
\]
Charge and flux stop being numbers. They become operations on the state of the
circuit, and the order in which you apply them matters.

\textbf{Uncertainty} (Chapter~\ref{ch:quantum}). Two operators that do not commute
have no eigenvectors in common, so no state of the circuit has a definite charge and a
definite flux at the same time. Robertson's inequality gives
$\Delta\Phi\,\Delta Q \geq \hbar/2$. The circuit can never be perfectly at rest, with
no charge and no current at once.

\textbf{Energies are eigenvalues} (Chapter~\ref{ch:quantum}, third postulate). The
energy is now the operator $\hat H$, and a measurement of it can only return one of
its eigenvalues. For the particle in a box they came out discrete because only waves
that fit between the walls are allowed. The circuit is a box with soft walls: the
energy $\Phi^2/2L$ grows on both sides, so the wave describing the state has to die
out on both sides, and only some shapes manage it. Solving for those shapes gives a
ladder of energies, equally spaced, $\hbar\omega$ apart, with
$\omega = 1/\sqrt{LC}$.

If one of these steps lost you, the problem is not in this box. You skipped
Chapter~\ref{ch:quantum}, or the end of Chapter~\ref{ch:classical}, where each step is
built from its history. Go back; this box only applies them to a circuit.

So what is a rung of the ladder? Not a place, and not a part of the circuit. Each rung
is one value that the total energy of the oscillation is allowed to have:
\[
  E_n = \hbar\omega\left(n + \tfrac12\right), \qquad n = 0, 1, 2, \dots
\]
The number $n$ counts how many quanta of oscillation the circuit is holding, and each
quantum carries exactly the energy of one microwave photon of frequency $\omega/2\pi$.
Climbing one rung means absorbing one such photon; dropping one means emitting it. There
is nothing in between: the circuit cannot hold half a quantum. Even the bottom rung is not
zero. The circuit keeps $\tfrac12\hbar\omega$ there, the \emph{zero-point energy}, because
charge and flux cannot both be exactly zero at once; that would violate the uncertainty
principle of Chapter~\ref{ch:quantum}.

Second, why nobody sees the steps in ordinary electronics. They are there, but tiny.
At 5~GHz one step is $3.3\times10^{-24}$~J. A capacitor of 1~pF charged to 1~V stores
$5\times10^{-13}$~J, which is 150 billion steps: a ramp made of grains of sand looks
smooth. At room temperature, thermal agitation is 1250 times larger than one step, so
the state is smeared over a wide band of rungs and no single one can be picked out. And
resistance drains the oscillation before any step could be noticed.

\begin{center}
\begin{tikzpicture}[x=1mm, y=1mm]
  % ---------- ordinary circuit, warm ----------
  \begin{scope}
    \foreach \i in {0,...,64} \draw[black!35, line width=0.25pt] (0,\i*0.5) -- (22,\i*0.5);
    \shade[top color=hfAlert!0, bottom color=hfAlert!0, middle color=hfAlert!55, opacity=0.75]
      (-0.5,9) rectangle (22.5,27);
    \draw[black!60, decorate, decoration={brace, amplitude=3pt}] (23.5,27) -- (23.5,9)
      node[midway, right=3pt, hf/note, align=left] {thermal\\spread:\\$\sim$1250\\rungs};
    \node[font=\sffamily\small, above, align=center] at (11,33.5) {ordinary circuit, 300~K};
    \node[hf/note, below, align=center] at (11,-1.5) {rungs too close to tell apart;\\the state is smeared\\(not to scale)};
  \end{scope}
  % arrow
  \draw[black!40, -{Stealth[length=2.4mm]}, line width=0.8pt] (44,16) -- (70,16)
    node[midway, above, hf/note, align=center] {superconducting,\\cooled to 15~mK};
  % ---------- same circuit, cold ----------
  \begin{scope}[xshift=92mm]
    \foreach \n/\e in {0/{\tfrac12}, 1/{\tfrac32}, 2/{\tfrac52}, 3/{\tfrac72}} {
      \draw[black!70, line width=1pt] (0,\n*9) -- (20,\n*9);
      \node[font=\small, left] at (0,\n*9) {$n = \n$};
      \node[font=\small, right] at (20,\n*9) {$\e\,\hbar\omega$};
    }
    \fill[hfConcept] (10,1.6) circle (1.6);
    \draw[hfConcept, -{Stealth[length=1.8mm]}, line width=0.8pt] (15,0.6) -- (15,8.4);
    \node[hf/note, text=hfConcept, anchor=west] at (15.5,4.5) {one photon};
    \node[font=\sffamily\small, above, align=center] at (10,33.5) {same circuit, cold};
    \node[hf/note, below, align=center] at (10,-1.5) {each rung is one allowed energy;\\the circuit sits on the bottom one};
  \end{scope}
\end{tikzpicture}
\end{center}

The transmon removes all three disguises. It is cooled far below the size of one step,
it is superconducting so nothing drains away, and it holds zero or one quantum, not
billions. What remains is the bare ladder.

The last reason it works is superconductivity itself. Below the critical temperature,
the electrons pair up and all the pairs share a single quantum state. Billions of
particles move as one, so the whole circuit has a single degree of freedom, and that
single degree of freedom behaves like one particle on a spring. In 1980 Anthony Leggett
asked whether quantum mechanics would still hold for a variable like that, describing a
macroscopic object. The experiments of Clarke, Devoret and Martinis in 1985 were the
answer: it does.
\end{reflection}

That ladder is useless as a qubit. A microwave pulse tuned to the jump from the bottom
rung to the next has exactly the frequency of every other jump too, so it pushes the
circuit up the ladder and out of any pair of levels you wanted to use.

Now replace the inductor with a Josephson junction. The junction behaves like an
inductor whose value depends on the current through it, a \emph{nonlinear} one. It
bends the oscillator so that the rungs get closer together as you climb. The first jump
and the second now have different frequencies, and a pulse tuned to the first leaves
the second alone.

\begin{figure}[htbp]
\centering
\begin{tikzpicture}[x=1mm, y=1mm]
  % ---------- (a) LC oscillator ----------
  \begin{scope}
    \draw[line width=0.8pt] (0,4) -- (0,0) -- (16,0) -- (16,3);
    \draw[line width=0.8pt] (0,6) -- (0,10) -- (16,10) -- (16,7);
    \draw[line width=1.2pt] (-3,4) -- (3,4) (-3,6) -- (3,6);
    \draw[line width=0.8pt, decorate, decoration={coil, aspect=0.6, segment length=1.6mm, amplitude=1.6mm}] (16,7) -- (16,3);
    \node[hf/note, left] at (-3,5) {$C$};
    \node[hf/note, right] at (18,5) {$L$};
    \node[font=\sffamily\small, above] at (8,11) {(a) LC oscillator};
    % ladder
    \foreach \lvl/\name in {0/0, 9/1, 18/2, 27/3} {
      \draw[black!70, line width=1pt] (0,-40+\lvl) -- (16,-40+\lvl);
      \node[font=\small, left] at (0,-40+\lvl) {$\ket{\name}$};
    }
    \foreach \a/\b in {0/9, 9/18, 18/27}
      \draw[hfConcept, -{Stealth[length=1.8mm]}, line width=0.8pt] (8,-40+\a+0.6) -- (8,-40+\b-0.6);
    \node[hf/note, below, align=center] at (8,-42) {every jump the same:\\a pulse climbs the whole ladder};
  \end{scope}
  % arrow between panels
  \draw[black!40, -{Stealth[length=2.4mm]}, line width=0.8pt] (27,5) -- (51,5)
    node[midway, above, hf/note] {swap $L$ for a junction};
  % ---------- (b) transmon ----------
  \begin{scope}[xshift=70mm]
    \draw[line width=0.8pt] (0,4) -- (0,0) -- (16,0) -- (16,10) -- (0,10) -- (0,6);
    \draw[line width=1.4pt] (-5,4) -- (5,4) (-5,6) -- (5,6);
    \draw[hfAlert, line width=1.2pt] (13.5,3.5) -- (18.5,6.5) (13.5,6.5) -- (18.5,3.5);
    \node[hf/note, left] at (-5,5) {large $C$};
    \node[hf/note, right, text=hfAlert] at (19,5) {junction};
    \node[font=\sffamily\small, above] at (8,11) {(b) transmon};
    % ladder
    \fill[hfConcept!10] (-1,-41.5) rectangle (17,-30);
    \foreach \lvl/\name in {0/0, 9/1, 16.5/2, 22.5/3} {
      \draw[black!70, line width=1pt] (0,-40+\lvl) -- (16,-40+\lvl);
      \node[font=\small, left] at (0,-40+\lvl) {$\ket{\name}$};
    }
    \draw[hfConcept, -{Stealth[length=1.8mm]}, line width=1pt] (8,-39.4) -- (8,-31.6);
    \draw[hfAlert, -{Stealth[length=1.8mm]}, line width=0.8pt, dashed] (8,-30.4) -- (8,-24.1);
    \draw[black!30, -{Stealth[length=1.8mm]}, line width=0.8pt] (8,-22.9) -- (8,-18.1);
    \draw[hfConcept, line width=0.8pt, decorate, decoration={brace, amplitude=4pt, mirror}]
      (18.5,-40) -- (18.5,-31);
    \node[hf/note, anchor=west, text=hfConcept, align=left] at (22,-35.5) {\textbf{the qubit}\\$\approx 5$ GHz};
    \node[hf/note, anchor=west, text=hfAlert, align=left] at (22,-27) {$\ket1\to\ket2$ is 200--300 MHz\\lower: the pulse misses it};
    \node[hf/note, below, align=center] at (8,-42) {each jump a bit smaller};
  \end{scope}
\end{tikzpicture}
\caption{What a superconducting qubit is. (a) A capacitor and an inductor make an
oscillator whose energy levels are equally spaced. (b) Replace the inductor with a
Josephson junction and the levels become uneven. The qubit is the two lowest levels of
this circuit.}\label{fig:ladder}
\end{figure}

So the qubit is not a particle or a little object on the chip. It is the two lowest
energy states of a whole electrical circuit, a few hundred micrometers across, with
billions of paired electrons moving as one. $\ket0$ is the circuit at rest in its
ground state. $\ket1$ is the circuit holding one quantum of oscillation, an energy of
about 20 micro-electronvolts, which is what a single microwave photon at 5~GHz
carries. A superposition $\alpha\ket0 + \beta\ket1$ is the circuit in both at once.

The first circuits of this kind came in quick succession. In 1999, at NEC in Tsukuba,
Yasunobu Nakamura, Yuri Pashkin and Jaw-Shen Tsai showed coherent oscillations in a
``Cooper-pair box'', a tiny island of superconductor; they lasted a few nanoseconds.
The same year Hans Mooij's group in Delft proposed the flux qubit, a loop with current
flowing both ways at once, seen in experiments in 2000. In 2002, at Saclay, Denis Vion,
Daniel Esteve, Michel Devoret and colleagues built the \emph{quantronium}, and at NIST
in Boulder, Martinis's group demonstrated the phase qubit. In 2007, at Yale, Jens Koch,
Robert Schoelkopf, Steven Girvin and colleagues proposed the \emph{transmon}: a
Cooper-pair box with a large capacitor added, which makes it nearly immune to stray
charges at the price of a smaller nonlinearity. The transmon became the dominant
design, used by IBM and Google among others. Now the checklist.

\subsection*{1. Characterized and scalable}

\emph{Characterized} means you know the Hamiltonian of your qubit: the energies of its
levels, how it couples to its neighbors, and what else is nearby. For a transmon that is
the ladder of Figure~\ref{fig:ladder}. The qubit frequency is typically around 5~GHz,
the next jump is 200 to 300~MHz lower, and the third rung is always there, waiting for
a pulse that is too strong or too short. A transmon that ends up in $\ket2$ has left the
computation, an error called \emph{leakage}.

Transmons are made, not found, and no two come out of fabrication identical. Each one
has its own frequency and must be measured and calibrated on its own, and the
calibration drifts. Characterization is not done once; it is a daily job.

\emph{Scalable} means adding qubits does not make the rest worse. The transmon's best
argument is here: chips are made with the lithography of ordinary electronics, so
drawing a thousand qubits costs little more than drawing ten. Its worst problem is also
here. Every qubit needs its own control lines, running from room temperature down into
the refrigerator, and each wire brings heat and noise with it. Hundreds of qubits fit;
a million wires do not.

\subsection*{2. Initialization}

Before a computation, every qubit has to be in $\ket0$. Nature does that for free if
the qubit is cold enough, because a system in equilibrium sits in its lowest level
except for thermal kicks. The kicks are small when the temperature is small compared
with the energy of the jump. For a 5~GHz transmon, that energy corresponds to a
temperature of 0.24~K. Boltzmann's law gives the chance of finding the qubit in $\ket1$:

\begin{center}
\begin{tabular}{lr}
\hline
temperature & qubit found in $\ket1$ \\
\hline
room temperature, 300~K & 49.98\% \\
liquid helium, 4~K & 48.5\% \\
50~mK & 0.8\% \\
15~mK & 0.00001\% \\
\hline
\end{tabular}
\end{center}

At room temperature the qubit is a coin toss. At 4~K, the temperature of liquid helium,
it is still a coin toss. Only well below 0.24~K does it settle in $\ket0$, and that is
why superconducting chips live in dilution refrigerators at 10 to 20~millikelvin, colder
than outer space. In practice the chip runs warmer than the refrigerator, and a fraction
of a percent of qubits start in $\ket1$, so machines often reset them actively before
each run: measure, and flip the ones that came out wrong.

\subsection*{3. Coherence}

A transmon in $\ket1$ eventually drops to $\ket0$, handing its quantum of energy to the
environment. The probability that it has survived after a time $t$ is $e^{-t/T_1}$, where
$T_1$ is the \emph{relaxation time}. The relative phase between $\ket0$ and $\ket1$ also
drifts and is lost, on a time scale called $T_2$. Interference needs both.

What the criterion asks for is a ratio, not a time. A transmon with $T_1 = 100\,\mu$s
running 50~ns gates loses about 0.05\% per gate to relaxation alone, and fits about 2000
gates into one $T_1$. The best transmons reported in 2025, with $T_1$ around a
millisecond, fit about 20\,000. In 1999 Nakamura's box lasted a few nanoseconds; in
twenty-five years coherence improved by about six orders of magnitude.

Shor's algorithm on RSA-2048 needs billions of operations. A ratio of $10^4$ is not
close, and no improvement in materials will close five orders of magnitude. That is why
error correction, the subject of Part~\ref{part:qec}, is not optional.

\subsection*{4. Universal gates}

A one-qubit gate on a transmon is a microwave pulse at the qubit frequency. While the
pulse is on, the state turns steadily from $\ket0$ toward $\ket1$ and back, a rotation on
the Bloch sphere; the length of the pulse sets the angle and its phase sets the axis.

\begin{history}{1937 \textperiodcentered{} Columbia}
Isidor Rabi sent beams of molecules through an oscillating magnetic field and saw their
nuclear spins flip when the field was tuned to the right frequency. The steady flipping
back and forth under a resonant drive is called a \emph{Rabi oscillation}, and Rabi's
method won him the 1944 Nobel Prize. It became nuclear magnetic resonance, then MRI.
Nakamura's result of 1999 was a Rabi oscillation in a circuit, and every single-qubit gate
on every platform since is one too.
\end{history}

Two-qubit gates need two transmons to interact, through a capacitor or a tunable coupler
between them. That only works between neighbors on the chip. A gate between distant
qubits has to be built from a chain of swaps, which adds time and error.

Is that menu enough? \emph{Universal} means that a small, fixed set of gates can
approximate any unitary to any precision, and the right set is not obvious.

\begin{history}{1998 \textperiodcentered{} the Gottesman--Knill theorem}
Daniel Gottesman, building on work with Emanuel Knill, proved that circuits made only of
Hadamard, CNOT and the phase gate $S = \mathrm{diag}(1, i)$, plus measurements, can be
simulated efficiently on a classical computer, no matter how many qubits they use or how
much entanglement they create. These are the \emph{Clifford} gates, and with them alone a
quantum computer is no more powerful than a laptop.
\end{history}

Add a single gate, $T = \mathrm{diag}(1, e^{i\pi/4})$, a $45^\circ$ turn about the vertical
axis of the Bloch sphere, and the set becomes universal. Robert Solovay and Alexei Kitaev
showed independently, in 1995 and 1997, that any one-qubit gate can then be approximated
to precision $\varepsilon$ by a sequence of $H$ and $T$ gates whose length grows only like
a power of $\log(1/\varepsilon)$. A transmon can do any rotation directly, so it passes this
criterion easily today. Under error correction it will not: rotations become sequences of
discrete gates, and the $T$ gate is by far the most expensive of them.

\subsection*{5. Measurement}

A transmon is read through a microwave resonator coupled to it. The resonator's frequency
shifts slightly depending on whether the qubit is in $\ket0$ or $\ket1$, so a weak pulse
bounced off it comes back changed. This scheme, \emph{circuit quantum electrodynamics},
was proposed by Alexandre Blais and colleagues at Yale in 2004. It takes hundreds of
nanoseconds, during which the qubit may relax and give the wrong answer.

Readout errors compound. If each qubit is read correctly 99\% of the time, all 100 qubits
of a register come out right only 37\% of the time. At 99.9\% the chance rises to 90\%.
Every fraction of a percent at the end of a computation is paid again for every qubit.

\subsection*{6 and 7. Moving qubits}

Here the transmon is weakest. Its quantum lives in microwave photons, which cannot leave
the refrigerator without being swamped by thermal noise; at room temperature, as the
table above shows, a microwave mode is a coin toss. Linking two refrigerators needs a
converter from microwave to optical photons that preserves quantum states, and in 2025
that was still a laboratory problem.

\section{Trapped ions}

\begin{history}{1995 \textperiodcentered{} Innsbruck and Boulder}
Ignacio Cirac and Peter Zoller, in Innsbruck, proposed the first concrete architecture:
ions held in a line by electric fields in a vacuum, with two levels of each ion as the
qubit, and laser pulses as the gates. The ions interact through their shared vibrations.
The same year, Christopher Monroe and David Wineland, at NIST in Boulder, demonstrated a
controlled-NOT gate on a single beryllium ion. Wineland shared the 2012 Nobel Prize, with
Serge Haroche, for manipulating individual quantum systems.
\end{history}

Here the qubit is an atom that has lost one electron. Its outermost electron has two
energy levels that differ by the orientation of its spin relative to the nucleus. In the
ytterbium ion $^{171}\mathrm{Yb}^+$, one of the most used, the two levels are 12.6~GHz
apart. Nature built the ladder; nobody had to bend it. The same checklist gives a very
different picture.

\subsection*{1. Characterized and scalable}

Every ytterbium ion in the universe is identical, so every qubit has exactly the same
frequency, known to many decimal places. Characterization comes for free, and this is the
ions' great advantage over anything fabricated.

Scaling is their problem. The ions sit in a line, held apart by their mutual repulsion and
coupled by the vibrations of the whole chain. Add ions and the chain gets longer, its
vibrations crowd together, and gates get slower and harder to aim. A single chain stops
being practical at a few dozen ions. The plan since 2002, proposed by David Kielpinski,
Monroe and Wineland, is to shuttle ions between zones of a trap, like trains in a yard, so
that gates always act on small groups.

\subsection*{2. Initialization}

Compare with the transmon. A 12.6~GHz qubit at room temperature is in $\ket1$ about half the
time, exactly the coin toss of the table above, and ion traps run at room temperature. They
do not need cold, because they do not rely on equilibrium. A laser tuned to excite the ion
from every state except $\ket0$ keeps kicking it until it falls into $\ket0$, where the laser
can no longer touch it. This \emph{optical pumping} prepares the qubit with probability above
99.9\%. The criterion asks for a known state, not a cold one; a one-way process is enough.

\subsection*{3. Coherence}

The two levels of the ion are barely disturbed by stray fields, and the ion floats in vacuum,
touching nothing. In 2021 a single ion kept its phase for over an hour. That is about a
billion times longer than a transmon.

But criterion 3 is about a ratio. A two-qubit gate on ions takes around 100~$\mu$s, a
thousand times slower than on a transmon, so the number of gates per coherence time lands at
about $10^4$, the same range. The hour-long coherence does not buy a longer computation; it
buys a slower one.

\subsection*{4. Universal gates}

One-qubit gates are laser or microwave pulses, Rabi oscillations again. Two-qubit gates use
the chain: a laser couples the internal state of two ions to the vibration they share, and
the vibration carries the interaction from one to the other. Any ion in the chain can be
coupled to any other, so there are no swaps. In 2024, Quantinuum reported two-qubit gates
with 99.9\% fidelity, the best of any platform at the time.

\subsection*{5. Measurement}

Ions are read by a method Hans Dehmelt proposed in 1975. Shine a laser tuned so that the ion
scatters light only when it is in one of the two states. In that state it glows, millions of
photons per second; in the other it stays dark. Count the photons and you have the answer,
with fidelity above 99.9\%.

\subsection*{6 and 7. Moving qubits}

Ions emit optical photons, which travel through ordinary fiber. In 2007 Monroe's group
entangled two ions held in separate traps a meter apart, using the photons each one
emitted. This is the ions' second advantage: modular machines linked by light are a natural
extension, not a research problem.

\section{The others}\label{sec:others}

The other platforms are easiest to understand through the criterion where they break. Topological
qubits come after them, in a section of their own, because they need more explanation.

\textbf{Nuclear magnetic resonance: criteria 1 and 2.} In 1997 David Cory, Amr Fahmy and
Timothy Havel, and separately Neil Gershenfeld and Isaac Chuang, proposed using the nuclear
spins of molecules in a liquid as qubits, driven by the radio pulses of NMR machines. In 2001
Lieven Vandersypen and colleagues at IBM used seven spins to run Shor's algorithm and factor
15. But a nuclear spin at 400~MHz, at room temperature, is in its upper state with probability
49.998\%: initialization fails, and badly. NMR worked around it with statistical tricks over
trillions of molecules, at the cost of a signal that falls exponentially with the number of
qubits, which fails criterion 1. In 1999 Samuel Braunstein and colleagues showed that the NMR
experiments of the time never created entanglement at all. It was useful for testing ideas, and
a dead end for computing.

\textbf{Neutral atoms: criteria 1 and 4.} Atoms without charge, held by tightly focused laser
beams called optical tweezers. Like ions, they are identical by nature, and they can be moved,
which gives flexible connections. Scale is their strength: in 2025 a Caltech group trapped more
than six thousand atoms in one array. Trapping is not computing, though. Neutral atoms barely
interact, so two-qubit gates excite them to huge Rydberg states, an idea proposed by Dieter
Jaksch and colleagues in 2000, and those gates are slower and less accurate than one-qubit gates.
Atoms also escape the tweezers, which is a qubit vanishing from the computation.

\textbf{Photons: criterion 4.} The qubit is the polarization or the path of a photon. Photons
barely interact with the environment, which is excellent for criterion 3 and makes them the
natural answer to criteria 6 and 7. For the same reason they barely interact with each other,
which is fatal for two-qubit gates. In 2001 Emanuel Knill, Raymond Laflamme and Gerard Milburn
showed that gates can be built from mirrors, beam splitters and measurements, but they only work
with some probability, and the overhead is large.

\textbf{Spins in semiconductors: criterion 1.} Daniel Loss and DiVincenzo proposed in 1998 to use
the spin of single electrons trapped in quantum dots, tiny traps in silicon. They would scale like
transistors, which is the whole appeal. But each dot is shaped by a few atoms of disorder, so every
qubit is slightly different, and characterizing and controlling them uniformly is still the
obstacle.

\section{Topological qubits}\label{sec:topological}

Every platform so far stores the qubit in something local: two levels of a circuit, of an ion,
of an atom, of a spin. Then it defends that local thing from the world, first by isolating it
and in the end by error correction, which Part~\ref{part:qec} will show costs hundreds of
physical qubits per logical one. The topological approach tries to store the qubit somewhere
local noise cannot reach. It has not yet passed criterion 1: there is no undisputed
demonstration that the qubit exists, and Chapter~\ref{ch:majorana} tells that story. This
section explains how the device is supposed to work, because the claims in that chapter only
make sense against it. It uses the bosons, fermions and anyons of Section~\ref{sec:identical}.

\begin{concept}{Topological protection}
Information stored in a global property of a system, such as the joint parity of two
distant zero modes, so that no local disturbance can change it.

\textbf{What it fixes:} criterion 3 by construction. Instead of fighting decoherence only
with many physical qubits per logical qubit, the hardware itself does part of the job.

\textbf{What it costs:} materials that nobody knew how to make, a protection that only
holds while an energy gap stays open and no electron jumps in from outside, and a gate set
that is not universal.
\end{concept}

\subsection*{The gap}

In a superconductor the electrons bind into Cooper pairs, and the pairs, being bosons, all
share one collective state. Leon Cooper showed in 1956 that the pairs form, and in 1957 John
Bardeen, Cooper and Robert Schrieffer built the full theory. Breaking a pair costs energy: at
least $\Delta$ for each electron set free, the \emph{superconducting gap}. Below the gap there
is nothing to excite, which is why a superconducting circuit is so quiet. The BCS theory gives
$\Delta \approx 1.76\,k_B T_c$, with $T_c$ the critical temperature. Aluminum, with $T_c =
1.2$~K, has a gap of about $180\,\mu$eV. Lead, with $T_c = 7.2$~K, has about 1.4~meV,
somewhat more than the BCS estimate of 1.1~meV.

At 15~mK, $k_B T$ is $1.3\,\mu$eV, a hundred and forty times smaller than the gap of
aluminum, so heat alone should break essentially no pairs. In practice a few pairs are always
broken anyway, by stray radiation and by heat leaking down the wiring, far more than the
temperature predicts. The free electrons they leave, called \emph{quasiparticles}, wander
through the metal. For a transmon they are one source of error among many. For a topological
qubit they are the main enemy, as the next figure shows.

\subsection*{Half of an electron}

An excitation of a superconductor is not quite an electron. It is a mixture of an electron and
a \emph{hole}, the absence of an electron, which behaves as a particle of opposite charge. A
hole is to an electron what an antiparticle is to a particle. In some superconductors a special
excitation can appear that is an exact half-and-half mixture, at zero energy. It is its own
antiparticle, which is the property Ettore Majorana found in an equation in 1937, and so it is
called a \emph{Majorana zero mode}. Its operator equals its own adjoint, $\gamma = \gamma^\dagger$.

One Majorana zero mode is half of an ordinary fermion state. Two of them, $\gamma_1$ and
$\gamma_2$, combine into one complete fermion mode,
\[
  c = \frac{\gamma_1 + i\gamma_2}{2},
\]
and an ordinary fermion mode can only be empty or full, by the exclusion principle. That
``empty or full'' is a bit, and its value is called the \emph{parity}.

The point is where the bit lives. It is not in either half. It is in the pair, and the halves
can be micrometers apart. Noise that touches one end of the wire, a stray field, a vibrating
atom, cannot read or change a value that only exists in the combination of both ends. One
thing can change it: an electron arriving from outside, such as a quasiparticle, which turns
an empty mode into a full one. This is called \emph{quasiparticle poisoning}, and the average
wait for it is the parity lifetime that Chapter~\ref{ch:majorana} is about.

\begin{figure}[htbp]
\centering
\begin{tikzpicture}[x=1mm, y=1mm]
  % where each half lives
  \draw[black!25, dashed, line width=0.6pt] (0,13) -- (90,13);
  \fill[hfConcept!15] (0,13) -- plot[domain=0:40, samples=60] (\x, {13 + 8*exp(-\x/5)}) -- (40,13) -- cycle;
  \fill[hfConcept!15] (50,13) -- plot[domain=50:90, samples=60] (\x, {13 + 8*exp(-(90-\x)/5)}) -- (90,13) -- cycle;
  \draw[hfConcept, line width=1pt, domain=0:40, samples=60] plot (\x, {13 + 8*exp(-\x/5)});
  \draw[hfConcept, line width=1pt, domain=50:90, samples=60] plot (\x, {13 + 8*exp(-(90-\x)/5)});
  \node[hf/note] at (45,22) {each half sits at one end;\\the middle holds nothing};
  % the device
  \draw[hfSource, fill=black!30, line width=0.8pt] (0,6) rectangle (90,10);
  \node[font=\sffamily\footnotesize, text=white] at (45,8) {superconductor};
  \draw[hfSource, fill=hfSource!10, line width=0.8pt] (0,0) rectangle (90,6);
  \node[font=\sffamily\footnotesize, text=black!70] at (45,3) {semiconductor nanowire};
  \draw[black!50, -{Stealth[length=2mm]}, line width=0.8pt] (52,30) -- (78,30);
  \node[hf/note, anchor=east] at (51,30) {magnetic field along the wire};
  % the two halves
  \fill[hfConcept] (5,3) circle (2.4);
  \fill[hfConcept] (85,3) circle (2.4);
  \node[font=\small, text=white] at (5,3) {$\gamma_1$};
  \node[font=\small, text=white] at (85,3) {$\gamma_2$};
  % together: one mode
  \draw[hfConcept, dashed, line width=0.8pt] (5,0.6) .. controls (30,-18) and (60,-18) .. (85,0.6);
  \node[hf/note, text=hfConcept] at (45,-17.5) {together, one fermion mode $c = (\gamma_1 + i\gamma_2)/2$:\\empty or full, and that is the parity bit};
  % local noise
  \draw[hfAlert, line width=0.8pt, decorate, decoration={snake, amplitude=1mm, segment length=3mm}] (-17,3) -- (-4,3);
  \draw[hfAlert, -{Stealth[length=2mm]}, line width=0.8pt] (-4,3) -- (2.2,3);
  \node[hf/note, text=hfAlert, anchor=north] at (-11,-1) {local noise\\touches one half:\\the bit is safe};
  % poisoning
  \draw[hfAlert, -{Stealth[length=2mm]}, line width=0.8pt] (101,3) -- (87.8,3);
  \node[draw=hfAlert, circle, inner sep=0.6pt, font=\scriptsize, text=hfAlert] at (104,3) {$e^-$};
  \node[hf/note, text=hfAlert, anchor=north] at (101,-1) {an electron\\from outside\\flips the bit};
\end{tikzpicture}
\caption{A bit stored in the two ends of a wire. Two Majorana zero modes, $\gamma_1$ and
$\gamma_2$, sit at the ends of a semiconductor nanowire covered by a superconductor, in a
magnetic field along the wire. Neither is a state by itself; together they form one fermion
mode, empty or full, and that is the parity bit. Noise acting on one end reaches only half of
it and cannot read or change the bit. An electron arriving from outside does change it, and
the average wait for that event is the parity lifetime.}\label{fig:majorana-wire}
\end{figure}

\subsection*{The recipe}

In October 2000 Kitaev wrote down the simplest model with a Majorana zero mode at each end: a
chain of sites in which electrons of a single spin direction pair up. Ordinary superconductors
pair electrons of opposite spins, so no known material did that. In 2010 two groups found a way
to fake it with three ordinary ingredients: Roman Lutchyn, Jay Sau and Sankar Das Sarma, and
separately Yuval Oreg, Gil Refael and Felix von Oppen.

\begin{itemize}
\item A semiconductor nanowire with strong \emph{spin--orbit coupling}, such as indium arsenide,
in which an electron's spin is tied to its direction of motion.
\item A magnetic field along the wire, which pushes one spin direction up in energy and out of
the way.
\item A coat of ordinary superconductor, such as aluminum, which by contact induces pairing in
the wire. This is the \emph{proximity effect}.
\end{itemize}

Past a critical magnetic field, the wire enters a \emph{topological phase}, and a Majorana zero
mode should appear at each end. The phase has its own gap, the \emph{topological gap}, which
protects the modes, and it is much smaller than the gap of the superconductor that induces it:
tens of microelectronvolts in Microsoft's devices, against $180\,\mu$eV for aluminum. Two
things go wrong easily. Disorder in the wire can close the topological gap in places, and then
ordinary states trapped near the ends, called \emph{Andreev bound states}, can sit near zero
energy and imitate Majorana modes in most measurements. And the two modes of a finite wire are
not perfectly independent: they overlap a little, by an amount that falls exponentially with
the length of the wire, and that overlap slowly scrambles the phase of the qubit.

\subsection*{The tetron}

One wire gives one parity, which is a bit, not a qubit. In 2017 Torsten Karzig and colleagues at
Microsoft published designs for qubits made of four or more modes, and the simplest is the one
Microsoft builds: the \emph{tetron}. Two parallel topological wires are joined by an ordinary
superconducting backbone into a single small island, with four modes, $\gamma_1$ to $\gamma_4$,
at the four ends.

The island is small, so adding one electron to it costs a noticeable amount of electrostatic
energy, the \emph{charging energy}. That fixes the total parity of the four modes and is the first
defense against poisoning. With the total fixed, the four modes have two states left, and those
are one qubit. Its $Z$ is the parity of the pair along one wire, $i\gamma_1\gamma_2$. Its $X$ is
the parity of a pair across the two wires, $i\gamma_1\gamma_3$. A line of algebra with
$\gamma^2 = 1$ and $\gamma_j\gamma_k = -\gamma_k\gamma_j$ shows that these two operators
anticommute, exactly like the $X$ and $Z$ of Chapter~\ref{ch:state}.

Reading a parity uses a \emph{quantum dot}, a tiny island in the semiconductor that holds a few
electrons, coupled to the two modes being measured. An electron can tunnel from the dot through
either mode, and the two paths close a loop threaded by a magnetic flux, so they interfere. How
easily charge moves on and off the dot, its \emph{quantum capacitance}, then depends on the
parity of the pair, and a radio-frequency signal reflected off the dot reads it in a single shot.
A qubit needs both loops: one along a wire for $Z$ and one across the wires for $X$.

\begin{figure}[htbp]
\centering
\begin{tikzpicture}[x=1mm, y=1mm]
  % wires
  \draw[hfSource, fill=hfSource!10, line width=0.8pt] (0,20) rectangle (70,25);
  \draw[hfSource, fill=hfSource!10, line width=0.8pt] (0,0) rectangle (70,5);
  \draw[hfSource, fill=black!30, line width=0.8pt] (32,5) rectangle (38,20);
  \node[hf/note, anchor=west] at (39,12.5) {superconducting backbone};
  \node[font=\sffamily\footnotesize, text=black!70] at (50,22.5) {topological wire};
  \node[font=\sffamily\footnotesize, text=black!70] at (50,2.5) {topological wire};
  % modes
  \foreach \x/\y/\n in {4/22.5/1, 66/22.5/2, 4/2.5/3, 66/2.5/4} {
    \fill[hfConcept] (\x,\y) circle (2.3);
    \node[font=\scriptsize, text=white] at (\x,\y) {$\gamma_{\n}$};
  }
  % Z loop along the top wire
  \draw[hfConcept, dashed, line width=0.9pt, rounded corners=4pt] (-3,18) rectangle (73,33);
  \fill[hfConcept!25] (35,33) circle (2.6);
  \draw[hfConcept, line width=0.8pt] (35,33) circle (2.6);
  \node[font=\scriptsize] at (35,33) {QD};
  \node[hf/note, text=hfConcept, anchor=west] at (75,29) {$Z$ loop: a dot reads\\the parity $i\gamma_1\gamma_2$\\along one wire};
  % X loop across the left ends
  \draw[hfReflect, dashed, line width=0.9pt, rounded corners=4pt] (-6,-3) rectangle (11,28.5);
  \fill[hfReflect!25] (-6,12.5) circle (2.6);
  \draw[hfReflect, line width=0.8pt] (-6,12.5) circle (2.6);
  \node[font=\scriptsize] at (-6,12.5) {QD};
  \node[hf/note, text=hfReflect, anchor=east] at (-10,12.5) {$X$ loop: a dot reads\\the parity $i\gamma_1\gamma_3$\\across the wires};
  % note
  \node[hf/note] at (35,-9) {one island, four modes, total parity fixed by the charging energy: one qubit};
\end{tikzpicture}
\caption{A tetron, the qubit Microsoft builds. Two topological wires, joined by an ordinary
superconducting backbone, hold four Majorana zero modes. With the total parity fixed, the
parity along one wire is the qubit's $Z$ and the parity across the wires is its $X$. Each is
read by a quantum dot (QD) coupled to the two modes, through an interference loop threaded by a
magnetic flux.}\label{fig:tetron}
\end{figure}

\subsection*{Gates by measurement}

Majorana zero modes are non-abelian anyons. Dmitri Ivanov showed in 2000 that exchanging two of
them applies a unitary on the shared states, not just a phase, which is what Kitaev's 1997
proposal needed. Physically moving modes along wires is hard, though, and in 2008 Parsa
Bonderson, Michael Freedman and Chetan Nayak showed that the same gates can be produced by a
sequence of parity measurements, with no motion at all. That is Microsoft's plan:
\emph{measurement-only} topological computing, with the dots of Figure~\ref{fig:tetron} as the
only tool.

The limit is the gate set. Exchanging Majorana modes produces only Clifford gates, quarter turns,
and the Gottesman--Knill theorem earlier in this chapter says that Clifford gates alone can be
simulated on a laptop. The $T$ gate has to come from \emph{magic states}, special states prepared
outside the topological protection and purified at a cost that error correction has to pay.
Topological protection buys criterion 3 and a smaller overhead per logical qubit. It does not
buy universality.

\subsection*{The checklist}

\textbf{1. Characterized and scalable.} Not shown. Whether the modes in the devices built so far
are topological or ordinary Andreev states is the open question of Chapter~\ref{ch:majorana}. If
they are topological, the wires are made by lithography, like transmons, and could scale the same
way.

\textbf{2. Initialization.} By measurement: measuring a parity leaves the qubit in a known state.

\textbf{3. Coherence.} Two clocks, one per parity. $Z$ survives until a quasiparticle poisons the
island. $X$ survives until the overlap between modes and the noise around them scramble it. In the
2025 devices the first was milliseconds and the second microseconds.

\textbf{4. Universal gates.} Clifford gates by sequences of measurements; the $T$ gate by magic
states.

\textbf{5. Measurement.} The quantum dot and its capacitance, in a single shot. In 2025 the $Z$
readout was wrong 0.5\% of the time and the $X$ readout 16\%.

\textbf{6 and 7. Moving qubits.} The same problem as transmons: everything lives in a
dilution refrigerator and talks through microwaves.

\section{No winner yet}

\begin{center}
\footnotesize
\renewcommand{\arraystretch}{1.25}
\begin{tabular}{>{\raggedright\arraybackslash}p{22mm}*{6}{>{\raggedright\arraybackslash}p{16.5mm}}}
\hline
 & \textbf{1.} characterized, scalable & \textbf{2.} initial state & \textbf{3.} coherence per gate & \textbf{4.} gates & \textbf{5.} readout & \textbf{6--7.} moving qubits \\
\hline
super\-con\-duct\-ing & fabricated, wiring & needs mK & $\sim10^4$ & neighbors only & $\sim$99\% & hard \\
trapped ions & identical; short chains & optical pumping & $\sim10^4$ & all-to-all, slow & $>$99.9\% & natural \\
neutral atoms & identical; thousands & optical pumping & good & Rydberg, slower & good & possible \\
photons & sources, losses & easy & excellent & proba\-bil\-is\-tic & detectors & natural \\
semicond. spins & non-uniform & needs mK & improving & neighbors & improving & hard \\
topological & not shown & by mea\-sure\-ment & -- & Clifford only & -- & hard \\
\hline
\end{tabular}
\end{center}

Every platform trades one criterion against another. The transmon is easy to make and hard to
wire; the ion is perfect and slow; the atom is plentiful and loosely coupled; the photon travels
well and does not interact. Charts that score each platform from 0 to 5 circulate widely, and they
are opinions dressed as data. The table above says something more modest and more durable: none
of them has yet met all five criteria at the scale that error correction needs, which is the
subject of Part~\ref{part:qec}.

\chapter{A Bit That Lasts Twenty Seconds}\label{ch:majorana}

\begin{etymology}
\textbf{Topology} comes from the Greek \emph{tópos}, ``place'', and \emph{lógos},
``study''. It is the part of mathematics that studies what does not change when a
shape is stretched or bent without tearing: a coffee mug and a doughnut have one
hole each, and no small push can change that number.

\textbf{Parity} comes from the Latin \emph{paritas}, ``equality'', from \emph{par},
``equal''. In arithmetic it says whether a number is even or odd. Here it says whether
a fermion mode is empty or full.
\end{etymology}

\section*{Before this chapter}

This chapter is the story of a claim, and it does not stop to explain the physics
behind it. That is done in two earlier sections, and without them most of what follows
will be hard to follow:

\begin{itemize}
\item Section~\ref{sec:identical}, on identical particles: what bosons, fermions and
anyons are, and why exchanging two particles can apply a gate.
\item Section~\ref{sec:topological}, on topological qubits: the superconducting gap and
quasiparticles, the Majorana zero mode and its parity, the recipe for the wire, the
tetron, and how a parity is read through a quantum dot.
\end{itemize}

It also uses the $X$ and $Z$ of Chapter~\ref{ch:state} and the Gottesman--Knill theorem
of Chapter~\ref{ch:hardware}.

\bigskip

On 16 November 2026 Microsoft posted an eight-minute video called \emph{Majorana 2
Explained: The road to 1000x improvement}. In it, the people behind the company's
quantum chip say that its qubits became a thousand times more stable in one year,
that one operation now fits tens of millions of times inside the life of a qubit,
and that a useful machine is coming in 2029. Chapter~\ref{ch:hardware} explained how
the device is supposed to work, and ended its checklist with ``not shown''. This
chapter follows the bet from its start, in 1937, to the video, and then reads the
video against the paper behind it.

\section{A particle that is its own antiparticle}

\begin{history}{1937 \textperiodcentered{} Rome and Naples}
Ettore Majorana, a Sicilian physicist from Enrico Fermi's group in Rome, found a
version of Dirac's equation whose solutions are their own antiparticles: a particle
and its mirror image are the same object. He had the result for years and did not
publish it. He brought it out in \emph{Il Nuovo Cimento} in 1937, when colleagues
convinced him to compete for a professorship, and was appointed at Naples for his
reputation alone. On 25 March 1938 he took a boat from Palermo to Naples and was
never seen again. His body was never found.
\end{history}

Whether any elementary particle is a Majorana fermion is still open; the neutrino is
the candidate, and experiments are still looking. The story of this chapter is about
something else with the same name: the half-electron, half-hole excitations at the ends
of a superconducting wire, which are their own antiparticles in the sense of
Section~\ref{sec:topological}.

\section{Kitaev's idea}

\begin{history}{1997 \textperiodcentered{} Landau Institute, Moscow}
Alexei Kitaev, the same Kitaev of phase estimation (Chapter~\ref{ch:fourier}), and of the
surface codes of Chapter~\ref{ch:qec}, posted \emph{Fault-tolerant quantum computation
by anyons} in July 1997. His proposal: store qubits in non-abelian anyons, whose shared
state depends only on how they were wound around each other, the topology of their
paths. Moving one around another applies a gate, and the gate is exact as long as the
winding is right, however shaky the hand that did it.
\end{history}

In October 2000, by then also affiliated with Microsoft Research, Kitaev wrote down the
wire model of Section~\ref{sec:topological}, with one Majorana zero mode at each end. The
promise attracted Microsoft early. By 2006 the company had a group at the University of
California, Santa Barbara, led by the Fields medalist Michael Freedman, devoted to
topological quantum computing. When the people in the 2026 video say ``we've worked on
topological quantum computing for 20 years'', this is the start they mean.

\section{The peak}

In 2010 two groups published the recipe of Section~\ref{sec:topological}: a semiconductor
nanowire, a superconducting coat and a magnetic field. Experimenters started right away.

\begin{history}{2012 \textperiodcentered{} Delft}
Vincent Mourik, Leo Kouwenhoven and colleagues followed the recipe with indium antimonide
wires and a superconducting contact. They sent current into the end of the wire and saw a
peak of conductance at exactly zero voltage, staying put as they changed the field and the
gate voltages. A Majorana mode at the end of a wire should produce exactly that peak. The
paper appeared in \emph{Science} and started a race.
\end{history}

The trouble is that the Andreev bound states of Section~\ref{sec:topological} produce the
same peak with no topology at all. The following decade was spent trying to tell the two
apart, and most of what follows in this chapter is that same argument, at higher and higher
stakes.

\section{The retraction}

\begin{history}{2018--2021 \textperiodcentered{} Delft}
In 2018 Kouwenhoven's group, partly funded by Microsoft, published \emph{Quantized Majorana
conductance} in \emph{Nature}: the zero-voltage peak had reached the exact height theory
predicts for a Majorana mode, $2e^2/h$. Sergey Frolov and Vincent Mourik, who had worked on
the 2012 paper, asked for the raw data and found that data contradicting the claim had been
left out. On 8 March 2021 the authors retracted the paper. After a new calibration the
plateaus sat 8\% higher than the quantized value, and the authors wrote that they could no
longer claim the observation. An investigation commissioned by Delft found no fabrication,
but found some selection in what had been published.
\end{history}

The retraction did not end the program. It changed its strategy. Microsoft moved most of
the work in house and replaced the single zero-voltage peak with a battery of measurements.

\section{Majorana 1}

In 2022 Microsoft posted the \emph{topological gap protocol}, published in \emph{Physical
Review B} in 2023: a sequence of transport measurements on both ends of a wire that, the
company says, certifies that a region of magnetic field and gate voltage is in the
topological phase. Only devices that pass the protocol move on to the qubit experiments.

On 19 February 2025 Microsoft announced \emph{Majorana 1}, a chip with eight qubits and,
according to the company, a path to a million. The same day \emph{Nature} published a
paper with the measurement behind it: a single-shot measurement of the parity of a wire,
read through a quantum dot whose capacitance changes with the parity. The paper did not
claim the announcement. The editors added a note to the peer review file saying that the
results ``do not represent evidence for the presence of Majorana zero modes in the reported
devices'', and that the paper was published for its device architecture. The press release
said the topological qubit had arrived. The paper said a parity had been read.

Henry Legg, at the University of St Andrews, posted two comments within a month. In the
first he argued that the topological gap protocol has no consistent definition of
``gap'' or of ``topological'', and that its verdict changes with the ranges of field,
voltage and resolution chosen for the analysis. In the second he showed that the same
protocol can classify the regions where the 2025 parity readout was done as either gapped
or gapless.

A qubit needs both loops of the tetron in Figure~\ref{fig:tetron}, because a bit that can
only be read in $Z$ is a classical bit. In July 2025 Microsoft posted the first measurement
of both on one device. The $Z$ parity switched about every $12.4$~ms. The $X$ parity
switched about every $14.5\,\mu$s, almost a thousand times faster, and its readout was wrong
16\% of the time.

\section{Majorana 2}\label{sec:majorana2}

On 2 June 2026 Microsoft announced \emph{Majorana 2} and posted the paper behind it, with 165
authors: \emph{20 Second Parity Lifetime in an InAs--Pb Tetron Device}. Three things changed
at once. The superconductor went from aluminum to lead, whose gap is about eight times larger.
The semiconductor went from indium arsenide to a stack of indium arsenide and indium arsenide
antimonide. And the way the device is tuned changed to a radio-frequency method precise to
about a microelectronvolt. The topological gap that the company reports rose from about 30 to
about $70\,\mu$eV.

The headline number is the $Z$ parity: it now switches on average every $22 \pm 1$~s, some
dwell times reach minutes, against 1 to 12~ms in the aluminum devices. That is where the
``1000x'' comes from; against the 12~ms device the ratio is 1833, and against the 1~ms device
it is 22\,000. The paper is a preprint. In November 2026 it had not been through peer review.

\begin{reflection}{A gap twice as large gives a lifetime a thousand times longer?}
The video says so directly: the material improved ``by a multiple of a few, two to three'',
and that turned into a factor of a thousand. It sounds like an exaggeration, and it is the
most ordinary part of the story.

The $Z$ parity flips when something supplies enough energy to cross the gap. Thermal
processes of that kind have a rate proportional to $e^{-\Delta/k_B T}$, with $\Delta$ the gap
and $T$ the temperature of whatever is causing the flips. Going from $\Delta = 30$ to
$70\,\mu$eV multiplies the lifetime by $e^{40\,\mu\mathrm{eV}/k_B T}$. For that factor to be
1000, $40\,\mu\mathrm{eV}/k_B T = \ln 1000 \approx 6.9$, which gives $k_B T \approx
5.8\,\mu$eV, or $T \approx 67$~mK. A dilution refrigerator sits at about 15~mK, and the stray
excitations in the device are usually hotter than the refrigerator. Sixty-seven millikelvin is
a reasonable number for them.

So a factor of two in an exponent gives a factor of a thousand outside it. That is the same
arithmetic that makes the Boltzmann factor in Chapter~\ref{ch:hardware} drop from a coin toss
at 300~K to almost nothing at 15~mK. The thousand is real. What it measures is the next
question.
\end{reflection}

Three weeks later, on 24 June 2026, \emph{Nature} published Legg's critique of the 2025 paper
as a \emph{Matters Arising}, with Microsoft's reply in the same issue. Legg had gone through
the code of the topological gap protocol, which Microsoft had made public. A plotting filter
was hard-coded to show only the largest region that passed the protocol, hiding a second one;
and the code flipped the bias voltage by position in an array rather than by its physical
value. With those fixed, he got different results from the same measurements, and he wrote
that the coding errors had led to incorrect statements to the peer reviewers. Microsoft
called it ``a minor off-by-one-pixel bug'', said that the protocol is only a tuning step with
no role in interpreting the qubit measurements, and said that Legg offered no alternative
physical model that reproduces their capacitance signal. On Majorana 2, Legg told \emph{The
Register} that it ``is not proven to even be a single qubit''.

The theorists answered too. In August 2026 Binayyak Roy, Jay Sau and Sumanta Tewari showed
that a short segment of the wire in the ordinary, non-topological state can produce the same
flux-dependent capacitance signal that Microsoft reads as the signature of Majorana modes, and
proposed a test to tell them apart. Jay Sau is a coauthor of the 2010 recipe. The people who
invented the platform are among the ones asking for more evidence.

\section{What the video says}

\begin{question}{Is a twenty-second qubit better than a millisecond transmon?}
It is not a twenty-second qubit. It is a twenty-second bit. The number measures how long the
$Z$ parity stays put, which is how long a classical 0 or 1 survives without an outside
electron jumping in. A qubit also has to keep its $X$ value, which is where the phase lives,
and Majorana 2 does not report $X$. The transmon's millisecond in Chapter~\ref{ch:hardware} is
measured on a working qubit with gates. The two numbers are not about the same thing.
\end{question}

The video does its arithmetic in a single sentence: with an operation every microsecond and a
qubit that lives a millisecond, there is one error every thousand operations; with tens of
seconds, ``you can do tens of millions of operations before an error''. With the measured
numbers and the same yardstick of one microsecond per operation:

\begin{center}
\small
\renewcommand{\arraystretch}{1.2}
\begin{tabular}{llrr}
\hline
device & parity & lifetime & operations per lifetime \\
\hline
aluminum tetron (2025) & $Z$ & 12.4 ms & 12\,400 \\
aluminum tetron (2025) & $X$ & 14.5\,$\mu$s & about 15 \\
lead tetron (2026) & $Z$ & 22 s & $2 \times 10^7$ \\
lead tetron (2026) & $X$ & not reported & -- \\
\hline
\end{tabular}
\end{center}

The tens of millions are real for $Z$. The last time anyone published $X$, the same sentence
would have said fifteen. Legg put it in one question: why do they not report the $X$
measurement? The lead may well have improved $X$ too; the video speaks as if it had, and the
paper does not show it. Saying ``qubits 1000 times more reliable'' about a measurement of one
of the two values a qubit needs is marketing.

The rest of the video is the usual. Quantum computers will feed ``the exact, ground-truth
answer'' into AI models so that Copilot can ``create materials and molecules''; they will find
superconductors that work at room temperature. Part~\ref{part:problems} is about the problems
whose cost is actually known, and these are not among them. The one honest line is about
skepticism: ``there's many things that are proprietary, that we can't share all the details''.
Legg found the bugs because the code of the protocol was public. A claim this size, from a
program with a retraction behind it, needs more data, not less.

\section{If it works}

Suppose the next paper shows a long $X$ lifetime, measurements that are right almost every time,
and a test that rules out the ordinary explanations. There would be a topological qubit, and it
would still not be a quantum computer. As Section~\ref{sec:topological} explained, measuring and
exchanging Majorana modes gives Clifford gates and nothing else, and the $T$ gate has to come from
magic states, prepared outside the protection and purified at the cost that Part~\ref{part:qec}
counts. Topological protection buys criterion 3 and a smaller overhead per logical qubit. It does
not buy universality.

Then there is the date. Microsoft's blog says the company ``cut our timeline in half'' and aims for
a scalable quantum computer by 2029. What the United States government has done is narrower. In
DARPA's Quantum Benchmarking Initiative, Microsoft is in Stage C, the final one, which it entered
through an earlier pilot program. Stage C is an independent check of whether the proposed machine
can be built as designed. On 7 October 2026 DARPA added Atom Computing, Diraq, IBM and IonQ to it,
alongside Microsoft and PsiQuantum, and the program's own target for a utility-scale machine is
2033. Being in Stage C means the design is being audited. It does not mean the qubit works.

The row for topological qubits in the table of Chapter~\ref{ch:hardware} still says ``not shown''.
After Majorana 2 it is still the honest entry. What changed in 2026 is the material: a lead
superconductor and a gap more than twice as large gave a bit that holds for twenty seconds, and
that is a real result. The qubit is still waiting for its $X$.


\part{Error Correction}\label{part:qec}
\chapter{Correcting What You Cannot Copy}\label{ch:qec}

\begin{etymology}
\textbf{Error} comes from the Latin \emph{errare}, ``to wander'', ``to stray''.
An error is a state that has wandered away from where it should be.

\textbf{Redundancy} comes from \emph{redundare}, ``to overflow'': \emph{re-},
``again'', and \emph{undare}, ``to surge like a wave''. More than is needed, kept
on purpose.
\end{etymology}

By the middle of 1994 quantum computing had an algorithm worth building a machine
for. Within weeks it also had an objection that seemed to make the machine
impossible. This chapter is about the objection and the answer, because the answer
is what turned the question ``can a quantum computer work?'' into the question this
book asks: ``how many qubits will it take?''.

\section{The objection}

\begin{history}{May 1994 \textperiodcentered{} Santa Fe Institute}
At a workshop at the Santa Fe Institute, a few weeks after Shor's result began to
circulate, Rolf Landauer of IBM argued that quantum computers could not correct
errors. Landauer was the physicist who had shown in 1961 that erasing information
costs energy, and he had spent decades studying what physics allows computers to
do. William Unruh published a similar warning about decoherence the next year.
\end{history}

Classical computers fight errors in three ways, and Shor himself later went through
them to show why the objection was so strong.

\begin{itemize}
\item \textbf{Checkpoints.} Save the state of the computation now and then; if an
error happens, go back to the last save. Saving the state is copying it.
\item \textbf{Error-correcting codes.} Add redundant parity bits that let you detect
and fix a flipped bit. The redundancy is, in effect, partial copies.
\item \textbf{Massive redundancy.} Run many copies, compare them, and take the
majority. If one of five copies fails, you rebuild the fifth from the other four.
That is copying too.
\end{itemize}

The no-cloning theorem (Chapter~\ref{ch:quantum}) forbids copying an unknown
quantum state. There is a second problem: classical errors are discrete (a bit
flips or it does not), while a qubit can drift by any small angle. And a third:
checking a qubit means measuring it, and measuring destroys the superposition.

Without error correction, every gate must be perfect. Factoring a number of
cryptographic size takes on the order of a billion steps, so each one would need an
error below one in a billion. Shor recalled that the best estimate at the time, by
Jeff Kimble, was that experimenters could do one or two quantum gates with about
80\% accuracy.

\section{Encoding without copying}

The simplest classical code repeats each bit three times: $0 \to 000$,
$1 \to 111$. If one bit flips, the majority still gives the right answer. Asher
Peres had already noted in 1985 that the same thing can be written for qubits:
\[
  \alpha\ket0 + \beta\ket1 \;\longrightarrow\; \alpha\ket{000} + \beta\ket{111}.
\]
This is not cloning. Cloning would produce
$(\alpha\ket0 + \beta\ket1)^{\otimes3}$, three independent copies of the
superposition. The encoded state is one entangled superposition spread over three
qubits.

The trick that gets around measurement is to ask only about parities. ``Are qubits 1
and 2 equal? Are qubits 2 and 3 equal?'' Those questions can be answered with extra
qubits without learning anything about $\alpha$ or $\beta$. The answers, called the
\emph{syndrome}, say which qubit flipped, if any, and a gate flips it back. The
superposition is never measured.

And the problem of continuous errors dissolves on its own. A small drift is a
superposition of ``no error'' and ``a flip''. Measuring the syndrome forces it to
choose: either nothing happened, or a full flip happened, which can then be fixed.
Measurement, the enemy of superposition, turns analog errors into digital ones.

\section{Shor's nine qubits}

The three-qubit code fixes bit flips, but it does nothing for phase errors, the
error that turns $\ket0 + \ket1$ into $\ket0 - \ket1$. Worse, a phase error on any of
the three qubits becomes a phase error on the encoded qubit, so the code makes phase
errors three times more likely.

\begin{history}{1995 \textperiodcentered{} Bell Labs}
Shor noticed that a phase error is a bit flip seen in another basis: apply a
Hadamard and $Z$ becomes $X$ (Chapter~\ref{ch:state}). So he protected against phase
errors with a three-qubit code in the Hadamard basis, and protected each of those
three qubits against bit flips with another three-qubit code. Three times three: one
logical qubit in nine physical ones, protected against any error on a single qubit.
The paper appeared in 1995. In 1996, Robert Calderbank and Shor, and independently
Andrew Steane at Oxford, showed how to turn pairs of classical codes into quantum
ones. Steane's code uses seven qubits. The family is called CSS, after the three of
them.
\end{history}

\begin{concept}{Logical qubit}
A qubit stored in a collective state of many physical qubits, so that errors on a
few of them can be detected and corrected without measuring the information itself.

\textbf{What it fixes:} the computation no longer needs perfect hardware, only
hardware that is good enough.

\textbf{What it costs:} many physical qubits per logical qubit, constant syndrome
measurements, and fast classical processing to decode them, for the whole duration
of the computation.
\end{concept}

\section{From possible to engineering}

A code that fixes one error is not enough. The extra qubits and gates of the code
can fail too, and if the code introduces more errors than it removes, it makes
things worse.

\begin{history}{1996--1997 \textperiodcentered{} the threshold theorem}
Several groups proved, independently, the result that settled the question: Dorit
Aharonov and Michael Ben-Or; Emanuel Knill, Raymond Laflamme and Wojciech Zurek;
Alexei Kitaev; and building on Shor's own 1996 work on fault-tolerant gates. If the
error rate of every physical operation is below a certain \emph{threshold}, codes
can be nested or enlarged so that the logical error rate falls as low as you want,
with an overhead that grows only polylogarithmically with the length of the
computation.
\end{history}

The threshold theorem is the reason quantum computing is an engineering problem and
not a dead end. It does not say the engineering is easy. It says how to measure the
distance: get the physical error rate below the threshold, then pay in qubits for
every factor of ten in reliability.

In 1997 Kitaev proposed codes laid out on a lattice, first on a torus and then, with
Sergey Bravyi in 1998, on a flat surface. Every check involves only neighboring
qubits, which suits chips where qubits can only talk to their neighbors.
These \emph{surface codes} have a threshold of about one percent, high by the
standards of the field, and became the default plan for superconducting hardware.
Their price is overhead: hundreds to thousands of physical qubits per logical qubit,
depending on how good the physical qubits are. That overhead is why the estimate for
breaking RSA-2048 in Chapter~\ref{ch:factoring} is counted in millions of physical qubits rather than
the few thousand logical qubits the algorithm needs.


\part{The Market}\label{part:market}

\end{document}
